\pdfoutput=1
\documentclass[11pt]{amsart}
\usepackage{amsmath,amssymb,amsthm,mathtools}
\usepackage[margin=1.2in]{geometry}
\usepackage{microtype,xcolor}
\usepackage[colorlinks=true,linkcolor=blue,citecolor=blue,urlcolor=blue]{hyperref}
\hypersetup{pdfauthor={Bernd Johannes Wuebben},pdftitle={Branched triple covers and wall crossing for local curve counts}}
\setlength{\emergencystretch}{2em}
\newtheorem{theorem}{Theorem}[section]
\newtheorem{proposition}[theorem]{Proposition}
\newtheorem{lemma}[theorem]{Lemma}
\newtheorem{corollary}[theorem]{Corollary}
\theoremstyle{definition}
\newtheorem{definition}[theorem]{Definition}
\newtheorem{example}[theorem]{Example}
\theoremstyle{remark}
\newtheorem{remark}[theorem]{Remark}
\newcommand{\PP}{\mathbb P}
\newcommand{\CC}{\mathbb C}
\newcommand{\ZZ}{\mathbb Z}
\newcommand{\QQ}{\mathbb Q}
\newcommand{\RR}{\mathbb R}
\newcommand{\cO}{\mathcal O}
\DeclareMathOperator{\sgn}{sgn}
\raggedbottom
\setcounter{tocdepth}{1}
\title[Branched triple covers and wall crossing]{Branched triple covers and wall crossing for local curve counts}
\author{Bernd Johannes Wuebben}
\subjclass[2020]{53D45 (primary); 14H30, 14N35, 32Q65 (secondary)}
\keywords{Local Gromov--Witten invariant, branched covering, sextic Galois closure, symplectic Calabi--Yau threefold}
\date{November 24, 2026\\This version: November 24, 2026}
\begin{document}
\begin{abstract}
A branched triple cover of a genus-two curve can change the local
curve count while preserving all determinant signs on unramified covers
of degrees at most seven. At the crossing the degree-three elementary
multiplicity in genus five changes by two in absolute value. At a generic
transverse crossing on a prescribed simply branched triple cover with
symmetric monodromy over a curve of genus at least two, we determine
the local Gromov--Witten change in degree seven and every source genus. Each nonempty subset of
branch values gives a nodal map joining the sextic Galois closure to the
original curve. Its contribution is the degree-one series of a smooth
curve, with sign determined by the number of nodes and the sum of
their branch weights, up to a common orientation sign. On a fixed
symplectic six-torus the other curves compensate each change in the
same integral homology class and elementary genus.
\end{abstract}
\maketitle
\newpage
\tableofcontents
\newpage
\section{Introduction}\label{sec:intro}
A connected triple cover of a genus-two curve, simply branched at two
values, has genus five. Suppose the curve is embedded in a symplectic
six-manifold and we vary its normal Cauchy--Riemann operator while
keeping the curve fixed. A deformation of the cover away from the
curve becomes possible when its pullback normal operator acquires a
kernel. We show that passage through this wall can change the local
curve count both in the class of the triple cover and in degree seven. Every
connected unramified cover through degree seven
can have the same determinant sign on both sides.

To specify the local count, let \(C\) have genus \(g\ge2\) and let
\(N\to C\) be a complex rank-two bundle of degree \(2g-2\), as for the
normal bundle of a curve in a threefold of vanishing first Chern class.
A real Cauchy--Riemann operator \(D\) on \(N\) defines an almost
complex structure on its total space. Over a holomorphic map \(u\),
a section \(\xi\) of \(u^*N\) gives a holomorphic graph precisely when
\(u^*D\xi=0\).
Call \(D\) rigid through degree \(m\) if every connected smooth finite
cover of degree at most \(m\) has zero pullback kernel, and super-rigid
if this holds in all degrees. On an unramified cover the normal
operator has index zero. At a rigid endpoint its determinant has
a sign, taken with the orientation continued from complex-linear operators.

The contribution of an isolated curve in class \([C]\) has an integral
decomposition into the elementary clusters of Ionel and Parker
\cite{ip,diw,WueProduct}. We write \(n_{d,h}\) for the signed multiplicity
of the elementary factor of genus \(h\) in class \(d[C]\) in the
isolating neighborhood. Its normalization is the local series of a
genus-\(h\) curve with opposite normal weights. Bai and Swaminathan coorient the wall of smooth double covers and
describe the bifurcation of a smooth cubic cover
\cite[Theorem~4.12(ii)]{bai}. We realize the cubic crossing on a
prescribed cover and compute the higher-degree change due to its nodal
joins. Doan and Walpuski study chambered invariants of real
Cauchy--Riemann operators \cite{dw}. Along paths of the realized unramified crossings, the local series on
intermediate covers determine the change on the original curve
\cite[Theorems~4.1 and~7.1]{WueUnramified}. We ask what the unramified determinant signs can miss
at a branched crossing.

At an ordinary crossing the pullback kernel is a real line with
embedded graph. Differentiating the pullback equation on a generator
in the branch values and the path parameter identifies those
parameters with the cokernel. We write \(\Delta\) for the value after
a crossing minus the value before it.

\begin{theorem}\label{thm:invisible}
On any genus-two curve there is an ordinary crossing on a prescribed connected \(S_3\) triple cover with two distinct simple branch values, with super-rigid endpoints arbitrarily close on opposite sides. All determinant signs on the connected unramified covers of degree at most seven agree at these endpoints. The local elementary changes satisfy
\[
       |\Delta n_{3,5}|=2,\qquad |\Delta n_{7,13}|=1 .
\]
\end{theorem}

The cover of degree three has trivial deck group. A displacement produces one embedded graph on each side of the
crossing. Their fixed-parameter normal signs are opposite. Subtracting
the incoming count therefore gives two contributions of the same sign. In degree seven a normal height on the
sextic Galois closure vanishes above each branch value, where its
graph can join the zero section.
The proof of Theorem~\ref{thm:invisible} is completed in
\S\ref{p14-sec:chambers}.

\subsection{The wall in degree seven}
For a rigid operator, matching across nodes also gives zero kernel
on connected stable maps in the same degree. Its oriented cokernels
form a real bundle \(V_D\) on the compact moduli stack
\(\overline{\mathcal M}_k(C,d)\) of unmarked connected degree-\(d\)
stable maps from curves of arithmetic genus \(k\). Its real rank
is twice the virtual complex dimension \(2k-2-2d(g-1)\). Set
\[
 L_{d,k}(D)=\int_{[\overline{\mathcal M}_k(C,d)]^{\rm vir}}e(V_D),
 \qquad L_d(D;v)=\sum_{k\ge0}L_{d,k}(D)v^{2k-2},
 \qquad x=(2\sin(v/2))^2.
\]
The normal obstruction comparison identifies this Euler integral with
the local threefold invariant \cite[Theorem~1.2]{zinger}. An elementary
factor of genus \(h\) contributes \(x^{h-1}\) in its first degree
\cite[Theorem~2]{pandhodge}. This distinction between the domain genus
\(k\) and the elementary genus \(h\) will matter whenever a stable map
has contracted components. Write \(H_d=1+d(g-1)\) for the genus of
an unramified degree-\(d\) cover.

Fix a connected simply branched cubic \(f:Y\to C\) with monodromy
\(S_3\). Write \(E\) for its set of branch values, \(b=|E|=2\ell>0\),
and \(h_j=1+7(g-1)+3\ell+j\). The sextic closure has genus
\(1+6(g-1)+3\ell\). Attaching it to \(C\) at \(j\) distinct points
gives precisely the arithmetic genus \(h_j\).

\begin{theorem}\label{p14-thm:wall}
At a generic transverse cubic crossing there is one punctured interval
on which every pullback to a connected stable map of degree seven is
injective, in every source genus. Choose super-rigid endpoints on
opposite sides in this interval and write \(\Delta\) for the value on
the positive side minus the value on the negative side. There are real
weights \(w_i\) indexed by \(E\), with all nonempty subsums nonzero, and
one orientation sign \(\kappa\) such that, with
\(a_I=\kappa(-1)^{|I|}\sgn(\sum_{i\in I} w_i)\),
\begin{align}
 \sum_{k\geq0}\Delta L_{7,k}\,v^{2k-2}
     &=\sum_{\varnothing\ne I\subseteq E}a_I x^{h_{|I|}-1},
            \label{p14-eq:main-series}\\
 \Delta n_{7,h_j}&=\sum_{|I|=j}a_I .\label{p14-eq:main-exponents}
\end{align}
All other degree-seven elementary multiplicities are constant across
the crossing. For fixed \(f\) and \(N\), every chamber of the real hyperplanes
where a nonempty subsum vanishes occurs in this way. The common sign can equivalently be absorbed
by reversing all the weights.
\end{theorem}

The interval assertion concerns degree exactly seven. For the endpoint
factorizations we also need rigidity in other degrees, which is a
residual condition on operators \cite{wendl,dwbn}. On a generic path
through the prescribed crossing one may ensure super-rigidity at a
countable dense set of noncentral times on both sides: evaluation at
each of these times is a submersion on paths with fixed first jet at
zero, by variations supported away from zero. Intersect the corresponding
countably many residual conditions. Rigidity in all degrees on the
whole punctured interval is unnecessary.


\subsection{The sign at the nodes}
The three sheet values of a cubic kernel section have sum zero. On
the sextic Galois closure their conjugates span the real standard
representation of \(S_3\). Fix one transposition \(\tau\). Over each branch value the three
ramification points of the sextic have the three different inertia
transpositions. Exactly one is fixed by \(\tau\). At each of these
chosen points a height in the negative eigenspace of \(\tau\) vanishes:
its value equals its own negative. A nonzero vector on this line has
trivial stabilizer, so its graph requires all six sheets. In
degree seven only one sheet remains to carry the zero height.
Selecting a nonempty set of joins connects the two components.

At a chosen node there are two complex variables: one moves the
attaching point along the graph and the other smooths the node.
Modulo the tangent of the graph, the leading matching term is
antilinear in the smoothing variable. It contributes one minus sign
per node. The other line of the standard representation, fixed by the
transposition, supplies an even height on the entire sextic component.
Every matching equation sees this same real parameter. After the
complex variables are eliminated there is one real equation left.
Pairing the next term at each node with the cubic adjoint gives a real
weight \(w_i\). The coefficient of the common height is, up to common
sign, \(\sum_{i\in I}w_i\). In particular one cannot recover the sign
of a two-node join from the individual node signs. For two branch
values on a genus-two target the four possible pairs
\((\Delta n_{7,12},\Delta n_{7,13})\) are
\[
             (-2,1),\qquad (0,1),\qquad (0,-1),\qquad (2,-1).
\]
A displacement replaces each isolated join by regular graphs.
Degree-one maps to these graphs account for contracted tails and bridges
and give the full factor \(x^{h_j-1}\).

\subsection{Compensation on a compact target}
When the normal model is realized on a closed threefold, deformation
invariance constrains the curves outside its isolating tube. At each
generic endpoint perturbation the bounded-energy cycle space splits
into parts supported in the tube and in its complement. Call the
elementary multiplicities of the latter part complementary.
\begin{theorem}\label{p14-thm:torus}
On the standard symplectic six-torus, fix a genus-\(g\) symplectic
curve of primitive class \(A\) as in \S\ref{p14-sec:compact}. For each
positive even \(b\) the cubics in Theorem~\ref{p14-thm:wall} with \(b\)
branch values can be realized on this curve. The changes in the
complementary elementary multiplicities in class \(7A\) are the
negatives of \eqref{p14-eq:main-exponents}. The sharp maximum of the sum
of their absolute values over all genera is \(2^b-1\).
\end{theorem}


The exterior change is determined in each integral class and genus
even when the curves realizing it differ at the two endpoints. With
all branch weights of the same sign, all subsets of a given size
contribute with the same sign, which gives the sharp bound.

The smooth-cover calculation and the first nodal comparison, where
a double cover joins the original curve, occupy
\S\S\ref{sec:chamber-statement}--\ref{sec:cubic}. The application to two spin endpoints in \S\ref{p10-sec:examples}
uses the absolute series of \cite[Theorem~1.1]{WueSpin}. For the
degree-seven theorem we classify the central kernels in
\S\ref{sec:heptic}, compute the local degree at the joins in
\S\ref{p14-sec:joins}, and count contracted components in
\S\ref{p14-sec:virtual}. Sections~\ref{sec:global}
and~\ref{p14-sec:compact} give the smooth degree-seven covers near a
change of sign and the compact compensation. Appendix~\ref{sec:normalization}
compares two normalizations at a double wall, using the compact
example of \cite{WueCounter}. Appendix~\ref{sec:quartic} treats the
first possible genera in higher degrees. Simultaneous crossings of
several branched covers remain open.

\section{Single walls in degrees two and three}\label{sec:chamber-statement}
Set \(c=g-1\), and let \(\mathcal U\) be the open set of operators invertible on \(C\). The sign \(\epsilon(D)\) is the sign of the determinant at the base. Along a generic path in \(\mathcal U\), the numbers of smooth covers crossing with a kernel are signed using the derivative on the moving branch values and path parameter.
Consider a generic oriented path from $D_-$ to $D_+$ in $\mathcal U$,
with both endpoints rigid through degree three. For $d=2,3$ let
$W_{d,l}$ count pairs $(t,f)$, up to isomorphism of $f$ over $C$,
for which the pullback of $D_t$ to the smooth connected cover $f$
has a kernel and $f$ has $2l$ simple branch values. The term $l=0$
counts unramified covers. The source genus is $1+dc+l$.
Each pair is counted once with the determinant sign of the
Hurwitz and path derivative on its kernel line
\eqref{p10-eq:cubic-intrinsic}. Put $W_d(x)=\sum_l W_{d,l}x^l$ and
$\Delta L_d=L_d(D_+;v)-L_d(D_-;v)$.

\begin{theorem}\label{p10-thm:chambers}
For a generic path as above the polynomials $W_2,W_3$ have finite
support. With $\epsilon=\epsilon(D_-)=\epsilon(D_+)$ one has
\begin{equation}\label{p10-eq:chambers}
\begin{aligned}
 x^{-2c}\Delta L_2 &= W_2(x),\\
 x^{-3c}\Delta L_3 &= 2W_3(x)
        +\epsilon\bigl(W_2(x(1-x)^2)-W_2(x)\bigr).
\end{aligned}
\end{equation}
In particular all the integers $W_{3,l}$ depend only on the endpoints.
\end{theorem}

A smooth cubic crossing of source genus $h$ contributes twice its
sign times $x^{h-1}$. At a double crossing with $2l$ branch values,
a connected cubic map may join the cover to the identity at any
nonempty subset of these values. The change for $j$ joins has sign
$(-1)^j\epsilon\eta_f$, where $\eta_f$ is the double crossing sign.
The factor $(1-x)^{2l}$ in \eqref{p10-eq:chambers} results from the sum over
these subsets. The sign at a join is the real degree of an antilinear
smoothing equation (Theorem~\ref{p10-thm:nodes}).


For a smooth simply branched or unramified covering \(f\) of genus \(h\) with a single transverse kernel and an embedded graph isolated in the first affected degree, the two real amplitudes give \(2\eta_f/|\operatorname{Deck}f|\) after quotienting by the covering symmetries acting on the line. The genus-\(h\) sine monomial already incorporates every higher-genus contracted decoration of the displaced embedded curve. For a double the deck involution exchanges the amplitudes. For the ramified \(S_3\) cubic the deck group is trivial. In the next arithmetic genus the coefficient is \(-(2\eta_f/|\operatorname{Deck}f|)(h-1)/12\). For the double with two branch values on a genus-two target it splits into \(-\eta_f/6\) on the moving-tail locus and \(-\eta_f/24\) at each branch value.
The double-cover displacement and the separate elliptic-tail terms
are proved in Appendix~\ref{sec:smooth-crossings}. We first establish
the corresponding smooth cubic calculation. The joined maps enter
the chamber formula in \S\ref{sec:cubic}.

\subsection{Smooth covers at a crossing}\label{p10-sec:incidence}

Let $f:Y\to C$ be a smooth connected cover, and write $D_f=f^*D$.
We regard covers up to isomorphism over $C$. Let $\mathcal H$ denote
the fixed-target Hurwitz space of the ramification type of $f$.
At a point $(f,t)$ with kernel $K$ and cokernel $P$ the differential
of the family of operators induces a map
\begin{equation}\label{p10-eq:cubic-intrinsic}
 W:T_f\mathcal H\oplus\RR_t\longrightarrow
                  \operatorname{Hom}_{\RR}(K,P).
\end{equation}
It is obtained by differentiating on $K$ and then dividing by the
image of $D_f$, so it is independent of the trivializations. When
$K$ is a line and $W$ is an isomorphism, write $\eta_f$ for its sign.
We orient the source by the complex orientation of the Hurwitz space
and increasing time. On the target we use the full Cauchy--Riemann
determinant orientation, continued from complex-linear operators.
Indeed $\dim_{\RR}P=2\dim_{\CC}\mathcal H+1$ in this case, and
$\det\operatorname{Hom}_{\RR}(K,P)$ differs from
$(\det D_f)^*$ by an even tensor power of $K^*$. This convention
also applies to a double cover.

\begin{proposition}\label{p10-prop:incidence}
Fix endpoints rigid through degree three. For a residual set of smooth
paths between them with fixed symbol and invertible base operator,
every smooth cubic cover with kernel has monodromy $S_3$, trivial deck
group, and only simple ramification over distinct branch values. Its
kernel is a real line, the graph of a generator is embedded, and
\eqref{p10-eq:cubic-intrinsic} is an isomorphism. At that parameter every
other smooth connected cubic cover and every smooth connected double
cover have injective pullback, in every source genus.
\end{proposition}

\begin{proof}
We first vary the operator and a nonzero kernel vector $u$ on one
cubic cover. On an unramified disk write $u_1,u_2,u_3$ for its three
values. Their sum is single-valued on the punctured target and
solves $D$ there. It is bounded at the branch points and hence
extends across them by removal of singularities. Base invertibility
gives $u_1+u_2+u_3=0$. On the connected Galois closure these three
sections span the real sum-zero representation of the monodromy
group. In fact the map from the permutation representation sending
the coordinate vectors to the three conjugate sections kills the
invariant vector, and the remaining real plane is irreducible for
either transitive group $C_3$ or $S_3$. In particular each pair of
sheet values is real-linearly independent on an open dense subset.
A relation $u_i=a u_j$ on an open set with $u_j\ne0$ and $a$ real
gives $(\bar\partial a)u_j=0$, so $a$ is constant. Unique
continuation on the Galois closure contradicts the representation
just computed.

Consider the universal equation $D_fu=0$. A covector $p$ annihilating
its differential in the section and operator directions solves the
homogeneous adjoint equation. On an unramified disk use the base
coordinate to write $p_i$ for its sheet values, with the density
in the real duality pairing expressed in that coordinate.
An arbitrary real zeroth-order variation $A$ of $D$ supported there
gives the variations $(Au_1,Au_2,Au_3)$. Where the first two vectors
are real-linearly independent, these range over all triples of form
values with sum zero. Thus $p_1=p_2=p_3$ on such a disk. Testing the
section variation $f^*w$ gives $\langle p,f^*(Dw)\rangle=0$ for
all smooth $w$. Every smooth base form equals $Dw$ by base
invertibility, so $\sum_i p_i=0$ on every unramified disk. Hence
$p$ vanishes there and identically by unique continuation. The
differential is surjective since its range contains the closed
finite-codimensional image of $D_f$.

Suppose the cover has $b$ simple and $q$ totally ramified values.
These exhaust the cubic profiles. In a chart with ordered branch
values, the $b+q$ distinct values are local complex coordinates.
Indeed the finite covering of their complement is locally fixed by
monodromy and has a unique completion by the prescribed power maps.
Riemann--Roch and Riemann--Hurwitz give
$\operatorname{ind}_{\RR}D_f=-2(b+2q)$. Quotienting the vector by
nonzero real scaling removes one kernel direction. The projectivized
incidence thus projects to the space of operators with index
\[
                    2(b+q)-2(b+2q)-1=-1-2q .
\]
Relative path transversality excludes $q>0$ and leaves a regular
zero-dimensional incidence for $q=0$. A kernel of dimension greater
than one would give a positive-dimensional projective fiber. A cyclic
cubic has no invariant kernel vector, by descent to the invertible
base, and every nontrivial irreducible real representation of $C_3$
has dimension two. Thus the monodromy is $S_3$. A point stabilizer
in $S_3$ is self-normalizing, so the deck group is trivial. At the
remaining covers, projection of the regular linearization to $P$ gives
the surjectivity of $W$. Both sides have dimension $2b+1$.

We give the marked transversality argument because separation of two
cubic covers will require a condition on their complex sheet spans.
At distinct points of a fiber evaluate the difference of their values.
At a ramification point take the complex-linear first derivative of
the section in a frame pulled back from $C$. Its antiholomorphic
derivative already vanishes at a solution since $df=0$ there.
Finally, over a fixed regular base point, evaluate the pair
\[
                     (u_1-u_3,u_2-u_3)\in N_x^2.
\]
Each of these linear maps $J$ on sections kills all base pullbacks.
To check surjectivity for the equation augmented by $J$, hold the
marks fixed and let $(p,\gamma)$ annihilate the section and operator
variations. Off the marks $p$ solves the homogeneous adjoint equation.
The same test with base pullbacks kills its trace, and the localized
variations $A$ force $p$ to vanish on an unmarked disk. By unique
continuation $p$ is supported on the marks. A distribution supported
at these finitely many points has finite order. At each mark we may
prescribe the value and the holomorphic first derivative of a test
section $w$ while arranging that $D_fw$ vanish to any fixed jet order.
In a local expression $D_f=\bar\partial+B$ the Taylor coefficients
with positive antiholomorphic degree are determined successively,
leaving the holomorphic coefficients free. Multiplying the resulting
polynomials by cutoffs gives the required global sections. Now
$\langle p,D_fw\rangle+\gamma(Jw)=0$ forces $\gamma=0$. Elliptic
regularity for the remaining adjoint equation gives $p=0$.

The coincidence of two values has real codimension four. The base
point in an unramified fiber varies in two real dimensions, and a
mark over a branch value, including the pair consisting of the
ramified and the unramified preimage, is discrete over its cover.
Vanishing of the derivative at a ramification point likewise imposes
four real conditions at a discrete mark. A generic path avoids all
three defects, so the graph is an injective immersion of a compact
curve. For the last evaluation the pair has complex rank one in real
codimension two and rank zero in codimension eight. Taking a countable
dense set of fixed base points shows that at every crossing there
is an unramified disk on which the two differences, and hence
$u_1,u_2$, are complex-linearly independent. The rank condition is
unchanged by relabeling the sheets.

It remains to exclude a second cover at the same time. Use separate
Hurwitz and section variables and a common variation $A$ for the
two equations. An annihilating pair of adjoints solves the homogeneous
adjoint equation on each cover and has zero trace on each by the
base-pullback test. For a nonzero cubic adjoint the three values on
an unramified disk span a real plane on an open dense set. Indeed,
view the adjoint there as a dual-bundle-valued differential in the
base coordinate. On the connected regular cover of the complement of
the branch set the trace-free conjugates form the same irreducible
plane as above. Real proportionality on a disk forces a constant
ratio by their common adjoint equation. The argument needs no
extension of the adjoint divided by the base differential across a
branch point.

The tensor pairing with $A$ on that disk is
$\sum_i u_i\otimes p_i$. Substituting the trace identities gives
coefficient matrix
$\left(\begin{smallmatrix}2&1\\1&2\end{smallmatrix}\right)$ in the
first two values of each triple. Thus, when both pairs are
real-linearly independent, the tensor has real rank two and image
$\operatorname{span}_{\RR}(u_1,u_2)$. For a double cover the
corresponding tensor is $2u_1\otimes p_1$, of rank one on an open
dense set for nonzero adjoint. Tensors of these different ranks
cannot cancel. This proves surjectivity for a cubic and a double
cover with marked kernel lines.

For two cubics restrict to the open locus where the first kernel
has a complex-independent pair on some disk. If both adjoints
are nonzero, cancellation for all $A$ forces the two real sheet
planes to agree on a smaller disk on which all the rank conditions
hold. Express any kernel sheet of the second cover as $a u_1+b u_2$
there, with $a,b$ real. Its equation gives
$(\bar\partial a)u_1+(\bar\partial b)u_2=0$. Complex independence
makes $a,b$ constant. On the connected regular cover of the common
punctured target defined by the intersection of the two monodromy
kernels, unique continuation extends all these relations. The two
planes of conjugate sections coincide as representations of the
punctured target's fundamental group. The sum-zero representation
of a transitive cubic monodromy group is faithful, so the monodromy
kernels agree. In their common quotient, all index-three subgroups
are conjugate, both for $S_3$ and for $C_3$. The two covers are
isomorphic on the punctured curve and then on its unique normal
completion. This excludes cancellation for distinct covers. If
only one adjoint were nonzero, its tensor would already have
positive rank. The paired universal equation is therefore
transverse off the diagonal. Its projection index is the sum of the
two single-cover indices, at most $-2$. Here the double-cover index
is $-1$ by the same branch count and Riemann--Hurwitz. A generic
path avoids these pairs.

All uses of transversality may be made on sheet-labelled Hurwitz
charts, with Sobolev order large enough for moving first jets and a
higher differentiability class for coefficient variations. Scaling
is a free kernel direction before taking projective lines. The
equations augmented by evaluations have closed finite-codimensional
range in the section variables. Consequently smooth coefficient
variations suffice for all the surjections above. More explicitly,
exhaust each marked chart by compact subsets away from the relevant
diagonals and normalize kernel vectors in $L^2$. On each compact
subset elliptic estimates give precompactness of the vectors for
smoothly converging paths. Transversality there is open and can be
achieved by finite-dimensional smooth perturbations supported away
from the endpoints, by the parametric transversality theorem. For the
open locus in the paired calculation one also exhausts the locus
where a specified complex minor is nonzero. Countably many such
charts, profiles and genera suffice. This yields the asserted residual
set in the smooth topology. For the analogous double-cover marked
incidence see \cite[Lemmas~5.11 and~5.15]{bai}.
\end{proof}

\begin{proposition}\label{p10-prop:smoothjump}
Suppose a path as in Proposition~\ref{p10-prop:incidence} meets the wall
at a smooth cubic cover $f:Y\to C$ of genus $h$, at time zero. There
is a punctured interval of operators rigid through degree three. For
any sufficiently small $t_-<0<t_+$ the series on the original curve
satisfy
\[
 L_3(D_{t_+};v)-L_3(D_{t_-};v)=2\eta_f x^{h-1},
 \qquad
 L_d(D_{t_+};v)=L_d(D_{t_-};v)\quad(d=1,2).
\]
At a simple double crossing of genus $h$, the corresponding
degree-two change is $\eta_f x^{h-1}$.
\end{proposition}

\begin{proof}
Fix a generator $\xi$ of $K$ at the cubic cover. For a simply
branched cover the quotient of $TY\to f^*TC$ is torsion of length
$b$, the number of simple branch values. Its $H^0$ is the Hurwitz
tangent and its $H^1$ vanishes. Thus the tangential deformation
problem is regular and contributes the complex orientation.

We realize the path in a neighborhood of the zero section by
prescribing the graph equation over a holomorphic projection
$\varphi:Y'\to C$ to be
\begin{equation}\label{p10-eq:cubic-graphequation}
 \varphi^*D_t u+\varphi^*B(u)=0 .
\end{equation}
Here $B$ is a smooth fiberwise real quadratic with values in
$\Omega^{0,1}_C\otimes N$ and $\varphi^*$ pulls back its form argument.
In a connection splitting, the horizontal-to-vertical block of the
almost complex structure prescribes precisely the zeroth-order term
in this equation, with standard complex structure on the fibers
\cite[Lemma~5.9]{bai}.
Cut off this block outside a disk to extend the structures to
$\PP(N\oplus\CC)$. A fiber Fubini--Study form extended using a
unitary connection, plus a sufficiently large multiple of a positive
base form, tames the compact family. The extension can be made for
all fiber dilations used below on a fixed disk of radius greater than
one, by applying the cutoff outside that disk.

Choose a nonzero adjoint $p$ annihilating
$W(T_f\mathcal H)(\xi)\subset P$. On an unramified disk where one
of its values $p_1$ is nonzero and the three $\xi_i$ lie on
distinct real lines, choose smooth real covectors $\alpha_2,\alpha_3$
with $\alpha_i(\xi_i)=0$ and $\alpha_i(\xi_1)=1$. Multiplying
$\alpha_2(u)\alpha_3(u)$ by a compactly supported vector-valued
form pairing positively with $p_1$ gives $B$ with
$\langle p,f^*B(\xi)\rangle\ne0$. Only the first sheet contributes.
This choice is possible for a cubic because its deck group is trivial.

Use local Hurwitz coordinates $y$ and a real kernel amplitude $z$.
Write $u=z(\xi+w)$ with $w$ in a fixed complement of $K$. After
division by $z$ the image equation determines $w=w(y,t,z)$ by the
implicit function theorem. The remaining equation is $F(y,t,z)=0$
in $P$. Its $(y,t)$ derivative at zero is $W(\cdot)(\xi)$, and
$F_z$ pairs nontrivially with $p$ by the choice of $B$. The terms
from differentiation of $w$ lie in the image of $D_f$. Thus the
solution is $(y(z),t(z))$, where $t(0)=0$ and $t'(0)\ne0$. The
undivided reduced equation is $zF=0$. There is one nonzero graph
on each nearby slice, and it is embedded since dilation by $1/z$
makes it converge smoothly to $(f,\xi)$.

For the fixed-time reduced section $zF$, order its columns by the
complex-oriented $y$ variables followed by $z$. Along the nonzero
branch the first $2b$ columns are $zF_y$, and modulo these the last
is $-zt'(z)F_t$. The identifications of the invertible image
complement and the tangential determinant are the same as those
used in \eqref{p10-eq:cubic-intrinsic}. The normal determinant sign of
the embedded graph is consequently
\[
             -\operatorname{sgn}\bigl(z^{2b+1}t'(z)\bigr)\eta_f
                    =-\operatorname{sgn}(t)\eta_f .
\]
Changing $\xi$ changes the orientations induced on $z$ and on $P$
together. The signed number of nonzero curves therefore changes by
$-2\eta_f$. This is the smooth-domain bifurcation of
\cite[Theorem~4.12(ii)]{bai}, realized here for the prescribed path
of first-order operators.

We next isolate these supports in all normalization genera. Suppose
there are additional simple curves of degree at most three contained
in tubes shrinking to the zero section at parameters tending to zero.
Divide their heights by their fiber maxima $\rho>0$. Their equations
on the unit disk have quadratic term $\rho B$ and converge smoothly
to the linear graph equation. Their image cycles have bounded mass
for a common taming form, since its class on the disk bundle comes
from the base and the projection degrees are bounded. Cycle compactness
gives a limiting holomorphic cycle with convergence of currents and
Hausdorff convergence of supports \cite[Theorem~2.11]{diw}. Integration
of a pulled-back area form preserves projection degree, and the limit
contains a point of height one. Every irreducible component has
positive projection degree, since a compact complex curve in an
affine fiber is constant. The normalization of a component reaching
positive height therefore carries a nonzero pullback kernel. By
Proposition~\ref{p10-prop:incidence} this component is the graph of a
nonzero multiple of $\xi$. Its degree exhausts the available three.
The limit is thus embedded of multiplicity one.

Near such a cycle the simple parametrizations are embeddings and
converge up to reparametrization
\cite[Theorem~2.20 and Proposition~2.21]{diw}. In particular their
genus stabilizes, and elliptic regularity gives smooth convergence.
After scaling back, the covers and sections belong to the Hurwitz
chart just analyzed. This proves the existence of one closed tube
and one interval with no additional simple supports. Shrink the
interval so the nonzero branches lie strictly inside. For the linear
endpoint equation use the same argument with unit-height kernel
graphs for a sequence of nonzero times tending to zero, retaining the
mapping multiplicity in each pushforward cycle. The embedded
multiplicity-one limit again gives convergence of the covers and
normalized sections. The reduced equation for this linear problem is
$F(y,t,0)=0$, whose unique nearby solution is $(0,0)$. Thus on
the punctured interval all smooth pullbacks through degree three are
injective. On a nodal domain every positive normalization then has
zero section and the matched constants on contracted components are
zero as well.

Fix an arithmetic genus $k$ and a positive degree $d\le3$. The
stable maps whose images lie in the closed tube form an open and
closed part of the parametrized moduli in the compact completion.
Indeed all their simple supports lie in the interior, and compactness
preserves containment in the closed tube. Their endpoint virtual
degrees agree by deformation invariance
\cite[\S3.1, (3.4)--(3.5)]{diw}. Maps supported on $C$ contribute
the fixed-target Euler number by
\cite[Theorem~1.2, (1.7) and~(1.10)]{zinger}. Using even one nonzero graph
exhausts degree three, so a map with that support has degree one onto
the embedded graph, with arbitrary contracted components.

For clarity, let $D'$ be the invertible normal operator of one such
embedded curve. A degree-one stable map to the curve has one effective
normalization mapping isomorphically, and each connected contracted
subcurve attaches at just one point. For any deformation of $D'$,
restriction to the effective normalization identifies the kernels
of the base and pullback operators and is a left inverse on their
cokernels. Thus the quotient of the two Fredholm complexes has
zero kernel throughout, and the stabilized quotient cokernels form a
bundle, including over nodal maps \cite[Appendix~B]{bai}.
At an invertible endpoint its orientation differs from the full
normal orientation by the determinant sign of $D'$. The complex
normal degree is $2h-2$, since the ambient first Chern class vanishes
on the disk bundle. Deform to the Dolbeault operator of two acyclic
lines of degree $h-1$. The quotient bundle there has complex
orientation, and its degree-one Euler series is $x^{h-1}$ by
\cite[Theorem~2]{pandhodge}. Consequently the change contributed by
the nonzero graphs is $-2\eta_f x^{h-1}$, including every contracted
configuration at its own fixed arithmetic genus. Conservation gives
the asserted change on $C$ and constancy in lower degrees.

For a double cover the involution negates its kernel generator and
identifies the two nonzero amplitudes in an odd cubic perturbation.
The fixed-target comparison then gives the change $\eta_f x^{h-1}$
in degree two (Theorem~\ref{p9-thm:wall}).
\end{proof}

\section{A double cover joined to the zero section}\label{sec:cubic}
For a double wall the first reducible support occurs already in degree three. A branch point supplies a joining flag, and different flags may occur together. This leads to a relation between the signed counts of smooth double and cubic covers.
\subsection{Smoothing a double cover joined to the identity}\label{p10-sec:nodes}

Let $D_t$ be a smooth path of real Cauchy--Riemann operators with fixed
symbol on a complex rank-two bundle $N$ of degree $2g-2$ over $C$.
Suppose $D=D_0$ is invertible, of determinant sign $\epsilon$.
Let $f:Y\to C$ be a connected smooth double cover with $b>0$ simple
ramification points over distinct values. Suppose that $K=\ker f^*D$
is a real line, that its generator $\xi$ is odd under the deck
involution, and that $(f,\xi):Y\to N$ is an embedding. Put
$P_f=\operatorname{coker}f^*D$. Write $s$ for the $b$ complex branch
parameters and the increasing real path parameter together. We assume
the double crossing is transverse, so that the derivative of the
pullback operator on its kernel induces an isomorphism
\[
  W_f:T_{(f,0)}(\mathcal H_2\times\RR)
          \longrightarrow \operatorname{Hom}_{\RR}(K,P_f).
\]
Here $\mathcal H_2$ denotes the fixed-target double Hurwitz chart.
Write $\eta_f$ for the sign of $W_f$ with the absolute Fredholm
orientation, obtained by deformation to a complex operator.

Choose a nonempty subset $I$ of the ramification points. Join each
$r_i\in I$ to $x_i=f(r_i)$ on one copy of $C$, and denote the resulting
finite cubic map by $h_I$. Its arithmetic genus is
$H=g(Y)+g+j-1$, where $j=|I|$. Near $r_i$ take coordinates
$w=z^2$ centered at $x_i,r_i$, and a local frame holomorphic for the
complex-linear part of $D$. Write $D=\bar\partial+B(w)\,d\bar w$ with
$B(w)$ antilinear, and put $v_i=\partial_z\xi(r_i)$. Since $df(r_i)=0$,
$d\xi(r_i)$ is complex linear, and embeddedness gives $v_i\ne0$.
The condition
\begin{equation}\label{p10-eq:node-jet}
                 B(x_i)v_i\ne 0\quad\text{in }N_{x_i}/\CC v_i
\end{equation}
is independent of the coordinates, complex frame and real generator.
Indeed the change of frame adds only a complex-linear term to the
zeroth-order coefficient, and $\xi(r_i)=0$.

\begin{theorem}\label{p10-thm:nodes}
With the double crossing and joining set $I$ as above, assume \eqref{p10-eq:node-jet} at every point
of $I$. The pullback kernel at $h_I$ is isolated in the full fixed-target
chart of genus $H$ times the path. Its local contribution to the change
of the absolutely oriented normal Euler number is
\[
                         (-1)^j\epsilon\eta_f .
\]
If \eqref{p10-eq:node-jet} holds at every ramification point, summing over
all choices of $j$ joins gives $\binom bj(-1)^j\epsilon\eta_f$.
Condition \eqref{p10-eq:node-jet} holds at all such crossings on a generic
smooth path of base-invertible fixed-symbol operators with rigid endpoints.
\end{theorem}

The attaching parameter moves the common value along the tangent of
the double graph. The smoothing parameter contributes in the
quotient by that tangent. Its leading term is conjugate-linear after
a radial change of variable, which accounts for the sign at a node.

The full local calculation is given in Appendix~\ref{app:double-equation}.
Its normalized equation has one complex tangent column from the
moving branch and one complex antilinear column from the smoothing
at each selected point. The projected coefficient in
\eqref{p10-eq:node-jet} makes both columns regular after the radial
change of variable. The base determinant and the absolute crossing
coorientation supply the two remaining signs. Unlike the sextic
calculation below, the double has only an odd kernel line; imposing
its normalization leaves no even height common to the nodes.

To pass from this local degree to the Gromov--Witten change, one must
also allow contracted curves at the joins and at smooth points of
the two branches. The displacement comparison of
Theorem~\ref{p10-thm:descendants} replaces the joined map by regular
embedded graphs of genus \(H\). In source genus \(k\) their degree-one
contribution is \([v^{2k-2}]x^{H-1}\) times the signed degree above.
Proposition~\ref{p10-prop:compact} gives the finiteness needed to sum
these contributions along a generic path.

\subsection{The chamber formula and the balanced product}\label{p10-sec:chamberproof}
\begin{proof}[Proof of Theorem~\ref{p10-thm:chambers}]
By Proposition~\ref{p10-prop:compact} and its double-cover counterpart
there are finitely many crossings on the path. At a smooth double
wall with $2l$ branch values the degree-two change is $\eta_f x^{2c+l}$
(Theorem~\ref{p9-thm:wall}). At a smooth cubic wall with the same
branch count the degree-three change is $2\eta_f x^{3c+l}$ by
Proposition~\ref{p10-prop:smoothjump}. The separation in
Proposition~\ref{p10-prop:incidence} applies to covers of every genus.

Consider the degree-three virtual cylinder at a double crossing.
A nonzero section must restrict nontrivially to a nonconstant
normalization, since matching determines its values on every
connected contracted block. There is no cubic normalization with
kernel at this parameter. There must therefore be a degree-two
normalization carrying the odd kernel and a copy of $C$ at zero
height. Both graphs are embedded. Their common image points are the
$2l$ ramification points, each with one preimage on each graph.
There is no self-node and a contracted block has at most two
nonconstant attaching branches. Thus at least one of these $2l$
points is used to connect the graphs, either directly or by a
contracted bridge. The set $I$ of such points determines the compact
locus of Theorem~\ref{p10-thm:descendants}; every other contracted
component belongs to its allowed decorations. When $l=0$ there
is no connected support of this type.

Put $j=|I|$. The arithmetic genus of the bare join is
$(1+2c+l)+(1+c)+j-1=1+3c+l+j$. Theorem~\ref{p10-thm:descendants}
therefore gives the total change at this double crossing as
\[
  \sum_{j=1}^{2l}\binom{2l}{j}(-1)^j\epsilon\eta_f x^{3c+l+j}
   =\epsilon\eta_f x^{3c+l}\big((1-x)^{2l}-1\big).
\]
For each fixed source genus the supported Euler change is additive
on the finite collection of disjoint compact loci just classified.
Outside them the ordinary oriented cokernel extends over the
parameter. Summing coefficientwise over all arithmetic genera proves
\eqref{p10-eq:chambers}. Substitution $x\mapsto (2\sin(v/2))^2$ is
injective on polynomials, since the substituted series begins with
$v^2$. The first equation determines $W_2$ from the endpoints and
the second then determines $W_3$.
\end{proof}

Let $H$ be Pardon's rational algebra of complex cycle families
\cite{pardon}, and $F_h$ the balanced antidiagonal local curve series
of genus $h$ in this algebra \cite{WueProduct}. In the comparison
algebra $S$ of almost complex cycle families, let $P_d(D)$ denote
the class of the isolated cycle $d[C]$ for the homogeneous normal
equation $D$. The variable $q$ records the degree over $C$. A
\emph{balanced lift through degree three} is a product of
$F_h(q^d)$ with integer exponents, finite at each $d\le3$, whose
image under the comparison $H\to S$ agrees with
$1+\sum_{d=1}^3P_d(D)q^d$ modulo $q^4$.

\begin{theorem}\label{p10-thm:product}
Each endpoint in Theorem~\ref{p10-thm:chambers} has a unique balanced lift
$P^H_{\le3}(D)$ through degree three. Set
$V(x)=2W_3(x)+\epsilon(W_2(x(1-x)^2)-W_2(x))$. Then
\[
 \frac{P^H_{\le3}(D_+)}{P^H_{\le3}(D_-)}
 =\prod_{l\ge0}F_{1+2c+l}(q^2)^{W_{2,l}}\,
  \prod_{l\ge0}F_{1+3c+l}(q^3)^{[x^l]V(x)}
       \pmod{q^4}.
\]
\end{theorem}
\begin{proof}
We apply the finite-degree factorization of \cite[Corollary~3.2]{WueUnramified}
using only the stated endpoint rigidity. Smooth the intersections of a
horizontal torus with $g-1$ parallel vertical tori in a product of
two symplectic two-tori, and include the product in a six-torus. The
resulting genus-$g$ surface represents a primitive class $A$,
as pairing in the first two factors with a vertical torus gives
one. Its symplectic normal bundle has degree $2g-2$. Identify the
normal germ with the given one along the surface and extend its graph
almost complex structure to a global tame structure, by interpolation
to a compatible extension in a sufficiently small tube
\cite[proof of Proposition~3.1]{WueUnramified}.

Through energy $3\omega(A)$ the only cycles contained in the
closed tube are multiples of the zero section. Indeed a component
in that tube projects holomorphically in positive degree; a component
of nonzero height would give a nonzero kernel on its smooth
normalization. The degree agrees with the positive energy label since
the tube retracts to $C$. The cycles in the interior give a compact
open and closed insertion through this cutoff. Compactness of cycles
and convergence of supports preserve separation from the tube
boundary under a small perturbation of the ambient structure
\cite{diw}. Transport to a generic perturbation and finite-energy
cluster decomposition express this insertion as a product of the
images of the balanced factors with integral exponents
\cite[Corollary~3.2]{WueUnramified}. Their coefficients generate the integral
cycle algebra by \cite[Theorem~1.1]{WueGenerators}. The construction at that energy involves
only labels through three. Lift each factor to $H$. If two products
first differed at label $d$, the difference at that label would be
a finite linear combination of $[q]F_h(q)$ whose image vanishes
in $S$. Independence of these images under fixed-genus detection
\cite[Theorem~5.1]{WueProduct} gives
uniqueness.

It remains to identify the exponents of these geometric products.
At each endpoint, fixed-target comparison in the compact realization
\cite[Theorem~1.2]{zinger} identifies the connected local series
in degrees $1,2,3$ with the integrals defining $L_d$. The series
of an elementary factor in its first label is $x^{h-1}$
\cite[Theorem~2]{pandhodge}. In degree one the endpoint factor is
$F_g(q)^{\epsilon}$ at both ends. At labels two and three the only
proper divisor is one. Its contribution to the logarithm of the
ratio cancels, and an exponent difference in genus $h$ at either
new label therefore contributes exactly its multiple of $x^{h-1}$.
The finite Laurent monomials in $x$ are linearly independent after
sine substitution. Equation~\eqref{p10-eq:chambers} determines all
the differences asserted above.
\end{proof}

Write $n_{d,h}(D)$ for the exponents in this lift. The original
fixed-target series in degree $d$ vanishes below genus $1+dc$ by
Riemann--Hurwitz. For $d=2,3$ the only proper-divisor contribution
comes from the genus-$g$ balanced reference in label one, whose
degree-$d$ term satisfies the same bound. The remaining monomials
$x^{h-1}$ have distinct leading orders, so $n_{d,h}(D)=0$ for
$h<1+dc$ \cite[Proposition~3.1 and Corollary~3.2]{WueUnramified}. Set
$N_d(D;x)=\sum_{l\ge0} n_{d,1+dc+l}(D)x^l$.
\begin{corollary}
The polynomial $N_3(D;x)-N_2(D;x+x^3)+N_2(D;x)$ modulo two is
constant between rigid endpoints joined in $\mathcal U$.
\end{corollary}
\begin{proof}
The two exponent changes are $\Delta N_2=W_2$ and $\Delta N_3=V$.
Modulo two one has $x(1-x)^2=x+x^3$, $2W_3=0$ and $\epsilon=1$.
\end{proof}

\subsection{Two spin endpoints}\label{p10-sec:examples}

Let $A,B$ be distinct theta characteristics on $C$ with
$H^0(C,A)=H^0(C,B)=0$. For a theta characteristic $S$, take the
positive spin operator $D_S$ associated with a lift of the
uniformizing representation to $\mathrm{SL}(2,\RR)$. In a spin frame
on the upper half-plane it is given by
\[
                D_S u=\bar\partial u+
                    \frac{i\bar u}{2\operatorname{Im}z}\,d\bar z .
\]
The choice of lift specifies the signs in the transition functions of
the frame. On the same underlying smooth complex rank-two bundle
consider the ordered endpoints
\[
               D_-=D_A\oplus D_A,\qquad D_+=D_A\oplus D_B .
\]
They are rigid in every degree. Indeed, on a finite cover $f$, the
equation for one spin summand makes $d(u^2 f^*dz)$ a nonnegative area
form whose density is positive wherever $u\ne0$ and $df\ne0$.
Integration forces $u=0$ on every positive-degree component, and nodal
matching gives vanishing on each attached contracted subcurve. This
also proves base injectivity when the antilinear term is multiplied by
any positive scalar. At scalar zero each base Dolbeault operator is
invertible by non-effectivity and Riemann--Roch. Hence the two endpoint
determinants have sign $\epsilon=+1$. To connect the operators in the
base-invertible locus, turn off the antilinear term on the second
summand, join the two Dolbeault operators in the complement of the
theta divisor in $\operatorname{Pic}^{g-1}(C)$, and restore the term
on $B$. The complement is path connected since the divisor has real
codimension two. Trivialize the smooth bundles along the interval and
perturb the path relative to its endpoints.

We evaluate the two endpoints using the absolute partition formula
of \cite[Theorem~1.1]{WueSpin}. Write $c=g-1$ and
$X_k=(2\sin(kv/2))^2$. At the repeated spin the disconnected
degree-$d$ coefficient is the sum over partitions $\lambda\vdash d$
with hook weight $\prod_{a\in\lambda}X_{h(a)}^c$. At the distinct
pair with positive determinant the same weights are summed over just
the self-conjugate partitions. These two absolute evaluations are
the input from spin theory. The signed number of crossings between
them follows from Theorem~\ref{p10-thm:chambers}.

Let $Z_d^-$ and $Z_d^+$ be the respective disconnected coefficients,
requiring positive degree on every connected component. For $d\le3$
the partition sums read
\[
\begin{array}{c|cc}
 d & Z_d^- & Z_d^+ \\ \hline
 1 & x^c & x^c \\
 2 & 2(xX_2)^c & 0 \\
 3 & 2(xX_2X_3)^c+(x^2X_3)^c & (x^2X_3)^c .
\end{array}
\]
Indeed the row and column of size three have hooks $1,2,3$, while
$(2,1)$ has hooks $1,1,3$. The logarithm of the disconnected series
has degree-two coefficient $Z_2-Z_1^2/2$ and degree-three coefficient
$Z_3-Z_1Z_2+Z_1^3/3$. Using $X_2=x(4-x)$ and $X_3=x(3-x)^2$ gives
the connected differences in the order $D_+$ minus $D_-$,
\begin{align*}
       x^{-2c}\Delta L_2 &=-2(4-x)^c,\\
       x^{-3c}\Delta L_3 &=2(4-x)^c\bigl(1-(3-x)^{2c}\bigr).
\end{align*}

\begin{corollary}\label{p10-cor:spin-counts}
For every generic path between these ordered endpoints in the
base-invertible locus, the signed polynomial of smooth cubic crossings is
\begin{equation}\label{p10-eq:spin-crossings}
       W_3(x)=\bigl(4-x(1-x)^2\bigr)^c
                         -\bigl((4-x)(3-x)^2\bigr)^c .
\end{equation}
The largest normalization genus with nonzero signed count is $6c$.
Every such path contains at least $8c$ smooth cubic cover--parameter
pairs in that genus.
\end{corollary}

\begin{proof}
Insert $W_2(x)=-2(4-x)^c$ and the cubic difference into
\eqref{p10-eq:chambers}. The terms in $(4-x)^c$ cancel, giving
\eqref{p10-eq:spin-crossings}. The two cubics on the right, before taking
$c$th powers, have common leading coefficient $-1$ and difference
$-8(x-2)^2$. Their difference of $c$th powers thus has degree $3c-1$
and leading coefficient $-8c(-1)^{c-1}$. By Riemann--Hurwitz, a
cubic cover with $2l$ simple branch values has genus $1+3c+l$.
Taking $l=3c-1$ proves both assertions.
\end{proof}

On a genus-two target one can take any two distinct even theta
characteristics. Each has degree one, hence at most one independent
section, and even parity forces non-effectivity. Here
$W_2(x)=-8+2x$ and $W_3(x)=-32+32x-8x^2$. The signed counts of smooth
cubics are $-32$, $+32$, and $-8$ in normalization genera four, five,
and six, respectively. In particular the half-coefficient of $x^2$ in
$x^{-3}\Delta L_3$ is $-10$, and subtracting the double-cover term in
\eqref{p10-eq:chambers} raises the signed four-branch count to $-8$.
Every such path therefore contains at least $72$ smooth cubic
cover--parameter pairs across these three genera. In higher genera
the signed count vanishes.

\section{Normal kernels on the sextic Galois closure}\label{sec:heptic}
In degree seven the normal kernel at a ramified cubic has two real directions on its Galois closure. Projectivizing this plane does not separate the equations at the different nodes. We identify the one remaining real equation after solving the complex matching equations.
\subsection{The cubic and its Galois closure}\label{p14-sec:kernels}

Fix a complex rank-two bundle $N$ of degree $2g-2$ over a closed Riemann
surface $C$ of genus $g\ge2$. A real Cauchy--Riemann operator $D$ on $N$
has the Dolbeault symbol and a real-linear term of order zero. All kernels,
cokernels and indices below are real. Let $f:Y\to C$ be a connected cubic
with monodromy $S_3$ and $b>0$ distinct simple branch values. Write
$Z\to C$ for the smooth Galois closure and $V$ for the real plane
$\{(a_1,a_2,a_3):\sum a_i=0\}$ with its permutation action. Locally
the involution at a ramification point acts by $z\mapsto-z$, with
$w=z^2$ on the base. The fixed-target Hurwitz space of $f$ has complex
dimension $b$.

Let $D=D_0$ be invertible on $C$. At an \emph{ordinary transverse
cubic crossing} of a path $D_t$, the kernel of $D_f=f^*D$ is a line
$\mathbb R\xi$, the graph of $\xi$ in $N$ is embedded, and differentiation
at $\xi$ in the branch positions and increasing time gives an isomorphism
\begin{equation}\label{p14-eq:lowertransverse}
 W_f:T_f\mathcal H_f\oplus\mathbb R\longrightarrow
          \operatorname{coker}D_f .
\end{equation}
Indeed $\operatorname{ind}D_f=-2b$. The trace of $\xi$ is zero by
removability on $C$ and base invertibility. Its three conjugates on $Z$
therefore span one copy of $V$. Over a regular disk they give three
solutions with sum zero. Evaluation of this plane will have real rank
two at regular points and complex span of dimension two on at least
one regular disk. We explain both the generic assertion and a construction
on any prescribed $f$.

\begin{lemma}\label{p14-lem:kernelplane}
Suppose two of the sheet values form a complex frame on a regular disk.
For a residual set of arbitrarily small smooth variations on this disk
preserving the three solutions, every connected smooth finite cover
$u:X\to C$ has
\[
                     \ker u^*D=V^{K_u}.
\]
Here $K_u\le S_3$ is the image in the lower monodromy of a sheet
stabilizer for $X$ in the joint monodromy of $X$ and $Z$. The
identification uses pullback to a common normal cover.
\end{lemma}

\begin{proof}
Use real variations annihilating the two frame vectors $s_1,s_2$.
On any smaller disk a dense set of these variations attains the
infinite-order Petri condition at a specified interior point. This
condition says that a tensor of formal solution and adjoint-solution
jets with zero pointwise product jet is zero. For unrestricted
variations use \cite[Corollary~5.10]{wendl} and
\cite[Theorems~1.6.17 and~2.5.3]{dwbn}. For a nearby unrestricted
operator $D'$ with such a jet, solve $D's'_j=0$ locally with $s'_j$
close to $s_j$: a Cauchy right inverse on a disk of radius $r$ has
norm $O(r)$ on bounded functions, and interior elliptic estimates
give closeness in any prescribed finite smooth norm. Conjugate by the
complex gauge taking $s_j$ to $s'_j$, then cut off the difference from
$D$ away from the point. This preserves the Petri condition there and
annihilates the original columns on the entire disk.

Consider an irreducible real finite monodromy system $W$ inequivalent
to the prescribed standard system, with its ineffective punctures
removed. A nonzero twisted solution provides columns $w_i$ which,
together with $s_1,s_2$, are real linearly independent as germs on a
common normal cover. This follows from inequivalence, irreducibility
and unique continuation. An annihilator of the section and operator
variations has adjoint columns $v_i$ with
\[
        \sum_i w_i\otimes v_i=s_1\otimes\alpha_1+s_2\otimes\alpha_2.
\]
This identity is the pointwise annihilator condition for maps killing
the two prescribed columns. Test it against $[D,\Phi]$, with $\Phi$ a
compactly supported complex-linear endomorphism of $N$. The integral
on the left vanishes. The sections $\Phi s_1,\Phi s_2$ are
independently arbitrary, so both $\alpha_j$ solve the adjoint equation
on the disk. The Petri condition and independence of the input jets
force all $v_i$ to vanish to infinite order. Adjoint unique
continuation gives surjectivity.

If the $a$ effective punctures move, the projective incidence has
projection index
\[
       2a-1-2\sum_x\dim(W/W^{I_x})\le -1 ,
\]
by the orbifold index formula \cite[Proposition~2.8.6]{dwbn};
$I_x$ denotes inertia at $x$. Exhaust each configuration chart by
compact blocks whose punctures avoid a common smaller test disk.
At a good Petri jet over a block, the unit-normalized solution set is
compact by elliptic estimates on the compactified normal covers.
Finitely many allowed variations give surjectivity near that set.
Finite-dimensional Sard gives arbitrarily small perturbations excluding
all such solutions over the block. Absence is open there. There are
countably many blocks and discrete monodromy types. The remaining
type $V$ is faithful, so each of its copies on a common normal cover
descends to $Z$. Each has a one-dimensional invariant line on $Y$;
openness of the given simple kernel on the fixed $Y$ permits just one
copy. Taking stabilizer invariants proves the formula, by bounded
descent and removability over the branch points.
\end{proof}

At an ordinary crossing on a generic path the required regular-fiber
evaluation is injective everywhere. Indeed the two independent
differences of the three heights at a marked regular value can be
varied independently. For completeness, the universal section
equation at a cubic with complex frame disk admits arbitrary trace-free
forcing there by variation of $D$. In the adjoint test the trace
vanishes by base invertibility, hence the annihilator vanishes off
the marked points by unique continuation. Any remaining adjoint
distribution has finite order and is supported at the marks. Prescribe
holomorphic value and derivative jets there and solve recursively
for mixed jets so that the image under $D_f$ vanishes to higher
order than that distribution. Taylor polynomials with cutoffs realize
these jets. Independence of the resulting two value differences
kills their multipliers, and elliptic regularity removes the supported
adjoint. The same argument with a holomorphic
linear jet at the ramifying point gives independent variations of the
height difference at the two points in a branch fiber and the
ramification derivative. One may perform these tests without first
assuming a complex frame disk: two real-independent sheet values
suffice for the forcing test. Such a disk of real independence exists,
since collinear nonzero solutions on an open disk have constant real
ratio by their equation, contradicting the standard monodromy.
Real rank one for a pair of real four-vectors has codimension three,
and rank zero has codimension eight. The regular mark contributes
only two parameters along a transverse crossing. The complex
dependence locus at a fixed regular mark has codimension two.
These tests give the claimed ranks, including complex independence
on some disk. On continued transverse crossing charts the preceding
residual perturbations are available by multiplication by a time
cutoff.

If a nonzero section on a positive component is described by
$a\in V^{K_u}$, the component factors through
$Z/\operatorname{Stab}(a)$. Each nonzero stabilizer orbit has size
three or six. At a branch value the heights have values $u,u,-2u$
with $u\ne0$ by embeddedness. The kernel of evaluation is the
negative eigenline of inertia, generated, in a suitable ordering, by
$(1,-1,0)$. A nonzero vector in this line has trivial stabilizer.
Thus a positive component with a nonzero section vanishing somewhere
factors through $Z$.

In total degree seven some positive component must have zero height.
Connectedness forces a nonzero component to meet height zero since
contracted components carry constant sections. The degree count
leaves exactly one copy of $Z$ and one copy of $C$. Up to the deck
group their nonzero height on $Z$ is a scalar multiple of a fixed
root $s$, negated by a transposition $\tau$. Over each branch value
$x_i$ there is a unique ramification point $r_i$ of $Z$ fixed by
$\tau$. Its height is zero. The other two heights over $x_i$ are
nonzero and opposite. At a regular value the six heights are
distinct by real injectivity. At all three ramification points the
vertical derivatives are nonzero, being proportional to the
difference of two distinct coordinates of the root times the
ramification derivative of the lower embedded graph. The sextic
graph is therefore embedded. Its intersections with the zero section
are precisely the $b$ points just described, with vertical and
horizontal tangent lines respectively. The direct nodes or connected
contracted subcurves joining the two components determine a unique
nonempty subset $I$ of these points. No point is used twice since
the positive normalizations are embedded. Riemann--Hurwitz gives
$g(Z)=1+6(g-1)+3b/2$. Joining at $j$ points before attaching any
contracted components gives genus $1+7(g-1)+3b/2+j$.

\subsubsection{Construction on a fixed covering curve}
Here and in Section~\ref{p14-sec:joins}, prescribing the crossing permits
variation of the operator and its solution, with $C,N,f$ fixed.

\begin{proposition}\label{p14-prop:realize}
Every connected simply branched cubic $f$ as above admits an ordinary
transverse crossing with the stated evaluation ranks and central kernels.
Its branch weights can be chosen in any prescribed chamber of the
nonempty subset-sum hyperplanes.
\end{proposition}

\begin{proof}
Over the regular locus map the flat $V$-bundle real-linearly to
$N_{\mathbb R}$. At a branch send an even and an odd basis vector to
$e_0$ and $zv_0$, where $e_0,v_0$ form a complex basis.
This is equivariant under $z\mapsto-z$. Fix also a complex frame on
a regular disk. Extend relatively over the complement by a smooth
bundle map; transversality avoids the rank-one and rank-zero strata
of codimensions three and eight over the surface. The maps of the
three sheet vectors define an embedded cubic graph. Start with a
normal operator equal to $\bar\partial$ in the branch frames. Off
the branches prescribe a real zeroth-order correction on the image
of the flat plane by sending each parallel column to minus its
image under this operator. Extend smoothly on a complement, using
the zero correction near the branch centers. This makes the
prescribed heights solutions.

We verify regularity on the fixed $f$. On the frame disk variations
killing the two heights can first make $D$ invertible. A nonzero
base kernel vector cannot belong to their real span on an open
disk: its two real coefficients would have zero $\bar\partial$ and
be constant, contrary to the absence of an invariant in $V$.
Testing such a vector against a nonzero adjoint therefore gives
a nonzero kernel-to-cokernel matrix entry, by variation of $D$
and unique continuation. In a finite reduction the rank increases
under a small multiple of this variation. Iterate until the
index-zero base problem is invertible.

Augment $D_f$ by $b$ complex pole columns at the ramifying points.
Each column is the image off the pole of $\chi a/z$, where $a$ is
in the complex derivative line of the height there and $\chi$ is an
even cutoff. The pulled-back antilinear coefficient contains a
factor $\bar z$, so the forcing lies in $L^p$, $p>2$. The derivative
of the inverse sheet $(w-\alpha)^{1/2}$ is $-\alpha/(2z)$, giving
exactly the branch-motion column modulo the regular domain. The
augmented index is zero and the known kernel line persists. Suppose
an additional augmented solution exists. Its trace extends by
cancellation of the two odd poles at each ramifying pair and is zero
by invertibility. The adjoint of the augmented problem also has
trace zero by testing base pullbacks. Two of its three components
are real-independent on an open dense subset: proportionality on
an open set would give a constant real relation by the adjoint
equation and hence on the normal cover, whereas the nonzero
trace-free conjugates span the irreducible standard representation.
Writing the three conjugates of the extra solution as $h_i$, annihilation
of the constrained variations would put both $h_1-h_3,h_2-h_3$ in
the real span of $s_1,s_2$ there. Their coefficients would again be
constant. By trace zero and continuation all $h_i$ would lie in
the original standard copy and have no poles, so its sheet-invariant
line is the only possibility. Thus a rank-increasing variation exists
until the augmented kernel is just the prescribed line. One unrestricted
time variation on the disk then pairs nontrivially with the remaining
adjoint: the available forcing on the three sheet heights is any
triple of sum zero. This proves \eqref{p14-eq:lowertransverse}.
Lemma~\ref{p14-lem:kernelplane} applies in its open regular neighborhood.

Section~\ref{p14-sec:joins} below proves the assertion concerning strict
branch weights by changing the coefficients at the branch points and
restoring the equation on this same $Y$. This step also supplies the
additional regularity needed for the signed local degrees.
\end{proof}

\section{One real obstruction at several joining nodes}\label{p14-sec:joins}

Fix the root $s$ on $Z$ and its transposition $\tau$, and choose once
and for all an even generator $s^+$ of the other eigenspace in $V$.
Over $x_i$ write $v_i=\partial_zs(r_i)$ and $e_i=s^+(r_i)$.
Use the quotient cubic $Y=Z/\langle\tau\rangle$. The image $y_i$
of $r_i$ is an unramified point of $Y\to C$. Trivialize $N$ in a
holomorphic frame for the complex-linear part of $D$, writing
$D=\bar\partial+B(w)\,d\bar w$, with $B$ antilinear. In that
frame put $B_i=B(x_i)$ and $L_i=N_{x_i}/\mathbb Cv_i$. We shall have
$[B_i v_i]\ne0$ and $[e_i]\ne0$ in $L_i$. Let $\psi$ generate
the adjoint line on $Y$ annihilating its branch columns in
\eqref{p14-eq:lowertransverse}. In coordinate $w$ it pairs
$a\,d\bar w$ with $\operatorname{Re}\psi(w)(a)\,dA(w)$.
Define
\begin{equation}\label{p14-eq:branchweights}
 2\overline{\beta_i}[B_i v_i]=[e_i],\qquad
 H_i(q)=\tfrac{\pi}{2}\operatorname{Re}\psi(y_i)(qv_i),
 \qquad w_i=H_i(\beta_i).
\end{equation}
Here the $w_i$ on the right are real weights. The local domain coordinate
$z$ and the target coordinate $w$ are those used in the smoothing chart
below. The same $s,s^+,\psi$ are used at all the joining flags.

\begin{proposition}\label{p14-prop:signs}
Join the sextic to the identity over a nonempty set $I$ of branch
values. Under the regularity conditions below, if
$\sum_{i\in I} w_i\ne0$, the joined map without contracted components (the bare map) is an isolated zero in the full
finite-cover chart times the parameter interval. Its local Euler change is
\begin{equation}\label{p14-eq:localsubsetsign}
              a_I=\kappa(-1)^{|I|}
                     \operatorname{sgn}\Bigl(\sum_{i\in I} w_i\Bigr),
\end{equation}
where $\kappa\in\{1,-1\}$ is common to all subsets at the crossing.
All nonempty subsums are nonzero at a generic crossing. Every
strict sign pattern of these sums can be realized on the fixed cubic.
\end{proposition}

\subsection{The derivative along the nodal stratum}
There are $3b$ independently movable simple ramification points on
$Z$, even when their images coincide. At each join there are two more
complex parameters $p,q$. In coordinates on the two branches the
family is $z y=q$, $w=(z+p)^2+y$ after translating the branch center.
Elimination gives the algebra with basis $1,z,z^2$ and relation
\begin{equation}\label{p14-eq:cubicpolynomial}
                   z(z+p)^2-wz+q=0 .
\end{equation}
At the node the relative deformation sheaf has length three: at the
central fiber its local algebra is the cokernel of the Jacobian of
$(zy,z^2+y)$, or equivalently $\mathbb C[z]/(z^3)$ after eliminating
$y$. The other simple ramification points have length one. The
relative deformation complex is supported at these finitely many
source points and the finite-map chart is smooth of complex dimension
$3b+2j$, $j=|I|$. Its central stabilizer is $\langle\tau\rangle$.
Write $P_Z=\operatorname{coker}D_Z$. Normalization and evaluation
give the joined kernel $\mathbb Rs$ and the joined cokernel $P_I$ with
\begin{equation}\label{p14-eq:joinedsequence}
 0\longrightarrow
 (\bigoplus_{i\in I}N_{x_i})/\mathbb R(e_i)_{i\in I}
 \xrightarrow{\iota}P_I\longrightarrow P_Z\longrightarrow0 .
\end{equation}
Use the convention that the ordinary forcing of a cutoff height
with value $a$ on the sextic and zero on the identity represents
$\iota(a)$.

To identify the image in $P_Z$, split under $S_3$. The trivial
part is zero. The sign part is the odd cokernel on the discriminant
double $T=Z/A_3$ and has dimension $2b$. The standard part has
multiplicity $2b+1$, its multiplicity space being
$\operatorname{coker}D_f$. These identifications follow by stabilizing
with forms away from the branches on each quotient. A bounded
invariant solution descends and extends across the punctures, and
the Fredholm indices agree by deformation to a complex-linear
operator. The odd operator on $T$ has no smooth kernel. Augment it
by a complex plane of odd poles $\chi a v_i/z$ at each ramification
point above $x_i$, as in the proof of Proposition~\ref{p14-prop:realize}.
Its resulting index is zero. We require it to be invertible. This
is open dense with the lower graph fixed. A nonzero odd polar
solution cannot be a real combination of the two lower heights on
an entire open disk: the coefficients there would be constant and
the sign representation on the punctured normal cover has zero
intersection with the standard one. An allowed zeroth-order
perturbation can prescribe arbitrary forcing on this solution on
a smaller disk. The two odd adjoint contributions add. Unique
continuation and rank increase in a finite reduction give the
assertion.

Write $B_0$ for the real space of sextic branch velocities and
increasing time. Their derivative on $s$ embeds $B_0$ as a hyperplane
in $P_Z$. To see this, first take the common shift of the three
ramification points over each $x_i$. These shifts and time preserve
$S_3$ symmetry. On the even vector their standard component is
\eqref{p14-eq:lowertransverse}; on the odd vector the same columns span
the odd standard summand since $\operatorname{End}_{S_3}V=\mathbb R$.
The additional separate shifts of the fixed flags give on the sign
part the discriminant pole columns computed above, with nonzero
factor $1/3$ from averaging over $A_3$. Indeed a branch shift $\alpha$
produces the principal part $-\alpha v_i/(2z)$, by differentiation
at fixed regular target point. Subtract the principal part and use
removability for the bounded remainder; changing the cutoff changes
the forcing by the image of a regular section. Finally the opposite
shifts at the two $\tau$-exchanged flags, applied to $(1,-1,0)$,
give one third of the common shift applied to $(1,1,-2)$.
At inertia $(13)$ and $(23)$ this follows by taking the ratios
$1/3$ and $-1/3$ of the coordinate differences. The even vector has
zero derivative at the remaining flag. These columns give exactly
the cubic branch columns on the invariant quotient. The remaining
functional $\Lambda$ is thus the half-trace to the cubic paired
with $\psi$. Extend it to $P_I$ by \eqref{p14-eq:joinedsequence}.
For the full $q_i=0$ derivative, each attaching column contributes
$v_ip_i$ in matching. They are independent modulo the one diagonal
real line since $[e_i]\ne0$. Denote this injective nodal derivative,
including time, by $J$, and its variables by $h$.


\subsection{The smoothing terms}
Let $F$ be the normalized obstruction equation, $q_i$ the complex
smoothing parameters, and $Q_i=q_i\log(1/|q_i|)$, continued by zero.
In the matching fibers the leading displacement is
$2\iota((\overline{Q}_i B_i v_i)_{i\in I})$. It vanishes under the
remaining real functional $\Lambda$. On the cone where all $|q_i|$
are comparable and the nodal parameters $h$ are
$O(R_0\log(1/R_0))$, with $R_0=\max_i|q_i|$, the next term is
$\Lambda F=\sum_i H_i(q_i)+o(R_0)$. It is the residue of the
variation of the two sextic sheets on the quotient cubic.
Appendix~\ref{app:joined-estimates} proves the two estimates,
including uniformity at the faces where some $q_i$ vanish.

\subsection{Degree and orientation}
Choose $n\in P_I$ with $\Lambda(n)=1$. The boundary degree is that
of the real-linear map
\begin{equation}\label{p14-eq:linearjoin}
 (h,q)\longmapsto Jh+2\iota\big((\bar q_iB_i v_i)_{i\in I}\big)
                                      +n\sum_i H_i(q_i).
\end{equation}
Indeed first homotope by linear interpolation to this model with a
common factor $\log(1/R_0)$ on the matching term. Replacing the separate
logarithms there by the common one costs $O(R_0)$, uniformly up to
the coordinate faces. At a sequence of noncentral zeros
\eqref{p14-eq:fullestimate} and injectivity of $J$ would force
$|h|=O(R_0\log(1/R_0))$, $R_0>0$. Projecting away the nodal
columns, then dividing by $R_0\log(1/R_0)$, puts each limiting
tuple $(q_i/R_0)$ on the real line spanned by $(\beta_i)$. Indeed
the remaining matching space is $(\bigoplus L_i)/\mathbb R([e_i])$.
Every entry of a nonzero tuple on that line is nonzero. Hence the
comparable cone applies, and division of \eqref{p14-eq:rowestimate}
by $R_0$ contradicts $\sum_i w_i\ne0$. On the model the common
positive logarithmic factor can then be changed to one. This
establishes both isolation and the asserted boundary degree.

Restore the even height as a real variable, projecting to its
evaluation vector $E_I=(e_i)$ before taking the quotient in
\eqref{p14-eq:joinedsequence}. Eliminate $B_0$ and order the remaining
rows as the matching fibers followed by the complementary normalized
line. The remaining matrix is
\[
 \begin{pmatrix}M_I&E_I\\ H&0\end{pmatrix},\qquad
 M_i(p_i,q_i)=v_i p_i+2\bar q_i B_i v_i,\qquad
 H(p,q)=\sum_i H_i(q_i).
\]
Each $M_i$ has real sign $-1$ by complex linearity in $p_i$ and
antilinearity in $q_i$ modulo $\mathbb Cv_i$. The Schur complement is
$-H M_I^{-1}E_I=-\sum_iw_i$. The sign from cancellation of the
branch-and-time block, the identity operator, and the even height is
common to all $I$: the normalization determinant is the determinant
on the two smooth components times the dual determinant of the
complex matching fibers. Every added pair $(p_i,q_i)$ and matching
fiber has real dimension four. Moving these past a fixed real line
or past each other preserves orientation. The projective Euler link
has Euler number twice the oriented degree, since its cokernel
quotient is the pullback of the tangent bundle of an even-dimensional
sphere under the normalized column. Dividing by the order-two
stabilizer leaves the signed unit \eqref{p14-eq:localsubsetsign}.

\subsection{Regularity and prescription of the weights}
The marked cubic height difference and the derivative test in
Section~\ref{p14-sec:kernels} give generically $[e_i]\ne0$. Adding the
value of $B_i$ to the marked test is submersive: after forcing off
the marks the adjoint vanishes there, the local section values and
holomorphic jets kill the marked multipliers, and the zeroth-order
coefficient can then vary freely at the center. The failure of
$[B_iv_i]\ne0$ has real codimension two.

With these jets fixed, adjoin to the cubic branch-augmented operator
one real column whose residues at the unramified $y_i$, $i\in I$,
are $\beta_iv_i$. A cutoff representative with residue $R$ there is
$R/w-B_iR\log|w|^2$. The off-pole forcing is locally
$O(1+|\log |w||)$ and its pairing with the remaining cubic adjoint
is $-2\sum_iH_i(\beta_i)$. Quotient out the known smooth kernel.
The resulting index-zero operator is invertible exactly when the sum
is nonzero. At a defect a polar solution has nonzero coefficient on
the new column. Its three conjugates at any chosen $x_i$, pulled
back to $Z$, have an order-two pole at three different flags, while
the other conjugates there have at most order one. Together with
the bounded $s_1,s_2$ these give five independent solution germs.
The constrained Petri argument of Lemma~\ref{p14-lem:kernelplane}
therefore gives a rank-increasing perturbation at a defect. More
precisely, first insert a good Petri jet on the test disk, then
apply the annihilator argument there to the three polar columns and
three local adjoint components. This gives open density with fixed
graph; intersect over the finitely many subsets. Open conditions on
a continued cubic crossing transfer to generic paths by a time
cutoff. For the countably many kernel exclusions of
Lemma~\ref{p14-lem:kernelplane} the same argument is applied to the
individual compact configuration blocks before taking the intersection.

Finally we explain how to prescribe the chamber on the fixed cover in
Proposition~\ref{p14-prop:realize}. The branch model used there starts
with $B_i=0$. Insert any bounded antilinear values with
$[B_iv_i]\ne0$ by cutoffs on disks of radius $\rho\to0$. The
perturbation tends to zero as an $L^p$ multiplier on the fixed
covers for each finite $p>2$. To restore the equation on the same
$Y$ add $2b+1$ smooth real operator variations on the regular
frame disk. Their values on $\xi$ may be chosen to span the
cokernel of $D_f$, by trace-free forcing and adjoint continuation.
The implicit function theorem with a thick real normalization now
gives coefficients tending to zero and a new normalized solution
converging to $\xi$ in $W^{1,p}$. At ramification the pulled-back
cutoff coefficient has the form $2\bar z\chi(z^2/\rho)B(z^2)$.
It is uniformly Lipschitz and $O(\sqrt\rho)$, hence small in
$C^\alpha$, $0<\alpha<1$. Interior estimates give $C^1$ convergence
of the heights at the flags. In particular the even value stays off
the complex derivative line, the odd height is $v_i z+O(|z|^{1+\alpha})$
uniformly after choosing $\alpha<1-2/p$, and regular-sheet separation
persists up to each ramification point. The branch pole columns
converge in $L^p$ since the factor $\bar z$ cancels the conjugate
polar denominator. Consequently \eqref{p14-eq:lowertransverse} remains
invertible, and its branch-annihilating adjoint, normalized by time,
converges uniformly on the cubic by the adjoint estimates. Apply the
preceding regular-disk tests to ensure the singleton $H_i(\beta_i)$
are nonzero and the discriminant augmentation is invertible.

Given real $c_i$ with all nonempty subsums nonzero, now choose
$\gamma_i$ with $H_i(\gamma_i)=c_i$. Insert bounded coefficients
$\widetilde B_i$ satisfying
$2\overline{\gamma_i}[\widetilde B_i v_i]=[e_i]$, and restore the
fixed-cover equation in exactly the same way. The normalized cubic
and discriminant pole columns and the time column converge on their
fixed sources. Thus the remaining orientation choices transport
continuously and $\beta_i\to\gamma_i$, $w_i\to c_i$. For sufficiently
small supports all desired strict inequalities hold. Restoring the
residual kernel exclusions on the regular disk preserves them.
This proves both propositions, including realizability on a
prescribed $f$. For every even $b>0$ such cubics exist over the same
$C$: send the $b$ peripheral loops to $(12)$ and one handle
generator to $(123)$, all other handle generators to the identity.
The surface relation holds and compactifying gives the required
connected holomorphic cover.

\section{Contracted components and the degree-seven wall formula}\label{p14-sec:virtual}

Let $Z\to C$ be the sextic Galois closure. Put $b=2l$ for the
number of branch values and $h_j=1+7(g-1)+3l+j$.
Fix a nonempty set $I$ of these values. Choose the negative root
line of one inertia transposition $\tau$ in the standard representation.
Over each value in $I$ join the $\tau$-fixed point of $Z$ to the
point of the zero section. Denote the resulting finite map by $a_I^0$
and its signed local degree from Proposition~\ref{p14-prop:signs} by $a_I$.
Its arithmetic genus is $h_{|I|}$. We write $G_{I,k}$ for the locus
obtained from $a_I^0$ by attaching one-pointed contracted curves at
smooth points or inserting two-pointed contracted curves at its nodes,
with total arithmetic genus $k$.

\begin{proposition}\label{p14-prop:virtual-joins}
At a cubic crossing satisfying the kernel and smoothing hypotheses of
Sections~\ref{p14-sec:kernels} and~\ref{p14-sec:joins}, the contribution of
$G_{I,k}$ to the change of the degree-seven local invariant is
\[
       a_I [v^{2k-2}]\bigl(2\sin(v/2)\bigr)^{2h_{|I|}-2}.
\]
A single interval about the crossing suffices to compute the changes
for all $k$ on the same two sides.
\end{proposition}

\subsection{The central locus}
By Section~\ref{p14-sec:kernels}, a central degree-seven map carrying a
projective height has positive normalizations \(Z\) and \(C\), each
with multiplicity one. Its matching height on \(Z\) spans the negative
line of one transposition. Its embedded graph meets the zero section
once over each branch value, with distinct complex tangent lines.

Each connected contracted subcurve at a node or at a smooth point has
one image. Maximal such subcurves have at most two positive branches,
one on each normalization when there are two. The joining values are
therefore recovered from the map with its projective kernel.
For fixed $k$, $G_{I,k}$ is compact. Take stable degree-one maps to
the two normalizations with marks evaluated at the indicated flags.
Glue their marks directly or insert stable two-pointed curves of
positive genus. Summing these proper gluing maps over the finitely
many genus distributions gives the stated locus, including a tail
specializing to a bridge. Distinct joining sets have disjoint images
in the projective incidence. The unique matching height line is
spanned by the chosen root on $Z$ extended by zero on $C$.

\subsection{Displacing the line of sections}
For a stable map $a$ to $C$ write $L_{a,t}=a^*D_t$. Near its
one-dimensional kernel, take a finite-dimensional stabilization and
projectivize its domain over the reals. The zero on the tautological
line represents the local degeneracy class: the derivative in the
projective directions cancels the image directions of the linear map.
In the transverse matrix space, with a generator of the kernel chosen,
the cokernel on the linking sphere is its tangent bundle. The rank off
the central locus is $r=4(k-1-7(g-1))$, the real virtual dimension of
the moduli of maps to $C$. The Euler number of this even-dimensional
link is two. Thus the relative projective class capped with the
fixed-target virtual class times the oriented parameter interval,
and multiplied by two, computes the terminal minus initial Euler
number \cite[\S6.2, Lemma~6.10 and Appendix~B]{bai}.

Choose a closed disk strictly inside a larger disk of regular values.
On the central map the seven sheets over the larger disk are distinct
as pieces of the normalization, including possible contracted
attachments. For nearby maps the preimage of the smaller disk stays
in seven disjoint neighborhoods. Each has degree one by conservation
of degree at an ordinary value avoiding all attaching points. Order
these pieces using an identification of the sextic with the fixed
model over $C$ taking a chosen joining flag to its fixed representative.
There are two such identifications, interchanged by $\tau$. In the
finite covering of all sheet orderings these choices and their
complement are disjoint closed sets over the compact central locus.
Shrinking the neighborhood extends the selection to a stack cover of
degree two, retaining the automorphisms of contracted components.

On one distinguished nonzero sheet choose a real covector density on
the disk with integral $\ell_a(s)\ne0$ on the chosen root $s$.
Integration uses pulled-back area, so constant components give zero.
Normalize by $\ell_a(u)=1$. For independent smooth vector-valued
$(0,1)$-forms on the seven sheets, compactly supported in that disk,
write $j_a$ for their sheetwise pullback. Replace $L_{a,t}u$ by
$L_{a,t}u-\ell_a(u)j_a$. On the normalized slice the equation is
\begin{equation}\label{p14-eq:virtual-affine}
                      L_{a,t}u=j_a,\qquad \ell_a(u)=1 .
\end{equation}
This defines a continuous perturbation in nodal stabilizations.
Indeed introduce $z=\ell_a(u)$, solve $D_t p=j$ locally on each
disk containing a nodal value on each sheet, and replace $u$ by
$u'+za^*p$, using cutoffs equal to one near the nodes. The new forcing
vanishes there, and the scalar row becomes
$z=\ell_a(u')+z\ell_a(a^*p)$. These substitutions are invertible
row and variable changes. On thick parts the remaining forcing converges
smoothly, and on a shrinking disk the area test on bounded sections is
bounded by seven times its area times the supremum of the integrand.
These are precisely the continuous perturbations of the stabilized
inverse in \cite[Appendix~B]{bai} and
\cite[Propositions~B.7.9 and B.10.4]{pardonatlas}.
To obtain a single ambient almost complex structure, split the tangent
bundle of the total space by a complex connection and write
\(D_t=\nabla^{0,1}+A_t\). Over the support disk take disjoint cutoffs
\(\chi_m\) equal to one near the seven normalized sheet graphs.
Write \(F_{t,j}(x,n)=A_t(x)n-\sum_m\chi_m(x,n)j_m(x)\), a fiber-valued
antilinear one-form in the base direction. The block formula
\[
 (X,V)\longmapsto\big(j_CX,\ iV+2iF_{t,j}(x,n)X\big)
\]
squares to minus the identity. With
\(\bar\partial u=(\nabla u+i\nabla u\, j)/2\), its graph equation
near these sheets is \eqref{p14-eq:virtual-affine} before imposing the
scalar row. The construction is independent of the source, including
its contracted components. It extends the homogeneous graph
correspondence of \cite[Lemma~5.9]{bai}. For a compact completion first cut off \(A_t(x)n\) radially,
identically one on a tube containing the normalized central graphs
and their sufficiently small displacements. Use the connection
structure near the divisor at infinity in the projective bundle.
A closed coupling form positive on the projective fibers, plus a
sufficiently large multiple of the base area form, tames this compact
family. All comparisons take place in the smaller tube. In the
bundle tube the area of a closed graph is bounded by its projection
degree times a constant.

On a finite-map chart the sheetwise forms surject onto the cokernel:
a dual solution annihilating them vanishes on an open set of each
positive normalization, and hence everywhere off its nodal flags by
unique continuation. The equation-space dual in the $L^p$ formulation
has no point masses. Sections with prescribed independent values at
the flags then eliminate the dual node-matching rows. Finitely many
forms suffice near the finitely many bare joins. On each smooth
stratum of the finite-map chart the displaced problem is transverse
for a regular choice. A persistent node or a collision of simple
branch values has positive complex codimension, whereas the full
normalized problem with time has index zero. Hence the displaced
zeros are finite, smooth, simply branched covers. Their graphs may
also be taken embedded. A double point over a common value imposes
four real conditions with at most two real parameters for that value.
At a simple ramification point a zero vertical derivative imposes four
conditions. For transversality of these marked problems the same
adjoint test vanishes off the marks. Prescribe arbitrary independent
values and holomorphic first jets there, solving the symbol equation
recursively for the antiholomorphic derivatives of a formal jet.
Testing against these jets eliminates the distributions at the marks.
Immersion excludes the diagonal in the marked-pair test, and finitely
many further forms then suffice by compactness. By conservation of
the local degree their signs in the labelled chart sum to $a_I$.
Here and below all displacements at the finitely many joins are made
using the same ordered disk and orientations.

\subsection{Isolation and contracted components}
Consider a convergent sequence of normalized graph solutions
approaching $G_{I,k}$, with displacements tending to zero. Take stable
limits of the positive normalizations separately, first retaining the
attaching marks. Their nonconstant limits together are the root graph
and the zero section, each once. Every positive component has positive
projection degree, so there are at most two such normalizations.
On forgetting marks, their individual stable limits have no contracted
component. Indeed the theorem of Ekholm and Shende
\cite[Theorem~1.1]{ekholm} would require a vanishing derivative,
equal complex tangent lines at two preimages, or three attaching
preimages of one point. The two central embeddings exclude each
possibility. The theorem applies to these bare maps because the
displaced equation is an ordinary graph equation. For finitely many
smooth perturbation parameters, put the varying structures and the
standard complex structure of an even-dimensional parameter ball
(adjoining one dummy coordinate if needed) in block diagonal form
on the product. Maps in constant slices are holomorphic for this one
structure, with the same attaching tangents.

If there are two positive normalizations, each converges smoothly
to one of the embeddings. Each nearby pair can meet at most once
near each central intersection: projecting the difference of the two
parametrizations to the real four-plane of the limiting tangents has
invertible derivative in the two local source parameters. Away from
these neighborhoods the limiting images are disjoint. Deleting the
one-ended contracted blocks and replacing the two-ended ones by nodes
now gives a finite map satisfying the same equation. Its joining
flags remain distinct. Properness of the marked factors shows that
their limiting subset is $I$. At zero displacement the local isolation
at each bare join forces precisely the central finite map. A regular
displacement has no finite nodal zero there.

With one positive normalization its unmarked limit has direct joins
at a nonempty set $J$. The bare isolation for $J$ again excludes a
sequence of smooth finite zeros at zero displacement. At a regular
displacement these zeros belong to the smooth embedded graphs just
constructed. Each remaining contracted block on such a graph attaches
once. On restoring marks to its limit one can insert rational trees
at smooth points or at existing nodes, and the one-pointed tails
can specialize to these trees. Distinct smooth points over an unjoined
intersection have disjoint fibers under marked stabilization. Hence
no additional joining value results: $J=I$. Conversely, every
degree-one expansion of a displaced graph solves
\eqref{p14-eq:virtual-affine}, and any of its limits at zero displacement
lies in $G_{I,k}$ by the marked description. This proves fixed-genus
isolation and identifies its entire displaced zero set. In particular
a bridge in a limit is counted with its own joining subset.

\subsection{The virtual class at a displaced graph}
Write $\Gamma$ for one of the embedded displaced graphs and $t_0$ for
its parameter. On its smooth projection the undisplaced homogeneous
operator is injective by isolation in the finite-map chart. This
remains true on each degree-one expansion by restriction and matching.
In a chart for maps to $C$ choose stabilizing marks, incidence slices,
horizontal perturbations and obstruction section $s_C$. Use the same
data for maps to the total space, lifting the slices by inverse
image and the horizontal forms by a complex-linear splitting of
projection. The splitting may be chosen affine in height near each
sheet by taking the complex-linear part of a connection lift.
Choose an equivariant supplement $BW$ to the image of $L_{a,t}$ at
the chart center, with representatives on the smooth parts of all
components including contracted ones. On $s_C=0$ the stabilized
vertical equation is
\[
                         L_{a,t}u-Bw=j_a .
\]
Its derivative in $(u,w)$ is invertible. Projection identifies the
regular simultaneous thickening with the horizontal one over the common
family of stabilized domains \cite[Theorem~3.5 and Lemma~5.10]{swam}.
The residual ambient obstruction rows are $s_C$ and $w$. For the
homogeneous projective equation introduce $z=\ell_a(v)$ as before.
Its stabilized solution is $(v,w')=z(u,w)$ and injectivity forces
$z\ne0$ at a projective zero. In the chart $z=1$ the joint rows are
therefore
\[
                       s_C=0,\qquad w=0,\qquad\ell_a(u)-1=0 .
\]
Their Thom product computes both the compactly supported projective
pairing and the constrained ambient class. To compare charts, take
the union of stabilizing marks and direct sums of perturbations.
Adding a vertical supplement pairs its new variable with the same
new obstruction coordinate. Changing horizontal lifts adds rows
vanishing on $s_C=0$, and eliminating a different vertical
supplement adds a multiple of $w$ to the scalar row. These changes
give the same localized Thom product on common refinements
\cite[\S9.2]{pardonatlas}. Both $r$ and the horizontal virtual
dimension are even. The scalar orientation fixed on the smooth
finite-map chart thus applies at every arithmetic genus.

In unparametrized normal coordinates at $\Gamma$ the augmented
linearization with time and the area test is
$(\xi,\dot t)\mapsto(A\xi+B_0\dot t,\lambda\xi)$.
It is invertible by transversality of the normalized finite-map
zero: the horizontal deformation complex of a smooth cover has
torsion cokernel sheaf and no obstruction. The fixed-time operator
$A$ has index zero. If it has a kernel, the augmented isomorphism
makes both kernel and cokernel lines. In an unramified disk where
their generators are nonzero, add a small zeroth-order term linear in
the height minus the height of $\Gamma$. The vertical tangent there
is the normal quotient, so the variation can have nonzero pairing
on these lines. This fixes the graph and makes $A$ invertible,
preserving the augmented isomorphism. In the affine presentation this
changes the homogeneous operator and the forcing by the corresponding
linear and constant terms. Compactness of the degree-one expansions
permits this homotopy inside the isolating chart. Any extra constrained
graphs approaching the retained one would have one nonconstant
normalization: in its retracting tube each positive component projects
in a positive multiple of seven. Ghost censorship at the embedded
multiplicity-one limit gives smooth convergence of the bare
normalization, contrary to uniqueness for the augmented equation.

The nearby ordinary graphs form a regular family $\Gamma_t$.
The same normalization and compactness argument shows that its
fixed-genus locus in a small tube consists of degree-one expansions.
On these maps the area test is the function
$\gamma(t)=\ell(\Gamma_t)-1$, with $\gamma'(t_0)\ne0$ by the
augmented isomorphism. Replacing the scalar row by $\gamma$ preserves
its Thom product with the two other rows since the functions agree
on their zero set. The time slice contributes
$\operatorname{sgn}\gamma'(t_0)$ times the ordinary fixed-time
degree-one class. The normal bundle of $\Gamma$ has degree
$2h_{|I|}-2$, since $c_1(TC)+c_1(N)=0$. Its regular normal
operator pulls back injectively to every degree-one map. Zinger's
comparison identifies the local invariant with the normal Euler class
on the moduli of degree-one maps \cite[Theorem~1.2]{zinger}.
To transport the sign to a complex-linear reference, quotient the
pullback Fredholm complex by the complex on the base curve $\Gamma$.
Restriction to the unique positive normalization splits the pullback
on cokernels, and matching identifies the two kernels throughout any
deformation. The quotient thus has vector-bundle cokernel with its
quotient determinant orientation \cite[Appendix~B]{bai}.
At an invertible end the remaining orientation is the base determinant
sign. At a Dolbeault reference on two non-effective degree
$h_{|I|}-1$ lines the Hodge integral gives the series
$(2\sin(v/2))^{2h_{|I|}-2}$ \cite[Theorem~2]{pandhodge}.
The product of the time and normal signs is the sign of the original
displaced zero, by evaluating in its own smooth genus. Summing these
signs and cancelling the factor two from the Euler link against the
degree two of the labelling proves the asserted coefficient at fixed
$k$. Isotropy on contracted curves is included in the degree-one
virtual class.

\subsection{One interval for exact degree seven}
We record the genus estimate needed to pass to a common interval. For
a compact smooth family of homogeneous normal operators on $C$ there
is a constant $K\ge0$ such that a simple graph projecting in degree
$d$ has normalization genus at most $1+d(g-1)+Kd(d-1)$. Indeed, take
the normalization of the off-diagonal part of the ordered fiber
square of its projection. The difference of the two heights is
generically nonzero on each component. For a fixed unitary connection
and $M$ a uniform norm bound on the zeroth-order term, it satisfies
$|\bar\partial_\nabla s|\le M|s|\,|df|$, up to a fixed metric
constant, for the induced projection $f$. Its isolated zeros have
positive integral orders by the similarity principle. Off these
zeros let $e=s/|s|$ and $\alpha=\operatorname{Im}\langle
e,\nabla e\rangle$. The curvature identity expresses $d\alpha$ as
the curvature evaluated on $e$ plus twice the difference of the squared
norms of the orthogonal $(1,0)$ and $(0,1)$ derivatives of $e$.
The latter $(0,1)$ norm is bounded by $M|df|$. Thus
$d\alpha\ge-f^*\kappa$ for one positive area form $\kappa$ on $C$
uniform over the family. Stokes on the complements of the zeros bounds
their total order by $d(d-1)\int_C\kappa/(2\pi)$.
A ramification point of index $a$ gives $a-1$ distinct off-diagonal
branches over the repeated point, all with zero difference. The genus
bound follows by Riemann--Hurwitz.

Suppose a connected stable map of exact degree $d$ has a nonzero
height. For its positive normalizations let $d_i$ be the degree of
the simple graph and $m_i$ the covering multiplicity, including graphs
of height zero. Thus $\sum_i m_i d_i=d$. Replace each covering by a
connected unramified cover of the simple normalization with the same
multiplicity. Such a cover exists since its base has genus at least
two. The new genus is at most
$1+m_i d_i(g-1)+K m_i d_i(d_i-1)$. Reconnect along a tree of equal
graph values, taken from the original matching after collapsing each
connected contracted block to its image. Lift each selected flag to
its replacement. If flags coincide there, place them on a constant
rational component attached at that point with one node for each of
the coincident flags. This keeps the dual graph a tree and the
resulting map stable. Its degree is $d$, its height is nonzero, and
its genus is at most $dg+Kd(d-1)$. The replacement is used only to
detect existence of a kernel.

Apply this test at $d=7$ on a compact subarc. In each of the finitely
many genera below the bound the central compact loci just described
exhaust all projective zeros. Fixed-genus isolation and compactness
give a kernel-free punctured interval. Intersecting the finitely many
intervals and applying the test proves the assertion in every genus
at exact degree seven. The Euler integrals are then constant on each
punctured side coefficient by coefficient. This finishes the proof of
Proposition~\ref{p14-prop:virtual-joins}. Since the only proper divisor
of seven is one, and the base operator stays invertible, the
degree-one factor has constant exponent. Subtraction of its unchanged
contribution in the connected series gives
$\Delta n_{7,h_j}=\sum_{|I|=j}a_I$ at super-rigid endpoints, with
all other degree-seven exponent changes zero. This proves Theorem~\ref{p14-thm:wall}.
\subsection{The possible changes in elementary multiplicity}\label{p14-sec:chambers}
Write \(c_j=\Delta n_{7,h_j}\). With a common choice of sign as in
Theorem~\ref{p14-thm:wall}, the set of attainable vectors on the given cubic
is the image of the complement of the real hyperplanes
\[
       H_I=\{w\in\RR^b:\, w_I=0\},\qquad
       w_I=\sum_{i\in I}w_i,\qquad \varnothing\ne I\subseteq E ,
\]
under \(w\mapsto ((-1)^j\sum_{|I|=j}\sgn w_I)_{1\leq j\leq b}\).
These hyperplanes form the resonance arrangement \cite{kuehne}. Several
chambers may determine the same vector. The sharp binomial bound used
in Theorem~\ref{p14-thm:torus} follows at once from this description:
\[
       |c_j|\leq\binom bj,\qquad
       \max_{w}\sum_{j=1}^{b}|c_j|=2^b-1 .
\]
Indeed choosing all weights positive gives equality for every \(j\)
simultaneously. Such a choice is realizable by
Proposition~\ref{p14-prop:realize} and the branch variations of
\S\ref{p14-sec:joins}.

\begin{proposition}\label{p14-prop:lattice}
The subgroup generated by the changes on a prescribed cubic is
\[
       2\ZZ^b+\ZZ\left(\binom b1,\ldots,\binom bb\right).
\]
Its index in \(\ZZ^b\) is \(2^{b-1}\).
\end{proposition}
\begin{proof}
Every summand has absolute value one, so the reduction of any
attainable vector modulo two is the displayed binomial vector.
For each nonempty \(I\) choose a point of \(H_I\) outside the finitely
many other hyperplanes. Crossing there changes only \(\sgn w_I\).
The difference of the two attainable vectors is a signed double
coordinate vector in coordinate \(|I|\). These differences give all of
\(2\ZZ^b\); the last binomial coefficient is one, so the reduction
of the generated subgroup has order two.
\end{proof}

There is a more restrictive relation on chambers meeting the full-sum
hyperplane \(H_E\) at a point on no proper-subset hyperplane.
Complementation pairs each proper subset \(I\) with a subset of
opposite sum on \(H_E\). Since \(b=2\ell\), the two cardinalities
have the same parity. On either side of this point it follows that
\[
             c_{b-j}=-c_j\quad(1\leq j<b),\qquad c_\ell=0 .
\]
Only \(c_b\) changes when crossing \(H_E\) there. Consequently neither
of these two chambers attains the unrestricted binomial bound.
For example, at \(b=2\) they give \((0,1)\) and \((0,-1)\). The
same-sign chambers give \((-2,1)\) and \((2,-1)\). Expanding
\(x=v^2-v^4/12+O(v^6)\) gives, for \(g=2\), the connected
genus-thirteen coefficient \(c_2-11c_1/12\). Thus the two chambers
with nonzero \(c_1\) give \(17/6\) and \(-17/6\).

The subgroup from chambers adjacent to \(H_E\) also has a simple
description. In the free coordinates \((c_1,\ldots,c_{\ell-1},c_b)\)
it is
\[
 2\ZZ^\ell+\ZZ\left(\binom b1,\ldots,\binom b{\ell-1},1\right).
\]
Indeed two proper-subset hyperplanes have the same restriction to
\(H_E\) only when their indexing sets agree or are complements:
their indicator vectors must be proportional modulo the constant vector.
Cross one generic restricted hyperplane in \(H_E\), and move both
endpoints slightly to the same side of \(H_E\). This gives twice any
coordinate with \(j<\ell\), up to sign. Crossing \(H_E\) itself gives
twice the last coordinate. Reduction modulo two finishes the argument
also for \(\ell=1\). Here the subgroup describes sums of changes over
several crossings; an individual crossing must satisfy the subset-sign
formula.

\begin{proof}[Proof of Theorem~\ref{thm:invisible}]
Such covers exist by the Riemann existence theorem: send the two peripheral loops to the same transposition and one handle generator to a three-cycle, with its partner trivial. The product relation and transitivity hold. The construction of \S\ref{p14-sec:kernels} realizes the crossing on any prescribed such cover. The standard two-plane of \(S_3\) is the central kernel on its normal closure. The fundamental group of an unramified cover, restricted to the punctured base, contains the normal closure of the peripheral loops. Its image therefore contains the normal closure of a transposition, which is \(S_3\). The two-plane has no invariants. Lemma~\ref{p14-lem:kernelplane} gives an invertible pullback on each unramified cover at the center. There are finitely many covers of the bounded degree, so their determinants remain invertible on a common interval. The smooth cubic calculation of \S\ref{p10-sec:incidence} gives the first change. Theorem~\ref{p14-thm:wall} gives the second from the single subset containing both branch values.
\end{proof}

\section{Smooth degree-seven covers near a change of sign}\label{sec:global}
Changing the sum of the weights at a fixed set of flags through zero reverses the degree of the local smoothing equation. Smooth maps near the join account for this difference.
\subsection{Smooth covers approaching a joined map}\label{p14-sec:smooth}

Fix a simply branched cubic $Y\to C$ with monodromy $S_3$ and
$b=2\ell>0$ branch values. Let $I$ be a nonempty subset of these
values. Write $j=|I|$, $h_j=1+7(g-1)+3\ell+j$, and $f_I$ for the
finite map obtained by joining its Galois closure $Z$ to $C$ at
the selected flags fixed by one transposition. We use increasing
$t$ as the positive direction across a cubic wall. A crossing on a
smooth cover is counted once for each isomorphism class over $C$ and
each parameter, with the sign of the derivative of the normalized
kernel equation in the complex-oriented Hurwitz directions and the
oriented path direction.

\begin{theorem}\label{p14-thm:smooth-covers}
On the prescribed cubic and normal bundle there is a two-parameter
family of operators $D_{a,t}$ with the following properties.
Along $t=0$ the prescribed cubic has a normal kernel, transverse in
the $t$ direction. There is a sequence of parameters approaching a point
of this segment for which smooth connected degree-seven covers with a
nonzero normal section converge to $f_I$ and its root section. Their source genus is $h_j$ and their monodromy
is $S_7$. If $I$ consists of all branch values, the graphs in the
convergent sequence are embedded after discarding finitely many terms.

There are also arbitrarily small Hurwitz neighborhoods $U$ of $f_I$
and rectangles $[a_-,a_+]\times[-\delta,\delta]$ as above with the
following signed count. Let $d_L,d_R$ denote the local degrees in the
uniformizing chart of the joined map along the left and right edges.
On the lower and upper edges, oriented by increasing $a$, take
sufficiently small generic perturbations relative to their ends.
Their smooth crossings in $U$ are regular, have simple branching over
distinct values, trivial deck group and embedded kernel graph. Their
signed counts $N_-,N_+$ satisfy
\begin{equation}\label{p14-eq:smooth-edge}
                 2(N_+-N_-)=d_R-d_L\in\{-2,2\}.
\end{equation}
\end{theorem}

\begin{proof}
\emph{Changing the degree.}
Take a short segment $c(a)\in\mathbb R^b$ crossing the hyperplane
$\sum_{i\in I}c_i=0$ at a point of no other subset hyperplane.
In the branch coordinates of Section~\ref{p14-sec:joins}, $H_i$ is
a nonzero real functional on a complex smoothing plane. Choose
nonzero $\gamma_i(a)$ smoothly with $H_i(\gamma_i(a))=c_i(a)$.
A right inverse of $H_i$ plus a fixed nonzero vector in its kernel
gives such a choice even when $c_i=0$. Assign the antilinear
coefficient at the $i$th center so that
$2\overline{\gamma_i(a)}[B_iv_i]=[e_i]$.
The local replacement and restoration on a disk of regular values in
Section~\ref{p14-sec:joins} apply with one cutoff radius for this compact
segment. Indeed the prescribed coefficients and the right inverses
in the restoration are uniformly bounded. Thus the weights on the
restored cubic are uniformly close to $c_i(a)$. Take the radius
small enough to retain opposite signs of the $I$-sum at the two ends
and to keep every other subsum away from zero. The ordinary and polar
operators on the fixed covers retain the ranks used in the nodal
derivative. The cubic remains transverse to one common time direction.
The kernel at $f_I$ remains the odd real line: evaluation at any
chosen flag removes the even section on $Z$.

Write $q_i$ for the complex smoothing parameters, and $h$ for the
remaining chart coordinates and time, centered at $f_I$, $t=0$.
Trivialize the reduced obstruction equation over the parameter
segment by stabilizing on disks of regular values and normalizing the
root section by a real linear functional. The center kernel and
cokernel have constant rank, so supplements and normalizations can be
continued along the segment. By Proposition~\ref{p14-prop:signs} the end
zeros are isolated and have opposite degrees in this uniformizing
chart. On the boundary of every sufficiently small ball in $(h,q)$
a zero occurs at some intermediate $a$, by homotopy invariance of
degree. Passing to a sequence of shrinking balls and a convergent
subsequence of parameters, put $r=\max_{i\in I}|q_i|$. The nodal
derivative is uniformly injective, hence $r>0$ at every noncentral
zero sufficiently close to the center and $|h|=O(r\log(1/r))$.
After division by $r\log(1/r)$ the matching equation of
Section~\ref{p14-sec:joins} places every subsequential limit of $(q_i/r)$
on the real line spanned by $(\beta_i)_{i\in I}$ at the limiting
parameter. Here $2\overline{\beta_i}[B_i v_i]=[e_i]$. Every coordinate
of the generator is nonzero. Thus all $|q_i|$ are comparable to $r$.
The remaining adjoint row, divided by $r$, now gives
$\sum_{i\in I}H_i(\beta_i)=0$ at the limit. These estimates are
uniform on the segment, since they use the bounds for the stabilized
inverse and the leading nodal and matching maps, all of which retain
their ranks at resonance. Every selected node has smoothed. The
domain is connected, of arithmetic genus $h_j$, and the normalized
sections converge to the odd section. This proves existence of the
asserted sequence.

\emph{Ramification and separation.}
On the three local sheets involved at a selected flag, the projection
in the neck coordinate $z$ is
\[
              w-c_i'=(z+p_i)^2+q_i/z .
\]
Here $c_i'$ is the displacement of the branch center.
Both $c_i'$ and $p_i$ are $O(r\log(1/r))$. Put $d_i^3=q_i/2$.
The critical equation is $2z^2(z+p_i)=q_i$, with three distinct
simple roots $z/d_i=\zeta+o(1)$, $\zeta^3=1$. At such a root,
\[
        w-c_i'=(z+p_i)(3z+p_i)=3d_i^2(\zeta^2+o(1)).
\]
These critical values have mutual distances comparable to $r^{2/3}$.
The two other local quadratic factors have critical values at
distance $O(r\log(1/r))$ from the original branch value. Consequently
each of the three displayed values supports exactly one simple
ramification point. At each unselected branch disk the three
quadratic factors persist. In total the ramification degree is
$B=3b+2j$. A solitary simple branch gives a transposition on the
seven sheets. The graph of the conjugates of that transposition has
components permuted transitively by the monodromy, so their common
size divides seven. There is an edge, hence this graph is connected
and its edge transpositions generate $S_7$.

For the full join the two limiting images coincide only at the
selected points. On a selected disk the normalized height on the
cubic factor has the expression $U_0(w)+zU_1(w)+z^2U_2(w)$. In the
fixed-base polynomial presentation its coefficients converge in
$W^{1,p}$ for some $2<p<4$, hence uniformly, by the stabilized
inverse in Section~\ref{p14-sec:joins}. Centrally $U_1(0)=v_i\ne0$ and
$U_2=0$. The difference of the heights at distinct roots $z,z'$
over the same $w$, divided by $z-z'$, is
\[
                          U_1(w)+(z+z')U_2(w).
\]
Its norm is bounded below on a smaller disk. At a simple critical
point with local coordinate $w=w_0+\eta^2$, the smooth height has
zero antiholomorphic derivative at $\eta=0$, by the pullback
equation. Taking the difference at $\eta$ and $-\eta$ gives vertical
derivative in the $z$ coordinate equal to
$U_1(w_0)+2z_0U_2(w_0)\ne0$. Thus this factor also immerses at
ramification. The other two quadratic factors stay separated by
their distinct nonzero limiting heights. Their normalized heights
converge in $C^1$ on smaller source disks by ordinary elliptic
regularity and have nonzero limiting vertical derivative at the
ramification points. Away from these disks the projections are
unramified and the central heights are pairwise separated on a
compact set. The resulting injective immersion of a compact source is
an embedding.

\emph{Isolating a neighborhood for paths of operators.}
To obtain signed counts, the neighborhood must lie in Hurwitz space
alone, independent of $a$ and $t$. Arrange the infinite-order Petri
condition at two distinct points of a regular disk where two cubic
heights form a complex frame, by the local gauge construction of
Section~\ref{p14-sec:kernels} and \cite[Corollary~5.10]{wendl}. Restore
the cubic, and choose its transverse time direction, on another
subdisk. Both Petri jets are then fixed along the family. Vary the
central segment relative to its ends on the frame disk by small
$a$-dependent tensors flat at the test points and annihilating the two
current cubic heights. This preserves the prescribed cubic.

For connected smooth covers sufficiently close to $f_I$, monodromy
outside the branch disks contains the regular $S_3$ on six sheets
fixing the seventh. Connectedness now implies double transitivity.
The reduced real sheet representation $W$ is absolutely irreducible
of dimension six: the complex commuting algebra of the permutation
representation has dimension two by the two orbits on ordered pairs.
The constant sector is invertible. The six components of a nonzero
solution in $W$ and the two independent cubic heights are linearly
independent as real solution germs on a regular disk. Indeed on the
joint normal cover they span inequivalent irreducible systems, and
unique continuation identifies their linear relations on the disk.

Here is the constrained transversality argument. An adjoint
annihilating the allowed variations gives on the disk a tensor
identity
\[
           \sum_{k=1}^{6}u_k\otimes v_k
                    =s_1\otimes\alpha_1+s_2\otimes\alpha_2,
\]
where $u_k,v_k$ are the test solution and adjoint components and the
$s_i$ are the prescribed cubic heights. The smooth frame gives the
$\alpha_i$ smoothly across the test points. Pair with $[D,\Phi]$ for
compactly supported complex-linear endomorphisms $\Phi$.
The left integral vanishes by the solution and adjoint equations.
Since $\Phi s_1,\Phi s_2$ can be chosen freely, both $\alpha_i$ solve
the adjoint equation. The Petri condition on infinite jets and
independence of the eight solution germs force the jets of all
$v_k$ to vanish. Unique continuation gives $v_k=0$. This is the
local-system Petri argument of \cite[\S\S1.6, 2.7--2.9]{dwbn}
with the two prescribed solutions retained.

A smooth branching profile with $s$ distinct values has complex
dimension $s$. The pullback operator has real index $-2B$ by
Riemann--Roch. After allowing $a$ and dividing by real kernel
scalings the expected dimension is $2s-2B$. The preceding test
gives transversality for a generic constrained central segment.
To apply it on the noncompact smooth strata, exhaust their
projective incidences by countably many compact subsets. Finitely
many supported variations suffice near the zeros on each such set,
and their parameter cutoffs preserve the end collars. The regularity
conditions on compact subsets are open, so the countable
intersection gives regular segments. Only simple branching at
distinct values remains, in dimension zero. The uniform matching
estimate above excludes partially smoothed zeros sufficiently near
the center throughout the segment. Thus all off-center zeros at
$t=0$ sufficiently near $f_I$ form a countable set.

Choose arbitrarily small invariant Hurwitz balls $U$ with radii
outside the set of radii of these zeros. Compactness gives a
$\delta>0$ for which $\partial U$ is free of kernels throughout the
rectangle. The uniform estimates exclude all nonsmooth zeros on
the horizontal faces after shrinking the ball and $\delta$, since
a zero with $q=0$ can only be the center at $t=0$.
The sign condition at each end leaves only the joined zero on each
vertical face. Small perturbations of the horizontal paths preserve
these boundary properties. Applying the preceding transversality
test there with unrestricted operator variations gives regular
zeros, with distinct simple branch values and one-dimensional
kernel. There is a transposition, hence full monodromy by the prime
degree argument. The deck group, the centralizer in $S_7$, is trivial.

For completeness, the marked tests for embeddedness are also
transverse. At two distinct preimages over a common value, equality
of heights has real codimension four, while a free base value has
two real parameters and a branch value is already determined.
Vanishing of the vertical complex derivative at a ramification
point gives another four real conditions. For either test a
putative annihilator solves its adjoint equation away from the
marked fiber. Its invariant part vanishes distributionally by base
invertibility: each test vanishes on sections pulled back from the
base. One of the two Petri points lies away from the marked value.
The unrestricted adjoint test there and unique continuation kill the
reduced part away from the marks. To test the remaining point-supported
distributions, vary a local section with prescribed value or
holomorphic first derivative. Its equation can be made zero to
any given finite order, since the coefficients involving a
$\bar\partial$ derivative are determined recursively by its pure
holomorphic jet. These variations kill the marked multipliers,
and the remaining homogeneous adjoint vanishes by elliptic regularity.
At a zero of a marked test it vanishes also on the infinitesimal
rescaling. The resulting negative expected dimensions exclude both
coincidences and nonimmersions. The signed zero sets on the compact
horizontal faces are therefore finite and have embedded graphs.

\emph{Boundary orientations.}
Orient the parameter rectangle by $da\wedge dt$. On the product of
its boundary with the uniformizing Hurwitz ball, the total signed
degree is zero because the section extends over the product and has
no zeros on the Hurwitz boundary. The right and left faces contribute
$d_R$ and $-d_L$. On each horizontal face a smooth zero has two
lifts, interchanged by the central involution. Their signs agree.
Indeed this involution acts complex linearly on the Hurwitz chart.
On the normalized obstruction target it preserves orientation,
as follows by equivariance from the nonzero local degree at either
vertical end. This orientation character persists across the
rectangle. The lower and upper faces thus contribute $2N_-$ and
$-2N_+$, respectively. This proves \eqref{p14-eq:smooth-edge}.
\end{proof}

When $I$ contains every branch value, complementary nonempty proper
subsets have opposite sums at the limiting weights. In particular the
size-$b/2$ terms cancel in either adjacent chamber. Weights of one
sign give the largest sum of absolute changes over all genera and
belong to different chambers. Equation~\eqref{p14-eq:smooth-edge}
determines the difference of the two counts in $U$. Extending that
identity to counts in the whole Hurwitz space would require control
of its complement.

\section{Compensation in a fixed integral homology class}\label{p14-sec:compact}
Let \((X,\omega)\) be a closed symplectic six-manifold with
\(c_1(TX)=0\). At an endpoint \(J\) of a smooth path of
\(\omega\)-tame structures, mark finitely many disjoint embedded
curves \(C_i\) in nonzero integral classes \(A_i\). Suppose each admits
arbitrarily small tubes in which the only cycles below a prescribed energy
bound are its multiples. Suppose also that each marked cycle series
has a geometric factorization into elementary clusters. These
hypotheses hold for the super-rigid homogeneous germs used in
Theorem~\ref{p14-thm:wall}: projection expresses each irreducible compact
curve in the bundle as a finite map with a height section, and a
vertical fiber contains no compact curve. The cluster construction
below then applies to the insertion consisting of the multiple zero
sections at each energy bound.

Let \(H\) be the virtual-dimension-zero, Chern-number-zero part of
Pardon's rational algebra of complex curves \cite{pardon}.
Write \(S\) for its almost-complex cycle-space counterpart and
\(\Phi:H\to S\) for the comparison of \cite[\S2]{WueProduct}. In \(S\)
a compact open and closed part of the relative cycle space defines an
insertion. Excision on disjoint supports gives multiplication, and
the two restrictions of such an insertion over a closed interval
agree. Denote by \(F_h\) the normalized balanced factor of a genus-\(h\)
elementary curve: it is the equivariant series on the sum of two
degree-\(h-1\) line bundles with product the canonical bundle and
circle weights \(+1,-1\). Its linear term is the multiplicity-one
class \(x_{h,1}\) and its image under \(\Phi\) is the elementary
cluster series \cite[\S\S3--4]{WueProduct}. We use the full integral
label \(\beta\in H_2(X;\ZZ)\) in the monomial \(Q^\beta\).
Products at a fixed energy cutoff are taken modulo terms above that
cutoff. For each marked curve send the local variable \(q^d\) to
\(Q^{dA_i}\), and let \(P_J\) be the product of its marked factors
over all the marked curves.

\begin{proposition}\label{p14-prop:exterior}
At each such endpoint there is a product \(R_J\) of elementary factors
with integral exponents, obtained from the connected supports outside
the marked sectors after a generic perturbation. It is independent
of the perturbation, the cluster decomposition, and sufficiently
small isolating tubes. For endpoints of a smooth tame path,
\begin{equation}\label{p14-eq:conservation}
                    P_{J_-}R_{J_-}=P_{J_+}R_{J_+}.
\end{equation}
The equality holds in the class-labelled completion with coefficients
in \(H\). It holds also when the marked curves differ at the two ends.
\end{proposition}
\begin{proof}
First take an integral taming form and fix an energy cutoff \(L\).
Choose smaller disjoint tubes \(V_i\) with closure in the respective
isolation neighborhoods. At the endpoint no nonempty cycle contained
in \(\overline V_i\) meets its boundary. This persists for nearby
structures: a contrary sequence would have a bounded-mass subsequence
converging as cycles with Hausdorff convergence of supports
\cite[Theorem~2.11]{diw}. A nonempty limit has positive energy, and
integrality prevents a sequence of positive cycles from limiting to
the empty cycle. Thus containment in each tube is open and closed
in the bounded-energy cycle space over a short path to a generic
endpoint perturbation.

At the perturbed endpoint distinct irreducible supports are disjoint.
The bounded-energy space of connected supports, counted with possible
multiplicities in the compactification, splits into compact open and
closed parts according to containment in the marked tubes and their
complement. To see closedness, a limit of connected supports is
connected, and disjointness of distinct irreducible curves forces it to
have a single irreducible support. The respective unions of supports
are disjoint compact subsets of \(X\). Restriction to disjoint open
neighborhoods splits each effective current into currents labelled by
the parts. This factors the entire insertion in \(S\), with mixed
configurations included by multiplying the part series. In particular
\(R_J\) is constructed from configurations all of whose connected
supports are complementary.

Subdivide each part by a finite cluster partition. Cluster isotopy,
excision and the fold relation decompose each cluster into a signed
elementary factor and corrections in classes \(d\gamma\), \(d\geq2\),
if the original class is \(\gamma\)
\cite[Propositions~3.18, 3.21 and Theorem~3.25]{diw}
\cite[\S4]{WueProduct}. The latter reference gives the identities of
cycle insertions needed here. Each step decreases
\(\lfloor L/\omega(\gamma)\rfloor\). This yields finite products for
the parts with integral exponents. The image of their product under
\(\Phi\) is precisely the insertion of the entire target at this
cutoff. For each integral class the entire class locus is compact
over the original path by cycle compactness; the class is locally
constant. The interval relation identifies its two endpoint
insertions. The argument requires isolation of the marked sectors
only at the endpoints.

To prove uniqueness, cancel factors in order of increasing positive
energy. At a first differing label \(\beta\), the difference of
the comparison images is a finite sum
\(\sum_h d_h\Phi(x_{h,1})\). All nonlinear terms at this label are
composed of smaller-energy contributions. Other labels at the same
energy have distinct monomials, also when they differ by torsion.
Independence of the multiplicity-one classes
\cite[Theorem~5.1]{WueProduct} gives \(d_h=0\) for every \(h\). This
proves equality of the two total products, uniqueness at an endpoint
and consistency on increasing the cutoff. Applying the same argument
in each isolating region fixes the marked factors. Equation
\eqref{p14-eq:conservation} and independence of the exterior follow.

For an arbitrary taming form, approximate its cohomology class on
the compact path by a rational class represented by a nearby closed
taming form and rescale the latter to integral periods. The energies
are uniformly comparable on the complex lines along this path. At a
fixed endpoint let \(M_J\) be the monoid of classes of effective cycles.
Any nonzero exponent of the full product has label in \(M_J\):
induct on positive integral energy. At an ineffective class the entire
insertion is zero and smaller effective labels have sums in \(M_J\).
Multiplicity-one detection forces all the new exponents there to
vanish. This also applies to the marked products in their tubes.
Cycle compactness and local constancy of class imply finiteness of
effective labels at bounded energy. Two choices of taming form
therefore give products in a common locally finite completion at this
endpoint, since their energies are comparable on \(M_J\).
Their images coincide at each label and the same induction identifies
their exponents. Subtracting the marked multiplicities identifies the
complementary ones as well.
\end{proof}

For example, in class $7A$ over the genus-two core with two branch
values the change in the marked Euler integral of source genus
thirteen is
\[
 a_{\{1,2\}}-\tfrac{11}{12}(a_{\{1\}}+a_{\{2\}}).
\]
Indeed $x^{11}=v^{22}-(11/12)v^{24}+O(v^{26})$, whereas
$x^{12}=v^{24}+O(v^{26})$. In the elementary product the two genera
give separate integral changes. The exterior multiplicities in
genera twelve and thirteen have the opposite changes, respectively,
by Proposition~\ref{p14-prop:exterior}.

\begin{proof}[Proof of Theorem~\ref{p14-thm:torus}]
On \(X=\RR^6/\ZZ^6\) with its standard symplectic form there exists a
constant tame linear complex structure admitting no nonempty compact
effective curve. To see this, identify the integral classes with
integral bivectors. For every nonzero integral bivector \(B\) the
equation \((\Lambda^2J)B=B\) defines a proper analytic subset of the
space of positive linear complex structures. Indeed this space is the
connected homogeneous space \(\mathrm{GL}^+(6,\RR)/\mathrm{GL}(3,\CC)\).
If the equation held on an open subset, the standard \(J\) would fix
the span of the entire \(\mathrm{GL}^+\)-orbit of \(B\) in the exterior
square. A diagonal matrix with distinct pairwise products of its
diagonal entries extracts the nonzero coordinate bivectors by
polynomial interpolation; oriented signed permutations then give every
basis bivector in the span. The standard \(J\) does not fix a wedge
of real coordinate vectors from two different complex lines.
Baire's theorem applied to the open tame subset gives a constant
\(J\) avoiding these equations for all \(B\ne0\). The structure descends
under translations. Every effective current would give an invariant
integral bivector by integration of its complex tangent planes, with
positive symplectic area. Hence the total series is one there.
The space of tame structures is path connected, so
Proposition~\ref{p14-prop:exterior} gives \(R_J=P_J^{-1}\) at every marked
endpoint on this torus.

In a product of two symplectic two-tori smooth the positive transverse
intersections of one horizontal torus with \(g-1\) vertical tori,
using \(zw=\delta\) in disjoint coordinate balls. The intersection
graph is a tree. The result is a genus-\(g\) symplectic surface of
class \([T_1]+(g-1)[T_2]\). Take its product with a point of the last
two-torus. The resulting class \(A\) is primitive by pairing with
the area class of the first factor. Its symplectic normal bundle has
Chern number \(2g-2\); unitary bundles of rank two over a surface are
classified by this number. A tubular identification thus realizes the
given homogeneous structure with its prescribed normal derivative.
It is tame near the zero section. Cutting off in a chart of tame
structures extends it over \(X\) while preserving the germ. At every
positive even \(b\) a connected cubic as required exists on the same
base: send all puncture loops to \((12)\), one handle generator to
\((123)\), and the rest to the identity, and fill the finite cover
over the punctures. Applying Theorem~\ref{p14-thm:wall} changes the
exterior multiplicities by \(-c_j\) at \(h_j\), and by zero in every
other genus at label \(7A\). The binomial bound of
\S\ref{p14-sec:chambers} proves sharpness.
\end{proof}

In particular at least one of the two endpoints attaining this
maximum has \(\sum_h|n^{\rm ext}_{7A,h}|\geq2^{b-1}\), by the
triangle inequality and integrality. The area and the class are fixed
as \(b\) increases. These endpoints may approach a smooth reference
structure whose marked section is itself super-rigid. Use, for example,
the homogeneous split spin model of \cite[Lemma~14.1]{WueUnramified} on the
same bundle. All the homogeneous structures agree with it along
\(C\). On a tube of radius \(R\), cutting off their difference in a
smooth chart gives sup norm at most \(K_b R\) and first derivative
bounded by \(K_b\); the constant may depend on the crossing. For each
\(b\) take \(R_b\) so small that \(K_b R_b^{1/b}\leq 1/b\).
Interpolation yields convergence in every \(C^{0,\alpha}\),
\(0<\alpha<1\). Real dilations of the full-join graphs in
\S\ref{p14-sec:smooth} likewise place their embedded representatives
in arbitrarily small tubes, using the operators of that construction.
Their genus is \(h_b=1+7(g-1)+5b/2\). It is unbounded along this
sequence of smooth tame structures at fixed symplectic area.

\smallskip
The compensation determines the total signed change outside the
isolating tube at the two endpoints. Allocating that change to the
individual intermediate curves of an arbitrary compact deformation
remains open. For simultaneous ramified crossings there is a further
local question: several independent heights can match on a common
component, and the resulting strata need not have a common
dimension or stabilizer. The single adjoint equation of
\S\ref{p14-sec:joins} treats the prescribed cubic crossing.

\appendix
\section{The double-cover displacement and its elliptic tails}\label{sec:smooth-crossings}
We prove the double-cover displacement used in \S\ref{p10-sec:incidence} and separate its first contracted elliptic contribution into the moving and ramification-point terms.
\subsection{The local contribution at a double-cover wall}\label{p9-sec:wall}

Let $C$ be a smooth curve of genus $g$, and let $N\to C$ be a complex
bundle of rank two and degree $2g-2$. Consider a smooth path $D_t$ of
real Cauchy--Riemann operators on $N$, with $D_0$ invertible. We orient
the parameter interval by increasing $t$. Write $\mathcal H$ for the
space of smooth connected double covers of the fixed target, locally
taken in a uniformizing chart. At a cover $f:Y\to C$ with $b$ branch
values, these values are local complex coordinates up to ordering.
The deck involution acts trivially on these coordinates.

Suppose $K=\ker f^*D_0$ is a real line, and put $P=\operatorname{coker}
f^*D_0$. Differentiating the family of pullback operators on a section
in $K$ and projecting to $P$ gives the intrinsic derivative
\begin{equation}\label{p9-eq:wall-differential}
 W_f:T_f\mathcal H\oplus\RR_t\longrightarrow
                  \operatorname{Hom}_{\RR}(K,P).
\end{equation}
We use the complex orientation of $\mathcal H$ and the absolute
Cauchy--Riemann determinant orientation of $f^*D_0$. Indeed,
Riemann--Hurwitz gives $\dim_{\RR}P=2b+1$, and the determinant of the
space on the right is $(K^*)^{\otimes 2b}\otimes(\det f^*D_0)^*$.
Write $\eta_f$ for the sign of \eqref{p9-eq:wall-differential} when it is an
isomorphism. This sign uses the full pullback operator, including its
invariant summand \cite[(4.59) and (5.59)--(5.60)]{bai}.

\begin{theorem}\label{p9-thm:wall}
Fix an arithmetic genus bound $k_*$. Suppose the smooth connected
double covers of genus at most $k_*$ with nonzero pullback kernel at
$t=0$ form a finite set up to isomorphism. For each such $f$, assume
that the kernel is an anti-invariant real line, the map
$(f,\xi):Y\to N$ is an embedding for a generator $\xi$, and
\eqref{p9-eq:wall-differential} is an isomorphism. For sufficiently small
$\delta>0$ the pullback operators at the two endpoints are injective on every
degree-two stable map of genus at most $k_*$. In each
genus $k\le k_*$ the change of the absolute normal Euler number is
\begin{equation}\label{p9-eq:wall-jump}
 L_{2,k}(D_{\delta})-L_{2,k}(D_{-\delta})
 =\sum_{f:Y\to C}\eta_f\,[v^{2k-2}]
                \bigl(2\sin(v/2)\bigr)^{2g(Y)-2}.
\end{equation}
Each summand is the change supported on stable maps that factor through a degree-one
stable map to the source of the indicated cover. In particular a single wall with smooth
source of genus $H$ contributes the series
$\eta_f(2\sin(v/2))^{2H-2}$ coefficientwise.
\end{theorem}

\begin{proof}
We first treat a single cover $f$ of genus $H$. Choose a complex
connection on $N$. On a disk bundle about the zero section we realize
$D_t$ as the normal operator of an almost complex structure with
holomorphic projection and standard complex structure on each fiber.
For a holomorphic map $h$ to $C$, prescribe the equation of a lift
$(h,u)$ by
\begin{equation}\label{p9-eq:wall-cubic}
 h^*D_t u+h^*B(u,u,u)=0 ,
\end{equation}
where $B$ is a smooth real symmetric trilinear form on $N$ with
values in $\Omega_C^{0,1}\otimes N$. Here pullback of the form includes
composition with $dh$. In a connection splitting, the difference of
$D_t$ from the connection Dolbeault operator and the cubic both specify
vertical-valued $(0,1)$-forms in the horizontal variable. Such a form
specifies the off-diagonal block of the almost complex structure, giving
\eqref{p9-eq:wall-cubic} for every $h$ \cite[Lemma~5.9]{bai}. We multiply
the change of horizontal structure by a radial cutoff equal to one in
the disk under consideration and extend across $\PP(N\oplus\CC)$.
A unitary connection gives a closed form restricting to the
Fubini--Study form on the fibers. Adding a sufficiently large multiple
of a positive area form from $C$ tames the extended family, since the
horizontal terms are bounded on the compact bundle.

\smallskip
\noindent\emph{The smooth branch and its sign.}
The invariant part of $f^*D_0$ has index zero. To see this, deform
equivariantly to a pullback Dolbeault operator and use
$(f_*\mathcal O_Y)^+=\mathcal O_C$, together with exactness of
invariants under the deck group in characteristic zero. The hypothesis
on $K$ thus implies that $P$ is entirely anti-invariant. Take
equivariant complements to the kernel and image in a Sobolev completion
of \eqref{p9-eq:wall-cubic} over the Hurwitz chart. The implicit function
theorem solves the image equation for the kernel complement. In local
coordinates $y$ on the chart and with kernel coordinate $a\xi$, the
remaining equation is
\[
                    a F(y,t,a)=0,\qquad F(y,t,-a)=F(y,t,a).
\]
Indeed the zero section solves the equation at all $(y,t)$ and the
deck change negates both $a$ and the obstruction in $P$. The
differential of $F$ in $(y,t)$ at the origin equals $W_f$ evaluated
on $\xi$. Its coefficient of $a^2$ at $(y,t)=0$ is the image of
$f^*B(\xi,\xi,\xi)$ in $P$. The solution in the kernel complement
starts in degree three there, and its linear image projects to zero.

Let $p$ be a nonzero adjoint section annihilating the image of the
Hurwitz directions in $P$. On an unramified disk where $p$ and $\xi$
are nonzero on one sheet, prescribe an arbitrary compactly supported
form on that sheet as $f^*B(\xi,\xi,\xi)$. This uses a real covector
taking value one on $\xi$ and its third tensor power. Such a disk
exists by unique continuation. On the other sheet both the prescribed
form and the adjoint change sign, so their two integrals agree.
Choose the form to give a nonzero pairing. The implicit function
theorem for $F=0$ now gives even functions $y(a),t(a)$ with
\[
                     t(a)=c a^2+O(a^4),\qquad c\ne0 .
\]
The two signs of $a$ give the same curve by the deck action. Dividing
its vertical coordinate by $a$ gives smooth convergence to
$(f,\xi)$, so the curve for small $a>0$ is embedded and has trivial
stabilizer.

At a nonzero solution the derivative of $aF$ in $(y,a)$ has first
columns $aF_y$. Its last column modulo these is $-a t'(a)F_t$ by
differentiation along the branch. It is invertible, and comparison
with $W_f$ gives the sign $-\operatorname{sgn}(t)\eta_f$. In this
comparison the tangential deformation of the cover has the complex
Hurwitz orientation. Indeed its holomorphic deformation complex reduces
to the torsion cokernel of $df$, supported at the simple ramification
points. At the zero section the vertical and tangential determinants
split, and continuation gives the determinant of the embedded graph.
This is the absolute sign convention of
\cite[Theorem~4.12 and (4.30), (4.59)]{bai}. Thus the change in the
signed number of embedded curves with nonzero vertical coordinate is
$-\eta_f$.

\smallskip
\noindent\emph{Isolation through genus $k_*$.}
Suppose there were other simple curves, of degree at most two over $C$
and normalization genus at most $k_*$, in arbitrarily small closed
tubes at parameters tending to zero. Pass to a sequence with fixed
normalization genus. There is no nonconstant compact curve in an
affine fiber. In degree one the projection of the normalization is an
isomorphism, and invertibility of $D_0$ excludes a small nonzero graph.
Equivalently, a sequence of such graphs can be divided by its maximum
norm, and elliptic compactness over the fixed base would give a
nonzero element of $\ker D_0$.

In degree two write the maps from the smooth normalizations as
$(h_n,u_n)$ and put $\rho_n=\max |u_n|>0$. The graphs
$(h_n,u_n/\rho_n)$ solve \eqref{p9-eq:wall-cubic} with cubic coefficient
$\rho_n^2B$. On a disk bundle of radius greater than one these
structures extend to a smooth two-parameter family in $(t,\rho)$
including $(0,0)$. A common taming form as above has bounded area on
the graphs, since its cohomology class on the disk bundle is pulled
back from $C$ and the projection degree is fixed. Gromov compactness,
with a point marked where the norm is one, gives a stable limit in
the linear structure of $D_0$ whose vertical coordinate is nonzero.
Forget the extra mark and stabilize. If two normalized components
had degree one, invertibility of $D_0$ would force both sections to
vanish and all constant components would have the same zero vertical
coordinate at their attachments. Thus the unique effective
normalization is a smooth double cover with a nonzero kernel section.
By hypothesis its graph is an embedded curve $\Gamma=(f,\lambda\xi)(Y)$,
with $\lambda\ne0$.

The stable limit factors in degree one through $\Gamma$. Its effective
normalization maps isomorphically to $\Gamma$, so no self-node
identifies two points on that normalization. Each maximal connected
contracted subcurve attaches at a single point, since its image has a
single preimage on the effective normalization. For the resulting
map $w$ to $\Gamma$ and every real normal Cauchy--Riemann operator
$D_\Gamma$, restriction and pullback give inverse isomorphisms
\begin{equation}\label{p9-eq:wall-degree-one-ker}
                     \ker w^*D_\Gamma \simeq \ker D_\Gamma .
\end{equation}
A section on a contracted component is holomorphic with values in a
fixed vector space, hence constant, and the node conditions prescribe
its value on the whole contracted subcurve. The compactness theorem
for sequences of simple maps to a finite-dimensional family of almost
complex structures \cite[Theorem~3.1 and Remark~3.2(b)]{bai} says that
a limit over an embedded curve, unless the limiting map is an
isomorphism to that curve, has a normal kernel section outside the
pullbacks from the embedded curve. By \eqref{p9-eq:wall-degree-one-ker}
the limit is an isomorphism. We thus have smooth convergence modulo
parametrization. The original graphs eventually belong to the
Hurwitz chart and its implicit-function neighborhood, contradicting
their choice.

We also need injectivity at the endpoints on the entire fixed-target
space in genus $k\le k_*$. Apply the same compactness argument to
graphs of normalized kernel sections of the linear operators at
$t_n\to0$. The invariant part of any smooth double-cover kernel
descends to a base section. It descends on the unramified locus by
invariance and extends across the branch values by bounded
removability for the base equation. Since $D_t$ remains invertible,
the kernel sections are odd. Their nonzero graphs are simple because
the two sheets have different values wherever the section is nonzero.
After smooth convergence in the rescaling argument, the linear problem
in the Hurwitz chart reduces to $F(y,t,0)=0$, whose only solution
near the origin is $(0,0)$ by \eqref{p9-eq:wall-differential}.
Thus all smooth covers through the cutoff have injective pullback for
small $0<|t|$. For a nodal source, restriction to its effective
normalizations and the constant matching on its contracted components
give injectivity as well. At $t=0$ the same reasoning shows that every
stable map with kernel lifts in degree one to $Y$ via its graph in
$\Gamma$. Conversely every such lift carries the pullback of $\xi$.

\smallskip
\noindent\emph{Conservation of the local degree.}
The isolation just proved gives a closed tube of some radius $r>0$
and an interval about zero in which the only simple curves in that
tube through the cutoff are $C$ and the embedded curves found above.
Shrink the interval so the latter lie in the tube of radius $r/2$.
Every effective component of a degree-two stable map contained in the
closed tube factors through a simple normalization of degree at most
two and genus at most $k_*$. Connectedness therefore prevents any
such map from meeting the boundary. In the compactified target these
maps form a compact open and closed part of the parametrized stable-map
space. Compactness gives closedness, and support in the interior gives
openness. Its virtual degrees at the two ends agree by fixed-genus deformation invariance and the splitting of the virtual
class over open and closed substacks \cite[\S3.1, (3.4)--(3.5)]{diw}.

At an endpoint the maps supported on $C$ contribute $L_{2,k}(D_t)$
by the normal Euler comparison \cite[Theorem~1.2, (1.7)]{zinger}.
The maps on the additional embedded curve, when present, have degree
one over it. We compute their degree including its orientation.
Let $D'$ be the invertible normal operator on that embedded curve
of genus $H$. For any deformation of $D'$ on the curve, the pullback and base operators
admit compatible finite-dimensional stabilizations. The quotient of the
two resulting complexes has zero kernel on every degree-one stable
map $w$. Indeed
\eqref{p9-eq:wall-degree-one-ker} identifies the kernels, and restriction
to the effective normalization is a left inverse on cokernels to the
pullback of base forms. The quotient cokernel is therefore a vector
bundle throughout the deformation, by the stabilized construction in
\cite[Appendix~B, (B.8)--(B.10)]{bai}. Its determinant is the
quotient of the pullback and base determinants. At an invertible
endpoint the Euler class in the absolute orientation is the Euler
class of this bundle multiplied by $\epsilon(D')$.

The normal bundle has degree $2H-2$, since the ambient first Chern
class on the tube vanishes. Deform $D'$ to the Dolbeault operator of
two line bundles of degree $H-1$ off the theta divisor, using
$\mathcal O(-1)^{\oplus2}$ if $H=0$. There the cokernel has its
complex orientation and its virtual Euler degree is the ordinary
degree-one local invariant. Pandharipande's degree-one formula
\cite[Theorem~2]{pandhodge} gives, in our genus variable, the series
$(2\sin(v/2))^{2H-2}$. This also gives the degree of the reduced
cokernel at $D'$ by deformation. The change from the embedded curves
off $C$ is consequently $-\eta_f$ times this series through the
cutoff. Conservation gives \eqref{p9-eq:wall-jump} with the opposite
sign on $C$.

\smallskip
\noindent\emph{Simultaneous covers.}
For finitely many covers $f_i$ a single cubic $B$ can be chosen
outside the finite union of hyperplanes in the adjoint tests above.
The compactness argument then gives one branch in each indicated
Hurwitz chart and the sum in \eqref{p9-eq:wall-jump}. To identify the
individual supported changes, we separate their times. With generators
$\xi_i$ and adjoints $p_i$ as above, the conormal to the $i$th wall
on real zeroth-order variations of the base operator is
\[
                   \ell_i(A)=\langle p_i,f_i^*A\,\xi_i\rangle .
\]
For distinct covers these functionals are nonproportional. Work on a
common unramified disk where the sections and adjoints in question
are nonzero. The two sheet integrals for each $\ell_i$ agree. Testing
with arbitrary real-linear $A$ on this disk, proportionality would
therefore imply proportionality of the rank-one tensor densities,
and in particular real proportionality $\xi_i=c\xi_j$ on chosen
sheets. The two sections solve the same equation there, so
$\bar\partial c=0$ and $c$ is a nonzero constant. Off the union of
the branch values the two covers have distinct nonzero mod-two
characters, since a cover of the complement has a unique finite
normal completion. Their characters are independent and their fiber
product is connected. Unique continuation extends the equality over
this fiber product. A deck transformation flipping only the first
sheet then negates exactly one side, a contradiction.

Choose $\chi(t)$ equal to one near zero and supported away from the
endpoints. Write $t_i(s)$ for the crossing time of $f_i$ in the
family $D_t+s\chi(t)A$. Then
\[
              t_i'(0)=-\ell_i(A)/\ell_i(\dot D_0).
\]
The denominators
are nonzero by \eqref{p9-eq:wall-differential}. Pairwise
nonproportionality lets us choose $A$ so all derivatives are distinct.
The compactness argument for the linear graphs, now including $s$,
confines all kernels through the cutoff to the same charts. Thus the
perturbation separates the crossings and preserves their signs.

For completeness, let $M=\overline{\mathcal M}_{k}(C,2)$, let $I$
be the time interval, and let $Z$ be the closed noninjectivity locus
in $M\times I$. On its complement the absolutely oriented pullback
cokernel has an Euler class $e$. The image of $e$ under the
connecting morphism of the pair $(M\times I,(M\times I)\setminus Z)$
pairs with $[M]^{\rm vir}\times[I,\partial I]$ to give the final
Euler degree minus the initial one. At $s=0$ the kernel supports of
the distinct covers are disjoint compact images of
$\overline{\mathcal M}_k(Y_i,1)$, since the effective normalization
recovers $f_i$. Choose disjoint neighborhoods. For sufficiently small
$s$ all kernels remain in them and avoid their boundaries, by
properness and the local stabilized operators
\cite[Appendix~B, (B.8)]{bai}. The relative class for the whole
$(t,s)$-family splits over these neighborhoods by excision. Pairing
over the product virtual family shows that each summand has the same
degree at $s=0$ as at small positive $s$. There each is the
single-cover contribution already computed. All compactness choices
were made at the prescribed finite cutoff.
\end{proof}

On a fixed-target path with finitely many degree-at-most-two wall
times in all genera, the coefficients in Theorem~\ref{p9-thm:wall} may
be evaluated at common regular endpoints on either side of a single
time. On a regular parameter interval the pullback cokernels form an
oriented bundle on each compact fixed-genus map space, so their
Euler degrees are constant there. For generic fixed-target paths the
finiteness assertion follows from \cite[Definition~2.6,
Proposition~2.7 and Theorem~3.40]{dw}.
\subsection{A contracted elliptic tail at a ramification point}\label{p9-sec:tail}

Let $C$ be a smooth curve of genus $g\geq 2$, let $N\to C$ be a complex
bundle of rank two and degree $2g-2$, and let $D_t$ be a smooth path of
real Cauchy--Riemann operators on $N$. Fix a smooth connected cover
$f:Y\to C$ of degree $m$ and genus $H$, with $b$ simple ramification
points over distinct values. Write $G=\operatorname{Aut}_C(f)$ and
$a=|G|$. The fixed-target Hurwitz stack $\mathcal H$ is smooth of complex
dimension $b$ at $f$. Indeed, the quotient of $T_Y\to f^*T_C$ is a
skyscraper sheaf of length $b$, with one summand at each ramification point.

Suppose $K=\ker f^*D_0$ is a real line and the intrinsic derivative
\begin{equation}\label{p9-eq:tail-core}
 T_f\mathcal H\oplus\RR_t
 \longrightarrow \operatorname{Hom}_{\RR}(K,\operatorname{coker} f^*D_0)
\end{equation}
is an isomorphism. Its sign $\eta$ uses the complex orientation of
$\mathcal H$, increasing time and the absolute determinant orientation
transported from complex-linear operators. For $0\ne\xi\in K$ assume
$d\xi_r\ne 0$ at each ramification point $r$. In a frame pulled back from
the target, $d\xi_r:T_rY\to N_{f(r)}$ is independent of the connection
because $df_r=0$, and is complex-linear by the pullback equation.

\begin{theorem}\label{p9-thm:tail}
In source genus $k=H+1$, the stable maps obtained by attaching a
contracted one-pointed curve of arithmetic genus one to $f$ form an
isolated compact part of the locus with nonzero normal kernel at the
crossing. The local degree of the change of the absolute normal Euler class
supported on this family is
\begin{equation}\label{p9-eq:tail-total}
 -\frac{2\eta}{a}\left(\frac{m(g-1)}{12}+\frac{b}{24}\right)
 =-\frac{2\eta}{a}\frac{H-1}{12}.
\end{equation}
The first term comes from the component on which the tail persists.
The second is the sum over ramification points on the component on
which the attaching node smooths. The corresponding local degree at the smooth cover in genus $H$ is
$2\eta/a$. When the endpoint pullbacks are injective in the genus under
consideration, each local degree is a summand of the difference of the
endpoint Euler numbers.
\end{theorem}

\begin{proof}
For fixed $k$ let $X=\overline{\mathcal M}_k(C,m)$ be the stack of
connected stable maps, let $I$ be the parameter interval, and let
$U\subset X\times I$ be the open set where the pullback normal operator
is injective. Its oriented real cokernel bundle $Q_U$ has rank
$d=4(k-1-m(g-1))$, the real virtual dimension of $X$. The relative class
\begin{equation}\label{p9-eq:tail-half}
 u=\tfrac12\delta e(Q_U)\in H^{d+1}(X\times I,U;\QQ)
\end{equation}
defines the local degrees in the statement. Here $\delta$ is the
connecting homomorphism. With injective endpoints,
$2u$ paired with the image of $[X]^{\rm vir}\times[I,\partial I]$
is the difference of the endpoint Euler numbers by the boundary formula.
Excision gives the summand of this pairing supported near an isolated
compact part of the complement of $U$. This local pairing also makes
sense when other covers have kernel at the same time.

Near a map with kernel line $K$ and cokernel $Q$, stabilization gives
a finite-dimensional presentation of the normal operator
\cite[\S6.1 and Appendix~B, (B.8)]{bai}. In its matrix space the
stratum with kernel line has normal space $\operatorname{Hom}_{\RR}(K,Q)$.
Its determinant is $(K^*)^{\otimes d}\otimes(\det D)^*$, where $\det D$
denotes the Fredholm determinant at that map. Thus the absolute
orientation coorients this stratum. On the link of a normal disc,
choosing a generator of $K$ identifies the universal cokernel with
the tangent bundle of $S^d$. Its Euler degree is two with the boundary
orientation. For $d=0$ the two invertible matrices have opposite
determinant signs. Hence \eqref{p9-eq:tail-half} pulls back the positive
Thom class of this stratum on the corank-at-most-one neighborhood.
Two stabilizations admit a common sum, which adds an identity block
after a change of splitting. The normal space and its orientation are
preserved. In particular \eqref{p9-eq:tail-core} gives local degree
$\eta/a$ for $u$ at the smooth cover.

Write $(E,e)$ for the pointed tail, $p$ for the attaching point on $Y$,
and $\mathbb E$ for the Hodge line on $\overline{\mathcal M}_{1,1}$.
The compact family over $(f,0)$ is
\[
 T=[(Y\times\overline{\mathcal M}_{1,1})/G].
\]
We describe the fixed-target virtual class near $T$. The cover
deforms unobstructedly, and evaluation of the constant value on $E$
eliminates the matching condition at the node. The remaining obstruction
space is $T_{f(p)}C\otimes H^1(E,\mathcal O_E)$. Parameters consist of
the cover, the point $p$, the pointed tail, and the attaching parameter
$q\in T_pY\otimes T_eE$. The node-preserving maps exist for every value
of the first three parameters, so their scalar holomorphic obstruction
equation has the form $q F=0$. On $q=0$ its derivative in the smoothing
direction is $df_p\otimes\delta_e$, where
$\delta_e:T_eE\to H^1(E,\mathcal O_E)$ is the boundary map for
$\mathcal O_E(e)$. Indeed, substitution of the plumbing relation into
the linear term of the core map gives the simple principal part on $E$.
The dual of $\delta_e$ is evaluation of a dualizing differential at the
mark, an isomorphism also on the irreducible nodal elliptic curve.
Thus $F$ is a unit away from ramification. At a simple ramification
point it vanishes simply as $p$ varies. Locally there are two smooth branches
$B=\{q=0\}$ and $M=\{F=0\}$ near such a point. The fixed-target
virtual class is the divisor of the scalar obstruction section and hence
the sum of their fundamental cycles, each with multiplicity one. At a
nodal tail the modulus of the pointed elliptic curve includes the
independent smoothing parameter of its self-node.

On $B$, normalizing at the attaching node identifies the normal kernel
with the core kernel. Constant tail sections evaluate onto the matching
fiber, so the cokernel over $T$ is the sum of the core cokernel and the
complex two-plane $V=f^*N\otimes\mathbb E^\vee$. Taking the Thom class
along \eqref{p9-eq:tail-core} leaves the Euler class of
$\operatorname{Hom}_{\RR}(K,V)$. On the cover
$Y\times\overline{\mathcal M}_{1,1}$, $K$ is constant and either
generator induces the complex orientation on this excess bundle. The
degree on $B$ is therefore
\begin{equation}\label{p9-eq:tail-moving}
 \frac{\eta}{a}\int_{Y\times\overline{\mathcal M}_{1,1}}
 c_2(f^*N\otimes\mathbb E^\vee)
 =-\frac{\eta}{a}\frac{2m(g-1)}{24}.
\end{equation}
Indeed $c_2=-c_1(f^*N)c_1(\mathbb E)$ and
$\int_{\overline{\mathcal M}_{1,1}}c_1(\mathbb E)=1/24$.
The latter degree follows from the Weierstrass presentation as the
weighted projective stack of weights $4,6$, with Hodge line
$\mathcal O(1)$. Transfer from the
cover accounts for the action of $G$, including its action on $K$.

It remains to restrict $u$ to $M$ at a ramification point $r$. Work on
a uniformizing chart before division by $G$. Choose core and target
coordinates so that $w=w_0+z^2$. On the pointed elliptic family choose
a meromorphic function $P$ with its sole pole of order two at $e$.
This choice extends locally across the nodal fiber by base change for
$\mathcal O_E(2e)$, which has two sections, with only the constants
in $H^0(E,\mathcal O_E(e))$. Take a coordinate $y$ with $P=y^{-2}$
near $e$. The plumbing
\begin{equation}\label{p9-eq:tail-plumbing}
 zy=q,\qquad
 w=w_0+z^2\quad\text{on the core},\qquad
 w=w_0+q^2P\quad\text{on the tail}
\end{equation}
gives maps agreeing on the overlap. Varying the marked ramification
cover and the elliptic modulus gives every tangent to $M\cap B$, and $q$
smooths the attaching node to first order. The inverse function theorem
on the smooth branch $M$ proves that these are local parameters there.

Use cokernel representatives supported on the core outside the attaching
disc and on the tail away from both nodes. The former span by unique
continuation for the adjoint equation, the latter by duality with the
Hodge line. Stabilization by these spaces is surjective with kernel a
real line near the central map. Normalize its section $s$ by a real
evaluation on the thick part of the core, giving $\xi$ at the center.
The stabilized kernels are continuous across the nodal fibers
\cite[Appendix~B]{bai}. In the plumbing norms, a continuous local
section is uniformly bounded and converges on compact subsets away from
the nodes \cite[\S\S4.3--4.4, Proposition~4.15]{siebert}. For the
pulled-back zeroth-order term the multiplication norm is bounded by the
sup norm of the section times the weighted $L^p$ norm of the map
derivative. A map holomorphic on the total space has bounded derivative
on each half-collar in its larger coordinate, so this latter norm tends
to zero on shrinking collars. On the remaining compact sets the
coefficients vary smoothly. The elliptic
obstruction coefficient $c_E$ vanishes identically along $q=0$, where
the stabilized core section extends by its constant value to the tail.

We compute $c_E$ to first order. In a holomorphic frame for a reference
Dolbeault operator on the target disc, the equation on the truncated
elliptic side $|y|\geq |q|/\rho$ reads
\[
 \bar\partial s+A_t(w)(s)\,\overline{dw}=\theta_E(c_E).
\]
Here $A_t(w)$ is real-linear on the fibers. Normalize the forcing
$\theta_E$ by pairing with a regular differential
$\omega=(a_0+O(y))\,dy$ near $e$. With the trace convention for
$\delta_e$, the boundary integral in Stokes' formula becomes
\[
 q a_0\,\frac{1}{2\pi i}\oint_{|z|=\rho}s(z)\frac{dz}{z^2}
       + O_\rho(|q|^2).
\]
For fixed $\rho$ the contour coefficient tends to that of $\xi$ as all
parameters approach the chosen point of $T$. Since $\bar\partial\xi(r)=0$,
it is $\partial_z\xi(0)+O(\rho)$ in the limit. For fixed small $c>0$,
the zeroth-order term on $|q|/\rho\leq |y|\leq c$ is bounded by
\[
 C|q|^2\int_{|q|/\rho}^c u^{-2}\,du\leq C\rho |q|.
\]
Outside this disc $dw=q^2dP$ and the integral of $|dP|\,|\omega|$ is
uniformly bounded, also on a smoothing $uv=\tau$ of the self-node.
Indeed $\omega$ is a bounded multiple of $du/u$, and the derivative
of the regular function $P(u,\tau/u)$ is bounded on $|u|\geq
\sqrt{|\tau|}$. The other half is covered by $v$. On the nodal fiber the two
residues in Stokes' formula cancel by matching of $s$ on the normalization.

Dividing by $q$, then sending $\rho$ to zero, gives the elliptic
coefficient to order $o(|q|)$ as all parameters tend to the wall. Its
linear term in the attaching direction is the intrinsic homomorphism
\begin{equation}\label{p9-eq:tail-jet}
 \begin{aligned}
 T_rY\otimes T_eE &\longrightarrow
 \operatorname{Hom}_{\RR}\bigl(K,N_{f(r)}\otimes H^1(E,\mathcal O_E)\bigr),\\
 v\otimes v'&\longmapsto
 \bigl[\xi\longmapsto d\xi_r(v)\otimes\delta_e(v')\bigr].
 \end{aligned}
\end{equation}
This limit holds along arbitrary simultaneous variations of the cover,
time and elliptic modulus. Compactness of $\overline{\mathcal M}_{1,1}$
gives a uniform estimate in finitely many charts. The map
\eqref{p9-eq:tail-jet} embeds the complex smoothing line since $d\xi_r\ne0$.
In a tubular neighborhood extend the splitting of the obstruction space
from the divisor. The straight-line homotopy from the elliptic
coefficient to the extended linear term stays nonzero for $q\ne0$
after shrinking the neighborhood. With this coefficient fixed, retract
the core coefficient to its restriction at $q=0$, where its derivative
is \eqref{p9-eq:tail-core}. These homotopies require only continuity of the
core coefficient in $q$. They apply over the whole ramification-marked
family: the leading map transforms as the tensor in \eqref{p9-eq:tail-jet},
and splittings can be chosen equivariantly using a metric and common
stabilizations.

After these two Thom isomorphisms the excess bundle on the deck cover
is the complex line
\[
 \mathbb E^\vee\otimes
 \bigl(N_{f(r)}/d\xi_r(T_rY)\bigr).
\]
Its degree is $-1/24$. The smoothing line has its complex orientation.
Interchanging it with the oriented real normal space
to the wall of smooth covers preserves the sign.
Thus the degree of $u$ on $M$ for the orbit of $r$ is
$-\eta/(24|\operatorname{Aut}_C(f,r)|)$. Summation over the orbits
gives $-\eta b/(24a)$. The transversality on $B$ and the nonvanishing
for $q\ne0$ on $M$ isolate $T$ in the kernel locus. Adding the two
virtual-cycle contributions, doubling by \eqref{p9-eq:tail-half}, and
using $2H-2=2m(g-1)+b$ proves \eqref{p9-eq:tail-total}.
\end{proof}

For a double cover of a genus-two target with two branch points,
$H=4$ and $a=2$. The smooth-cover jump is $\eta$. In genus five
the moving attachment contributes $-\eta/6$, and smoothing at either
ramification point contributes $-\eta/24$, giving $-\eta/4$ in all.
In general the ratio $-(H-1)/12$ of the two successive jumps is the
quadratic coefficient of $(\sin(v/2)/(v/2))^{2H-2}$, as in the
degree-one formula \cite[Theorem~2]{pandhodge}.

\section{Estimates and virtual classes at a joined double cover}\label{app:lowerjoins}
\subsection{The local equation for Theorem~\ref{p10-thm:nodes}}\label{app:double-equation}
Fix the joining set and the branch frames of \S\ref{p10-sec:nodes}.
\subsubsection{The deformation space and the matching equations}

At one selected point the nodal map and its deformation are given by
\begin{equation}\label{p10-eq:node-family}
          zy=q,\qquad w=w_0+(z+p)^2+y .
\end{equation}
When $q=0$, the coordinate on the normalized double is $z+p$ and
its attaching mark has coordinate $p$. Translating $w_0$ to zero and
eliminating $y$ gives the free algebra on $1,z,z^2$ with relation
\begin{equation}\label{p10-eq:node-cubic}
                    z(z+p)^2-wz+q=0 .
\end{equation}
To see versality, use the relative cotangent complex of this
hypersurface over the target \cite[Tag~08SL]{stacks}. Its dual on the
central fiber is multiplication by $3z^2-w$ between two copies of
the local ring. This element is a nonzerodivisor on the two branches.
Its cokernel is $\CC\{z\}/(z^3)$, with basis represented by the
variations of $q,w_0,p$, respectively. An unselected branch contributes
one complex parameter. The first deformation sheaf has finite support
and the complex has no other cohomology, so the fixed-target chart is
smooth of complex dimension $b+2j$ and carries its ordinary fundamental
class. Distinct branch disks supply independent pairs $(p_i,q_i)$.
The stabilizer is the single involution of $Y$ over $C$. In
\eqref{p10-eq:node-family} it acts by $(z,p,q)\mapsto(-z,-p,-q)$ at every
join and fixes $s$.

Put $P_I=\operatorname{coker} h_I^*D$. A kernel section restricts to
zero on the identity and to a multiple of $\xi$ on $Y$. Every such pair
matches because $\xi$ vanishes at the marked points. The normalization
sequence for the operator thus identifies $\ker h_I^*D=K$ and gives
\begin{equation}\label{p10-eq:node-matching}
       0\longrightarrow\bigoplus_{i\in I}N_{x_i}
       \longrightarrow P_I\longrightarrow P_f\longrightarrow0 .
\end{equation}
The smoothing calculation takes place in the individual matching fibers
of this sequence. We give it at one point, suppressing $i$.

\subsubsection{The pullback operator on the fixed target}

For a map $h$ in the finite chart, let $F=h_*\mathcal O$ on $C$. It is
a rank-three bundle by \eqref{p10-eq:node-cubic}, the ordinary quadratic
algebra at the remaining branch points, and the sheet decomposition off
the branch disks. Translate the moving branch centers on smaller disks
and extend the translations by diffeomorphisms supported in larger
disks. Identify the family of bundles smoothly, using the polynomial
frames on the smaller disks. On the intervening collars the sheets are
unramified and the identifications are smooth in the parameters.
Fix $2<a<4$. In these trivializations consider the operator
\[
  \mathcal A_h:W^{1,a}(C,N\otimes F)\longrightarrow
       L^a(C,\Omega_C^{0,1}\otimes N\otimes F)
\]
which acts by $D_t$ after evaluation on each regular sheet. The
complex-linear part is the tensor product Dolbeault operator. In the
polynomial frame the antilinear part applies $B$ to the coefficients
and replaces $1,\bar z,\bar z^2$ by their quadratic interpolants at
the three roots. This determines a measurable matrix $M_{p,q}$ almost
everywhere. At an ordinary simple branch the sole quadratic-sheet
factor in the interpolation of $\bar z$ is $\bar z/z$, which is bounded.

Let $d$ be the diameter of the three roots and $m$ their maximum
modulus. In a second divided difference order the two most distant
roots first and last. The first divided differences of conjugation have
modulus one and those of $\bar z^2$ have modulus at most $2m$. Hence
the respective second differences are bounded by $2/d$ and $4m/d$.
Newton interpolation bounds every entry of $M_{p,q}$ by $C(1+d^{-1})$.
Subtracting the mean of the roots depresses \eqref{p10-eq:node-cubic} to a
monic cubic with linear coefficient $-(w+p^2/3)$, so
$|w+p^2/3|\le C d^2$. It follows that the matrices have uniformly bounded
$L^{a'}$ norm for $a<a'<4$. Away from the limiting branch values they
converge pointwise, hence they converge in $L^a$. Since
$W^{1,a}\hookrightarrow C^0$ is compact, the antilinear part is a
continuous family of compact perturbations of the differential part.

We identify its finite-dimensional reduction with that of the actual
pullback operator. A polynomial section in the kernel evaluates to a
bounded solution on the normalization off the branch points and node
preimages. Removability extends it there, and the common root gives
equal values at the two sides of a node. Evaluation is injective by
interpolation on the open unramified locus. This argument also applies
to solutions with smooth forcing supported in that locus. At the
central map the generator on the two local components is represented
by $z u_1(w)$, where $u_1(w)=\xi(z)/z$ on the double. Oddness and
$\bar\partial\xi(0)=0$ give $\xi(z)=v z+O(|z|^3)$ with derivative
error $O(|z|^2)$. Thus $u_1(0)=v$ and its first weak derivatives in
$w=z^2$ are bounded. The same division on each ordinary branch disk
shows that this is a global coefficient solution. The two central
kernels agree.

The real indices agree by deformation to the complex-linear operators
and the finite projection formula, including on the nodal fiber.
Consequently the central cokernels have equal dimension. Smooth forms supported on unramified disks span both cokernels: a
dual annihilator vanishes on these disks and hence off the nodes
by unique continuation. In the normalized matching presentation,
testing arbitrary endpoint values then kills the matching multipliers.
Inside a finite common span surjecting onto both, a generic tuple of
minimal length projects to a basis of each. Use its span $E$ as a
stabilization. Both stabilized operators are surjective near the
central map and their kernels have rank one, using the continuity of
the pullback reduction on nodal curves
\cite[Appendix~B, (B.6)--(B.9)]{bai}. Evaluation identifies the two
lines and commutes with projection to $E$. Normalize a generator by a
fixed real functional supported on an unramified disk and taking value
one on $\xi$. The two obstruction maps to $E$ then coincide. Enlarging
$E$ by further such forms gives the same identification of their
cokernel classes, since evaluation again identifies the stabilized
kernels of equal dimension.

\subsubsection{The logarithmic term}

Set $\ell=\max(|p|,|q|^{1/3})$ and write $Q=q\log(1/|q|)$ for $q\ne0$,
with $Q=0$ at zero. The only entry of $M_{p,q}$ whose $q$-variation
has nonintegrable first derivative is the coefficient of $z^2$ in the
interpolant of $\bar z$. We compute its variation against a H\"older
continuous test using Euclidean area $dA(w)$.

In the disk $|w|\le A\ell^2$, for fixed large $A$, the roots have
size $O(\ell)$ by rescaling \eqref{p10-eq:node-cubic}. The bound by a
constant times $1+|w+p^2/3|^{-1/2}$ gives an integral $O(\ell^3)$
there for the varying and central matrices. On the complementary
annulus choose $r^2=w$ and set $z=rZ$. The roots in $Z$ are analytic
near $0,1,-1$ in $p/r$ and $q/r^3$, which are uniformly small. Their
interpolation matrix is real analytic in these parameters, with bounded
second derivatives. The coefficient in question is multiplied by
$\bar r/r^2$ on returning to $z$. All other nonconstant entries scale
by factors with nonnegative power of $|r|$.

At $p=q=0$ the interpolant of $\bar z$ is $(\bar r/r)z$. To first
order in $q$, the root increments at $0,r,-r$ are $q/w,-q/(2w),-q/(2w)$.
Subtract the old interpolant applied to those increments from their
conjugates. The three differences are $a_0,-a_0/2,-a_0/2$ with
$a_0=\bar q/\bar w-(\bar r/r)q/w$. Their quadratic coefficient is
\begin{equation}\label{p10-eq:node-modes}
            -\frac{3a_0}{2w}=-\frac32\left(
                 \frac{\bar q}{|w|^2}-\frac{\bar r}{r}\frac{q}{w^2}
                              \right).
\end{equation}
The second term has angular mode $-3$ in $w$ and hence zero mean. In
either term subtracting the value of the test at zero costs $O(|q|)$
in the integral by H\"older continuity. Against the constant test one,
\eqref{p10-eq:node-modes} integrates to
$-6\pi\bar q\log(1/\ell)+O(|q|)$.

For completeness the Taylor remainders can be estimated before angular
integration. Writing $R=|r|$, the largest is bounded by
$C R^{-1}(|p|/R+|q|/R^3)^2$. Integration against $O(R^3\,dR)$ gives
\[
       O\bigl(|p|^2+|pq|/\ell+|q|^2/\ell^3\bigr)
                        =O(|p|^2+|q|).
\]
The $q$-linear terms of all other entries have integral $O(|q|)$ in
absolute value. The $p$-linear terms are bounded by $C|p|/|w|$ on
the disk, and deleting its inner part costs $O(|p|\ell^2)$. For
$\ell=|q|^{1/3}$ the logarithmic integral above is $-2\pi\bar Q+O(|q|)$.
For $|q|<|p|^3$ both $|q|\log(1/\ell)$ and $|Q|$ are uniformly
$o(|p|)$. Since $\ell^3\le |p|^3+|q|$, this proves, also on the axes,
that $M_{p,q}-M_{0,0}$ has $L^1$ norm $O(|p|+|Q|)$ and that its
pairing with any fixed H\"older test is its integral of its first-order
$p$-variation plus the point mass
\begin{equation}\label{p10-eq:node-mass}
              -2\pi\bar Q\,\delta_x
         \quad\text{in the }(\bar z,z^2)\text{ coefficient},
\end{equation}
up to $o(|p|+|Q|)$. Translating the branch center before this
calculation leaves the matrix on the smaller disk unchanged in form.
Variation of its smooth factors and of the data on the collars adds
their ordinary linear terms in $s,p$ and an error
$o(|s|+|p|+|Q|)$ in these pairings.

To pass to the obstruction map, pair with a central adjoint solution.
At the central parameter the conjugation matrix is bounded, sending
$\bar z$ to $(\bar r/r)z$ and $\bar z^2$ to $(\bar w/w)z^2$.
The adjoint solution, initially in $L^{a/(a-1)}$, thus lies in
$W^{1,k}$ for every finite $k$ by elliptic estimates and Sobolev
iteration. In particular it is H\"older continuous. The normalized
stabilized generator converges in $W^{1,a}$ to the central one.
In pairing the operator variation with the difference of these
generators the multiplier has $L^1$ norm $O(|s|+|p|+|Q|)$ and the
generator difference is $o(1)$ in $C^0$. The smooth variation of the
principal part on the collars gives the same little-oh estimate from
$W^{1,a}$ convergence. The central operator applied to the difference
pairs to zero. A varying stabilizing form contributes only its smooth
variation times its coefficient, which tends to zero. We may therefore
apply \eqref{p10-eq:node-mass} to the central generator in the actual
reduced equation. Its coefficient at the node is $B(x)v$.

\subsubsection{Identification of the residue and the local degree}

The distribution $a_1 z^2\delta_x$ represents in
\eqref{p10-eq:node-matching} a common nonzero real multiple of the matching
class of $a_1\in N_x$. To verify this, choose a local real-linear frame
$S(w)$ of $D$-homogeneous sections, with $S(0)=1$. Such a frame follows
by solving the local equation for the constant real basis vectors with
the Cauchy integral right inverse on a sufficiently small disk, and
renormalizing their values at zero. On the punctured disk the section
$z^2 S(w)a_1/w$ takes the value $S(w)a_1$ on the double and zero on the
identity. Multiply by a cutoff equal to one near zero. Its
distributional image is the ordinary cutoff forcing plus
$\pi a_1 z^2\delta_x$, since $\partial_{\bar w}(1/w)=\pi\delta_x$.
The ordinary forcing represents the matching class up to the convention
for the evaluation difference and projects to zero on $P_f$. The total
image pairs to zero with every adjoint by integration by parts, using
$1/w\in L^k$ for $k<2$ and the adjoint regularity just proved. This
identifies the residue in its own matching fiber, with the same real
scalar for every $a_1$.

On the nodal stratum $q_i=0$ the normalized component maps are
independent of the $p_i$. Their obstruction in $P_f$ has derivative
$W_f$ in $s$. The $i$th evaluation difference has derivative
$\pm v_i p_i$ in its moving mark, and may have a linear term in $s$.
These derivatives follow by stabilizing the normalized component
operators off the marks and varying the evaluation row in
\eqref{p10-eq:node-matching}. The comparison with the fixed-target
stabilization preserves the derivatives at the zero, since a continuous
change of trivialization only multiplies their linear part by its
central value.



\begin{proof}[Proof of Theorem~\ref{p10-thm:nodes}]
Choose the stabilizing forms away from all selected disks. Each disk
contributes \eqref{p10-eq:node-mass} with its own cutoff, so none of its
residue projects to another matching fiber. With $Q_i=q_i\log(1/|q_i|)$
extended by zero, the combined error is
\[
             o\left(|s|+\sum_{i\in I}(|p_i|+|Q_i|)\right)
\]
because there are finitely many disjoint disks. This holds on the
partially smoothed strata as well. After choosing a generator of $K$
and filtering each $N_{x_i}$ by $\CC v_i$, the linear part in these
coordinates is block triangular. Its diagonal blocks are $W_f$
evaluated on the generator, the complex-linear maps
$p_i\mapsto\pm v_i p_i$, and the maps
$Q_i\mapsto c_i\bar Q_i[B(x_i)v_i]$ to $N_{x_i}/\CC v_i$ for
nonzero real $c_i$. By \eqref{p10-eq:node-jet} the last maps are
antilinear isomorphisms of complex lines. The radial coordinate
change $q_i\mapsto Q_i$ preserves orientation near zero. The
little-oh estimate gives isolation and a homotopy on every sufficiently
small link to the invertible linear part.

Normalization gives the Fredholm determinant isomorphism
\[
  \det h_I^*D\simeq \det f^*D\otimes\det D\otimes
       \bigotimes_{i\in I}(\det_{\RR}N_{x_i})^* .
\]
It preserves the orientations obtained by deformation to the complex
normalization sequence. The evaluation fibers have complex orientation
and the one invertible identity component contributes $\epsilon$.
Replacing the generator of $K$ negates an even number $4j$ of real
matching coordinates. Each attaching block has positive sign and each
smoothing quotient negative sign. Thus the degree with the absolute
orientation upstairs is $(-1)^j\epsilon\eta_f$. On its link the
cokernel is the quotient by the image of the normalized obstruction
vector, hence pulls back the tangent bundle of the sphere of dimension
$2b+4j$. Its Euler number gives twice the local degree. Dividing by
the order two of the single deck group gives the asserted change.
Distinct subsets give distinct maps since the branch values in the
fixed target are distinct.

It remains to verify genericity. On a fixed double cover mark a simple
ramification point $r$ over $x$ and consider, on odd sections, the map
\[
          (D,\xi)\longmapsto(f^*D\xi,\partial_z\xi(r),B_D(x)) .
\]
It is a submersion at a nonzero solution. Indeed an annihilator of its
differential has an adjoint distribution solving the homogeneous adjoint
equation off $r$. On an unramified disk disjoint from $x$ where $\xi$ is
nonzero on one sheet, real zeroth-order variations prescribe arbitrary
forcing on that sheet and its negative on the other. Oddness and unique
continuation therefore force the distribution to be supported at $r$.
For any prescribed first holomorphic jet, choose a section whose equation
vanishes to higher order than the distribution there. This is obtained
by prescribing the pure holomorphic Taylor coefficients and recursively
solving the equation for coefficients involving $\bar z$, then cutting
off and antisymmetrizing. These variations kill the first-jet multiplier.
The adjoint distribution now solves the homogeneous equation across $r$,
so is smooth and zero. Free variation of $B_D(x)$ kills the remaining
multiplier. In Sobolev spaces of sufficiently high order the linearized
range is closed, since the section operator is Fredholm and the two
evaluations are finite dimensional.

At $v\ne0$, variation of $B$ surjects onto $N_x/\CC v$. Thus failure
of \eqref{p10-eq:node-jet} has real codimension two. The marked projectivized
odd incidence has index $-1$ over the operator space. Indeed, the odd
operator has real index $-2b$, since the invariant part has index zero
by Dolbeault descent. Varying the Hurwitz parameters adds $2b$ real
dimensions and projectivizing removes one. The branch mark has finite
choices. Adding a path parameter leaves the bad incidence of negative
expected dimension. Parametric transversality, followed by countable
intersection over covering genera and charts, gives the assertion
relative to rigid endpoints.
\end{proof}
\subsection{A ramification bound and compactness}\label{p10-sec:compactness}
Let $N$ for the moment be a Hermitian bundle of any rank over $C$.
Fix a unitary connection $\nabla$ and a metric on $C$ in its
conformal class. The graph equation for a lift $u$ of a finite
holomorphic map $f:Y\to C$ is
\[
  \bar\partial_{f^*\nabla}u + f^*P(\cdot,u)=0.
\]
Here $P$ is a smooth fiber-dependent $(0,1)$-form on a bounded
fiberwise convex tube. For a positive orthonormal pair $X,JX$ on $C$,
let $M$ bound the fiber Lipschitz constant of $P(\cdot,\cdot)(X)$,
uniformly in $X$. Put
\[
  \kappa=\bigl(\lambda_{\max}(iF_\nabla(X,JX))+2M^2\bigr)\,dA .
\]
We use the Hermitian product conjugate-linear in the first variable.
For a compact family of linear normal equations one may use a common
$M$, independent of the height of the graph.

\begin{lemma}\label{p10-lem:curvature}
Suppose $(f,u)$ is generically one-to-one on its smooth
normalization, with distinct component images. Let $m=\deg f$,
let $R$ be its total ramification, and suppose $n$ disjoint pairs
of distinct normalization points are identified as nodes at which
the graph matches. Then
\[
 R+2n\le \frac{m(m-1)}{2\pi}\int_C\kappa .
\]
In particular a connected graph with these nodes and no contracted
components has arithmetic genus at most
$1+m(g-1)+\frac{m(m-1)}{4\pi}\int_C\kappa$.
\end{lemma}
\begin{proof}
Normalize the closure of the off-diagonal pairs in $Y\times_C Y$ over
regular values of $f$. The resulting smooth, possibly disconnected
curve $Z$ has projections $p_1,p_2$ and a map $b=fp_i$ of degree
$m(m-1)$ to $C$. The section $s=p_1^*u-p_2^*u$ is nonzero somewhere on
every component. Subtracting the two equations gives
$|(\bar\partial_{b^*\nabla}s)(T)|\le M|s|\,|db(T)|$ for each unit
tangent $T$. Integrating the vertical derivative of $P$ along
the segment between the two heights writes the difference as a
smooth real-linear zeroth-order operator applied to $s$. Its zeros
are isolated and have finite order by the similarity principle.
In a smooth unitary frame the leading Taylor term of order $a$ is
$z^a v$ with $v\ne0$, since the principal equation on that term is
$\bar\partial=0$.

Away from the zeros, set $e=s/|s|$ and
$\alpha=\operatorname{Im}\langle e,b^*\nabla e\rangle$.
Project the derivative perpendicular to $\CC e$ and denote the
result by $\nabla^\perp e$. Differentiation in a positive unit
frame $T,jT$ gives
\begin{align*}
 d\alpha(T,jT)={}&\operatorname{Im}
   \langle e,(b^*F_\nabla)(T,jT)e\rangle \\
  &+2\bigl(|(\nabla^\perp e)^{1,0}(T)|^2-
                            |(\nabla^\perp e)^{0,1}(T)|^2\bigr).
\end{align*}
The components parallel to $e$ have purely imaginary coefficients and
drop out of the exterior product. The bound on $\bar\partial s$ gives
$|(\nabla^\perp e)^{0,1}(T)|\le M|db(T)|$, whence
$d\alpha\ge-b^*\kappa$. On a positively oriented small circle about
a zero of order $a$, the integral of $\alpha$ tends to $2\pi a$.
Stokes' theorem on the complement of the disks bounds the sum of the
orders by $(2\pi)^{-1}\int_Z b^*\kappa$.
At a point of index $k$ there are $k-1$ off-diagonal branches over
the repeated point, one for each nontrivial $k$th root of unity.
Each forces a zero. Each identified pair forces a zero in each
order of the pair, at points distinct from the ramification pairs.
This proves the inequality. Riemann--Hurwitz on the components
and the arithmetic genus formula for a nodal curve give the last claim.
\end{proof}

\begin{proposition}\label{p10-prop:compact}
Over a generic compact path as in Theorem~\ref{p10-thm:chambers}, the
projectivized incidence of smooth connected cubic covers with kernel
is compact and finite, with the genera of all covers included.
\end{proposition}
\begin{proof}
Proposition~\ref{p10-prop:incidence} gives discreteness. For a sequence of
kernel lines choose generators of fixed small maximum norm and mark
points attaining that norm. Their embedded graphs lie in a compact
subtube. The family of graph almost complex structures agrees there
with a smooth family on a compact completion, tamed by a common form
as in the proof of Proposition~\ref{p10-prop:smoothjump}. Projection
determines the homology class and hence gives a uniform bound on
area. Lemma~\ref{p10-lem:curvature} bounds the genus, so a subsequence
has fixed genus. Stable-map compactness gives a limit whose marked
point has nonzero height. Forget that mark and stabilize.

A nonconstant component projects in positive degree, since the
portion of a single fiber in the tube lies in a complex vector space.
Invertibility on $C$ excludes a degree-one component of nonzero
height. If a component has degree three and nonzero height, its
parametrization is simple, since the alternative factors through
a nonzero graph of degree one. This cubic normalization is embedded
by Proposition~\ref{p10-prop:incidence}. Otherwise the positive components
are a nonzero double graph and one copy of the zero section,
both with multiplicity one. The double generator is odd. It can vanish
only at ramification points, since a zero over a regular value would
identify the two sheets of its embedded graph. Its derivative at a
ramification point gives a vertical tangent, distinct from the tangent
to the zero section.

For a maximal stable contracted subcurve in a limit of smooth
holomorphic curves, the attaching branches include a critical branch,
two branches with coincident tangent lines, or at least three
branches \cite[Theorem~1.1]{ekholm}. This theorem applies to the
varying structure by extending the parameter interval to a complex
disk, smoothly extending the family in its real coordinate, and
using the fixed block structure $J_t\oplus i$ on the product.
The curves are constant in the parameter direction. Each candidate
limit above has at most two branches through any image point, and
they are immersed with distinct tangent lines when both occur.
Consequently there is no contracted subcurve. Neither embedded
normalization admits a self-node. A single cubic limit therefore
converges in the smooth-cover chart and retains its kernel line.
In the two-component case connectedness requires direct joins at a
nonempty subset of the ramification points. Theorem~\ref{p10-thm:nodes}
isolates the projective kernel there even in the full smoothing
chart. Such a point cannot be approached by smooth cubic solutions.
Rigidity excludes endpoint limits. This proves compactness.
\end{proof}

For degree two a normalized nonzero graph has only the embedded
degree-two possibility in the preceding proof. The same argument
gives finiteness, also proved in \cite[Proposition~3.73]{dw}.
The contracted components excluded in Proposition~\ref{p10-prop:compact}
do enter the Euler class on the full stable-map space.
\subsection{Contracted components at a reducible cubic cover}\label{p10-sec:descendants}

Let $f:Y\to C$ be the double cover at a crossing as in
Theorem~\ref{p10-thm:nodes}, and take the crossing parameter to be $0$.
Write $r_1,\ldots,r_b$ for its ramification points, $x_i=f(r_i)$,
and $\xi$ for a nonzero odd kernel section. For a nonempty set
$I\subset\{1,\ldots,b\}$, let $h_I:X_I\to C$ be the finite map
obtained from $Y\amalg C$ by identifying $r_i$ with $x_i$ for
$i\in I$. Put
\[
 H=g(Y)+g(C)+|I|-1,\qquad d_I=(-1)^{|I|}\epsilon\eta_f .
\]
In arithmetic genus $k\ge H$, let $G_{I,k}$ consist of the stable
maps obtained from $h_I$ by attaching connected contracted curves
at smooth points and by inserting connected contracted curves between
the branches of its nodes. We include all stable configurations in
these operations. In particular an inserted curve may have rational
components. We regard $G_{I,k}$ as a locus at $t=0$ in the product
of the ordinary fixed-target moduli stack with the parameter interval.

\begin{theorem}\label{p10-thm:descendants}
Suppose the hypotheses of Theorem~\ref{p10-thm:nodes} hold with its
projected jet condition at every ramification point. Then $G_{I,k}$
is compact and isolated in the locus of stable maps with nonzero
pullback kernel.
Its contribution to the change of the absolute normal Euler number in
arithmetic genus $k$ is
\begin{equation}\label{p10-eq:desc-contribution}
 d_I\,[v^{2k-2}]\bigl(2\sin(v/2)\bigr)^{2H-2}.
\end{equation}
Thus the contribution of the specified set of joining nodes to the
connected series is $d_I x^{H-1}$, where $x=(2\sin(v/2))^2$.
\end{theorem}

We displace the nodal kernel to a signed collection of embedded
graphs on smooth domains. Near each such graph the Euler integral is
a constrained virtual count in the total space. Cutting its parameter
interval at one regular value leaves the ordinary degree-one series.


\subsubsection{The compact locus and the choice of sheets}
A presentation by proper gluing maps makes explicit the allowable
specializations at an attachment. Choose nonnegative integers
$e_Y,e_C,e_i$ with $e_Y+e_C+\sum_{i\in I}e_i=k-H$.
Take the two evaluation fibers
\[
 \overline{\mathcal M}_{g(Y)+e_Y,|I|}(Y,1)_{(r_i)_{i\in I}},
 \qquad
 \overline{\mathcal M}_{g(C)+e_C,|I|}(C,1)_{(x_i)_{i\in I}} .
\]
At $i$, join the corresponding marks directly if $e_i=0$ and
insert a member of $\overline{\mathcal M}_{e_i,2}$ if $e_i>0$.
Compose the maps to $Y$ with $f$. The union of the images over these
finitely many choices is $G_{I,k}$. Indeed a degree-one stable map
to a smooth curve has one nonconstant normalization, mapping
isomorphically, and every connected contracted subcurve maximal for
inclusion meets this normalization just once. Its markings may lie
on contracted components. This accounts also for tails whose attaching
points come together or specialize to a joining point. Conversely,
cutting out all maximal contracted subcurves that join the two
positive components, and marking the ends of each direct joining node,
recovers such gluing data.

In the lift to $N$, the positive components are the graph of $\xi$
over $Y$ and the zero section over $C$. They are embedded and meet
precisely over the $x_i$, with distinct complex tangent lines.
Every contracted connected subcurve has a single image point. Thus
its external positive branches number at most two, and two can occur
only on opposite components. The direct nodes and the contracted
subcurves with two ends recover $I$ intrinsically. The proper gluing
images for distinct $I$ are therefore disjoint. Restriction to the
positive components identifies the kernel throughout $G_{I,k}$ with
$\RR\xi$, since the base operator is invertible and on contracted
subcurves the sections are constants with prescribed attaching values.

Choose concentric closed disks $U\subset\operatorname{int}U'$ of
regular values disjoint from the $x_i$. Over $U'$ the three heights
$0,\xi,-\xi$ are separated. Even if a contracted curve lies over a
point of $U$, a nearby stable map has three pieces over $U$, each of
total degree one. In fact at a central map the three preimages of $U'$
with their contracted components have disjoint neighborhoods. The
collapse maps of converging domains keep the preimage of $U$ in
these neighborhoods, and evaluation at a regular value away from the
limiting nodal images gives degree one in each. Order the pieces with
the identity sheet first. These two orderings over the central locus
extend to a finite étale cover of a neighborhood, by taking all sheet
orderings locally and continuing the two selected ones. Its stack
degree is two. The deck involution of $f$ exchanges the orderings
and fixes each $r_i$. This removes the double-cover symmetry at the
finite map and retains the isotropy of the contracted curves. We work
on this cover until the last step of the proof.

We recall the local class computing the Euler change. In a finite
stabilization $S:E\to F$ of the normal family, projectivize $E$
over $\RR$ and restrict $S$ to its tautological line. The pushforward
of the resulting relative Euler class is supported on the kernel
locus. Near $G_{I,k}$ we can exclude kernels of dimension two by
upper semicontinuity. In the space of matrices the normal space at a
matrix of kernel dimension one is $\operatorname{Hom}(K,P)$ for
its kernel $K$ and cokernel $P$. The pushforward gives the Thom
generator there, since the projective vertical differential cancels
the image of the matrix. On the normal sphere the oriented cokernel
is its tangent bundle after choosing a generator of $K$. The sphere
has even dimension, and its Euler integral is two
\cite[Lemma~6.10]{bai}. Thus the connecting class of the cokernel
Euler class off the kernel locus equals twice the projective class
by the Thom isomorphism in the matrix space and pullback. Its pairing
with the fixed-target virtual class times the oriented interval is
the terminal Euler number minus the initial one. Both this identity
and its restriction by excision are preserved by common stabilization.

\subsubsection{Displacement on the finite-map charts}
On one labelled double sheet let $\ell_a(u)$ be the integral of a
smooth real covector applied to $u$, against the pullback by $a$ of a
cutoff area form supported in $\operatorname{int}U$. Choose the test
so that $\ell_a(\xi)\ne0$. It gives the same normalization of the
kernel through every attachment on that sheet, since pullback area
vanishes on contracted components. Let $j_a(e)$ denote a linear
combination of triples of smooth $N$-valued $(0,1)$-forms supported
in $\operatorname{int}U$, pulled back separately on the corresponding
labelled pieces. For the variable stable map $a$ put $L_{a,t}=a^*D_t$
and consider the finite-rank perturbation
\begin{equation}\label{p10-eq:desc-perturb}
 T^e_{a,t}u=L_{a,t}u-\ell_a(u)j_a(e).
\end{equation}
For continuity of the reduction at contracted curves we explain how
to use the nodal stabilization of \cite[Appendix~B]{bai}. Add a
scalar $z$ with $L_{a,t}u=zj_a(e)$ and $z=\ell_a(u)$. Fix a central
stable map. Enclose the distinct nodal values meeting the supports of
the forcings in small disjoint disks inside $U'$. On each disk and
each labelled piece choose a smooth primitive $b_i$ solving
$D_t b_i=j_i(e)$ on the base. These primitives can be chosen
smoothly in $t$ and linearly in $e$: on a small disk the Cauchy
right inverse for $\bar\partial$ followed by the zeroth-order term
has norm less than one after rescaling. With cutoffs equal to one on
smaller disks about the nodal values, pull the primitives back on
their respective labelled pieces and extend by zero. Write $b_a$ for
their sum. The substitution $u=u'+zb_a$ changes the forcing to
$j_a(e)-L_{a,t}b_a$ and the scalar equation to
$z=\ell_a(u')+z\ell_a(b_a)$. The new forcing vanishes on
neighborhoods of all the nodes. This uses the pullback identity on
maps to $C$, including the constant components.

On shrinking the chart, the remaining forms vary smoothly on compact
thick subsets away from the necks in the sense of convergence of
stable maps. In a domain trivialization they converge there in every
smooth norm; pullback by $a$ involves only its first derivative.
One may freeze these forms at the center in the nodal stabilization:
the error tends to zero in operator norm and is absorbed by the
uniformly bounded stabilized inverse
\cite[Propositions~B.7.9 and~B.10.4]{pardonatlas}. At the center a minimal cokernel
supplement $B$ gives the invertible augmented operator
$(v,w)\mapsto(L_{a,t}v-Bw,\ell_a(v))$ since $\ell_a$ is nonzero on
the kernel line. The area test on shrinking base disks has size at
most a constant times their area for bounded sections, by the
degree bound. This proves continuity also of the scalar row. On
overlaps the substitutions are translations by the difference of
the two sums of primitives times $z$, with the induced scalar row,
and hence give the same projective class. Near a central normalized line, its stabilized sections
converge on thick regions and are uniformly bounded on the domains.
Thus the area test is continuous, and the normalization $\ell_a(u)=1$
is available near the central projective kernel. On this slice the
equation becomes
\begin{equation}\label{p10-eq:desc-affine}
                  L_{a,t}u=j_a(e),\qquad \ell_a(u)=1 .
\end{equation}
The first equality is a local graph equation. In the three disjoint
fiber neighborhoods over the support of the forms, multiply the
corresponding forcing by a smooth cutoff equal to one near that
sheet. The resulting affine term in the horizontal-to-vertical
block defines an almost complex structure on $N$ with holomorphic
projection and standard fibers \cite[Lemma~5.9]{bai}. On the
normalized graphs it gives the indicated equation even when a
contracted curve lies over the support. Cutting off outside a larger
disk bundle extends the structure over $\PP(N\oplus\CC)$.
The fiberwise Fubini--Study form extended by a connection, plus a
sufficiently large pullback of an area form on $C$, tames the family.

Fix $k$. Use the same sheet data on the finite-map charts for all
$h_J$ of genus at most $k$, $J\ne\varnothing$. By
Theorem~\ref{p10-thm:nodes} each has an isolating closed neighborhood in
its own genus. Independent forms on the labelled sheets span the
cokernel at every one of their centers. To see this, present the
operator on the normalization with one evaluation-difference row
at each join. A dual annihilator of the sheet forms vanishes on
a disk of each positive component. The adjoint equation and unique
continuation give vanishing off the joins. The equation-space
multipliers lie in the dual of $L^p$ for $p>2$, so none is supported
at an isolated point. Sections with independently prescribed endpoint
values then kill the remaining evaluation multipliers. Hence one
finite-dimensional space of triples suffices for all these centers.

In a minimal stabilization the normalized obstruction section takes
values in this space of representatives, and \eqref{p10-eq:desc-affine}
translates the section by the constant parameter $e$. If more forms
are used, the derivative in the forcing parameters is still
surjective. On each stratum of fixed nodal and branching type the
section is smooth by normalization with the matching equations.
Persistence of a node has positive complex codimension by the local
smoothing parameters of Theorem~\ref{p10-thm:nodes}. On smooth sources
collision of branch values also has positive complex codimension by
the Hurwitz dimension formula. The unrestricted normalized section
has expected real dimension zero at fixed $e$, with the path
parameter included. Regular values therefore avoid these strata and
give isolated zeros on smooth simply branched coverings. Compactness
of the isolating charts makes their number finite.

We may arrange also that their graphs are embedded. On a smooth
source the vertical derivative at a simple ramification point takes
values in a complex two-plane. Equality of heights at two distinct
points over a regular value imposes four real conditions with two
real parameters for the value. At a branch value the marks are
discrete. The corresponding marked evaluation and first-derivative
conditions are transverse under the sheet forcings. Indeed an
annihilating distribution vanishes on the open test sheets away
from the marks by the independent variations. The adjoint equation
there eliminates any multiple of the normalization density by testing
where that density is nonzero. Unique continuation leaves a
distribution supported at the marks. In a local trivialization one
can prescribe independent values and complex first derivatives at
the marks with all jets of the linearized equation vanishing up to
any specified finite order. This follows by solving recursively for
the derivatives involving $\bar z$ from the Cauchy--Riemann symbol.
Testing on these sections removes the point multipliers. One first
obtains immersion, which excludes coincidences in a neighborhood
of the diagonal. On the compact remaining marked pairs, finitely many
additional forms give the required transversality near the finitely
many regular solutions. Arbitrarily small regular choices in the
smooth space of triples result, each in a finite-dimensional
perturbation. Write $\sigma_\Gamma$ for the oriented projective
degree of a resulting smooth solution in the chart of $h_I$.
Conservation of the degree of the normalized obstruction section in
that chart gives
\begin{equation}\label{p10-eq:desc-sum}
                        \sum_\Gamma\sigma_\Gamma=d_I
\end{equation}
by the oriented local calculation of Theorem~\ref{p10-thm:nodes}.

\subsubsection{Removing the contracted components}
Consider a sequence of projective zeros tending to $G_{I,k}$ with
$e\to0$, allowing also $e=0$. Use the normalized graph equation and
take stable limits of the nonconstant normalizations separately,
first retaining their nodal flags. Their limiting nonconstant images
together are the graph of $\xi$ and the zero section, each with
multiplicity one. This follows by gluing the marked limits and
forgetting the flags. Each nonconstant normalization retains positive
projection degree, since every compact curve in an affine fiber is
constant. Thus there are either one or two such normalizations.

In the unmarked stable limit of each normalization there are no
contracted components. In fact ghost bubble censorship requires a
vanishing derivative at a positive attaching branch, equal complex
tangent lines at a pair of attaching branches, or at least three
positive attaching branches at one point
\cite[Theorem~1.1]{ekholm}. The two central embeddings exclude
all three possibilities. To use the fixed-structure theorem in a
smooth finite-dimensional family, extend the parameters locally to
an open subset of $\CC^m$ and put the block-diagonal almost complex
structure on the product. The graphs on constant parameter slices
are holomorphic for this single structure, and the attaching
directions are unchanged. If a sequence uses successively larger
spaces of smooth forcings, choose a subsequence whose first $n$
seminorms differ from the limit by at most $2^{-n^2}$ at step $n$.
Disjoint smooth bumps of width $2^{-n-3}$ centered at $2^{-n}$
interpolate it to a smooth arc flat at zero. The same product argument
then applies.

If there are two positive normalizations, their individual unmarked
limits are smooth embedded copies of $Y$ and $C$. An identification
on either limit would give a self-intersection of its embedding.
The normalizations thus converge smoothly and remain embeddings.
Near a point of intersection of the two limiting images, project
the difference of their two local parametrizations onto the real
four-plane spanned by the two tangent lines. Its derivative in the
two source variables is invertible. On smaller neighborhoods the
inverse function theorem gives at most one intersection for
each nearby pair, and there are none away from these neighborhoods.
Consequently every maximal contracted subcurve has one end on
a positive normalization or one on each. Delete the one-ended
subcurves, replace the two-ended ones by nodes, and retain the
direct nodes. Matching and the thick integral show that the finite
nodal map so obtained satisfies the same normalized equation.
After subselection its joining sites approach a fixed subset $J$.
The flags for different connections have distinct limits on either
component. Marked stabilization and properness of the one- and
two-pointed contracted factors recover a limit in $G_{J,k}$, even
when other flags coalesce with a joining flag. Hence $J=I$ and the
finite zeros converge in the chart of $h_I$. At $e=0$ they are
$h_I$ itself, by Theorem~\ref{p10-thm:nodes}. The regular displacements
just chosen avoid all finite nodal zeros there.

With one positive normalization, its unmarked limit is a connected
finite union of the two central embeddings with direct nodes at some
nonempty subset $J$. Its normalized equation is unchanged on deleting
the other components, which are constant. For $e=0$, smooth finite
zeros approaching $h_J$ contradict Theorem~\ref{p10-thm:nodes}. For the
chosen regular displacements they are smooth embedded graphs from
the chart of $h_J$. Embeddedness excludes self-nodes and forces each
contracted connected subcurve maximal for inclusion to attach once.
Mark these attachments before passing to the limit. Forgetting the
marks on the positive limit recovers $X_J$. Restoring them can insert
a rational tree over a smooth point or between the two branches of an
existing node. Proper limits of the one-pointed tails attach to these
trees. In particular flags approaching two distinct smooth points
above an unjoined graph intersection stay on distinct fibers of the
stabilization. The limit belongs to $G_{J,k}$, so again $J=I$.

The zero-displacement argument proves isolation. Properness in fixed
genus and closedness of the kernel locus give a relatively compact
neighborhood with kernel-free boundary, which persists through
sufficiently small perturbations. The two cases show that within it
the displaced zeros are precisely the stable degree-one expansions
of the smooth graphs from the chart of $h_I$. For the reverse inclusion,
every such expansion solves \eqref{p10-eq:desc-affine} by constant
extension. If expansions escaped the neighborhood as $e\to0$, their
marked stable limits would lie in $G_{I,k}$ by the same gluing
argument. This also separates the compact families over the individual
displaced graphs.

\subsubsection{The virtual class of the normalized graph equation}
Fix one displaced graph $\Gamma$, at $t_0$, and write $a_0$ for its
projection. The homogeneous operator $L_{a_0,t_0}$ is injective.
Indeed $a_0$ has smooth source in the isolating finite-map chart of
the original nodal zero. Its pullback to every degree-one expansion
is still injective, by restriction to the positive normalization
and matching constants on the other components. On this compact
family the cokernel has real rank
\[
                   r=4\bigl(k-1-3(g(C)-1)\bigr),
\]
the real virtual dimension of $\overline{\mathcal M}_k(C,3)$.
We compare the projective Euler pairing for \eqref{p10-eq:desc-perturb}
with the virtual degree of ordinary maps to the almost complex
total space subject to the scalar equation in \eqref{p10-eq:desc-affine}.

Take an equivariant chart for maps to $C$, with stabilizing markings,
incidence slices, horizontal perturbation forms, and obstruction
section $s_C$. Use the same markings for maps to the total space,
taking inverse images of the target slices. Lift the horizontal forms
using a complex-linear splitting of projection for the graph
structure. On a horizontally thickened map the vertical equation
is still affine in $u$. Indeed in connection coordinates the
off-diagonal term of the almost complex structure is affine in height
where a sheet forcing acts, and linear elsewhere near the graph.
The horizontal lift can be chosen affine in height there by projecting
a connection lift to its complex-linear part. Taking the
$(0,1)$-part on the source then gives an affine vertical equation
with the usual Dolbeault symbol, restricting to $L_{a,t}u-j_a(e)$ on
$s_C=0$.

Choose vertical forms $B w$, $w\in W$, inducing an isomorphism from
$W$ to the cokernel of $L_{a,t}$ at the chart center. Forms on the
smooth interiors of every component, including contracted ones, are
available. They span by density off the nodes in the equation spaces
for sections with matching values. Choose equivariant representatives
by averaging and extend them over the stabilized universal domain.
A contracted component can require a vertical stabilizer even though
the forcing $j_a(e)$ vanishes there. In these conventions the
vertically stabilized equation on $s_C=0$ is
\begin{equation}\label{p10-eq:desc-stabilized}
                      L_{a,t}u-Bw=j_a(e).
\end{equation}
The derivative $(L_{a,t},-B)$ in $(u,w)$ is an isomorphism at the
center. In the simultaneous thickening the first row is the regular
stabilized horizontal equation, independent of the height. The
second has this isomorphism as its vertical derivative. Projection
therefore identifies the ambient thickening locally with the
horizontal one, with $(u,w)$ uniquely determined on it. This holds
over a nodal fiber by the regular perturbed-map chart and inverse
function theorem relative to the common base of stabilized domains
\cite[Theorem~3.5 and Lemma~5.10]{swam}. The ambient obstruction
section on the common chart is $(s_C,w)$.

On the horizontal zero set introduce $z=\ell_a(v)$ for a homogeneous
section $v$. Eliminating the variables in
$L_{a,t}v-Bw'=zj_a(e)$ by \eqref{p10-eq:desc-stabilized} gives
$v=zu$, $w'=zw$. At a projective zero $w'=0$ implies $z\ne0$ by
injectivity of $L_{a,t}$. Thus in the chart $z=1$ the residual
equations are exactly
\begin{equation}\label{p10-eq:desc-rows}
                           w=0,\qquad \ell_a(u)-1=0 .
\end{equation}
Capping the horizontal virtual class with the localized projective
Euler class amounts on this common thickening to multiplying the
Thom classes of $s_C$ and of \eqref{p10-eq:desc-rows}. Their product is
the class of the constrained ambient equation.

For compatibility, on overlaps take the union of stabilizing marks
and the sum of the horizontal and vertical perturbation spaces,
transporting earlier forms by the forgetful maps of marked domains.
The lifted slices and projection commute with these forgetful maps.
An enlarged vertical space adds free variables in
\eqref{p10-eq:desc-stabilized} paired with identical additional
obstruction coordinates. Switching the minimal space $W$ changes
cokernel coordinates and adds a multiple of $w$ to the scalar row.
A change of horizontal lift adds terms vanishing on $s_C=0$.
The resulting stabilizations and triangular row operations preserve
the localized Thom product, since throughout the homotopies the
joint zeros stay in the compact interior. The same sums refine triple
intersections, equivariantly throughout, as in the ordinary
Gromov--Witten implicit atlas \cite[\S9.2, Theorem~9.2.10]{pardonatlas}.
The determinant orientation under projection is the product of the
horizontal and vertical orientations, by the triangular
linearization and homotopy to its complex symbol. Moving the scalar
row past the horizontal virtual dimension or the vertical rank $r$
has even parity. Its orientation convention can thus be fixed at
the genus-$H$ crossing and used in every arithmetic genus.

\subsubsection{Restriction to one parameter value}
At $\Gamma$ let $A$ be the fixed-time normal linearization of the
ordinary ambient equation. Its real index is zero since the ambient
first Chern number on the graph is zero. The derivative with time and
the integral constraint, in unparametrized normal coordinates, is
\begin{equation}\label{p10-eq:desc-augmented}
 (\zeta,\tau)\longmapsto
          (A\zeta+b_0\tau,\lambda\zeta).
\end{equation}
Here $b_0$ is the parameter derivative and $\lambda$ the differential
of the time-independent integral test. This operator is invertible.
Indeed the horizontal equation at a smooth finite cover is regular
with variable domain, as its deformation complex reduces to the
torsion cokernel of the derivative of the cover. Eliminating that
equation and the vertical complement leaves the transverse normalized
equation that defined $\sigma_\Gamma$. Hence $\dim_{\RR}\ker A\le1$.

If $A$ is singular, choose an unramified subdisk of a labelled sheet
on which generators of its kernel and adjoint cokernel are nonzero.
Such a subdisk exists by unique continuation. If $u_\Gamma(w)$
is its height, add to the affine equation a small term
\[
                     \delta B_1(w)\bigl(u-u_\Gamma(w)\bigr)
\]
with support in this subdisk and a fiber cutoff equal to one near
the sheet. On an unramified graph the vertical tangent projects
isomorphically onto the normal quotient. We can therefore prescribe
the zeroth-order normal variation so that its pairing between the
two generators is nonzero. The one-dimensional Fredholm reduction
makes $A$ invertible for sufficiently small $\delta\ne0$, leaves
$\Gamma$ a solution, and preserves invertibility of
\eqref{p10-eq:desc-augmented}. Near the graphs this is the homogeneous
deformation \eqref{p10-eq:desc-perturb} with $L$ and $j$ replaced by
$L+\delta B_1$ and $j+\delta B_1u_\Gamma$ on the chosen sheet.
By compactness and continuous stabilization its size can be chosen
to preserve vertical injectivity near the expansions and projective
nonvanishing on the boundary of their isolated neighborhood, both in
genus $k$ and in genus $H$. Thus it preserves the localized degrees.

To see that no part of this degree lies away from the retained graph
after the perturbation, consider a sequence of other constrained
solutions as $\delta\to0$ in the original isolated neighborhood.
Its compact limit is an expansion of $\Gamma$. In a tube retracting
to that multiplicity-one curve only one nonconstant normalization
can occur, since a positive component projects with degree a
positive multiple of three. The unmarked normalization then converges
smoothly by ghost bubble censorship at the embedded limit. The
implicit function theorem for the augmented operator
\eqref{p10-eq:desc-augmented}, with $\delta$ as parameter, gives the
unique nearby constrained curve, namely the retained $\Gamma$
at $t_0$. All remaining components attach to it at one point each.

At the regular endpoint the nearby ordinary embedded graphs form
a smooth family $\Gamma_t$. In a small tube retracting to
$\Gamma$, the projection of every positive class represented by a
holomorphic component has degree a positive multiple of three.
A map of total projection degree three has one nonconstant component,
of multiplicity one. At fixed genus, a sequence of its unmarked
normalizations approaching $\Gamma$ has smooth limit. Indeed the
limiting positive map is the embedding with multiplicity one, and
ghost bubble censorship excludes any contracted component of the
unmarked limit. Smooth convergence places the normalization in the
regular chart of $\Gamma_t$.
Compare also the embedded-cycle compactness and openness results
of \cite[Propositions~2.21 and~2.29]{diw}. Embeddedness then implies
that every other component belongs to a one-ended contracted subcurve.
Consequently, after shrinking the tube and parameter interval, the
entire fixed-genus ordinary locus near the compact family consists of
degree-one stable maps to $\Gamma_t$.

On this locus the integral constraint equals the function of time
$\gamma(t)=\ell(\Gamma_t)-1$, since every attaching subcurve has
zero pullback area. Invertibility of \eqref{p10-eq:desc-augmented} gives
$\gamma'(t_0)\ne0$ by applying its scalar row to the tangent of
the implicit family. Interpolating between $\ell_a(u)-1$ and
$\gamma(t)$ on the common thickening preserves the localized product,
as the two functions agree wherever $s_C=w=0$. The real Thom class
of $\gamma$ cuts the parametrized ambient virtual class down to the
ordinary fixed-time class with sign $\operatorname{sgn}\gamma'(t_0)$.

Write $D_\Gamma$ for the regular normal operator on $\Gamma$.
Its normal bundle has degree $2H-2$, since
$c_1(TC)+c_1(N)=0$. Every degree-one pullback of $D_\Gamma$ is
injective by restriction and matching. Zinger's comparison therefore
expresses the local degree as the normal Euler pairing on the ordinary
degree-one stable-map class of $\Gamma$ \cite[Theorem~1.2]{zinger}.
To record its orientation, deform $D_\Gamma$ on the fixed bundle.
For each degree-one map the pullback and base operators have the same
kernel throughout the deformation. On cokernels restriction to the
positive normalization is a left inverse to pullback of forms.
Hence the quotient of their two stabilized Fredholm complexes is
injective with vector-bundle cokernel throughout
\cite[Appendix~B]{bai}. Its determinant is the quotient determinant.
At the regular endpoint its Euler class gives the absolute pairing
multiplied by the base determinant sign $\epsilon(D_\Gamma)$.
At the Dolbeault operator of a sum of two non-effective line bundles
of degree $H-1$ the orientation is complex. Pandharipande's
degree-one Hodge integral \cite[Theorem~2]{pandhodge} gives its
generating series $\bigl(2\sin(v/2)\bigr)^{2H-2}$. The oriented
fixed-time series at the regular endpoint is therefore
$\epsilon(D_\Gamma)x^{H-1}$.

\begin{proof}[Completion of the proof of Theorem~\ref{p10-thm:descendants}]
The product of the time-slice sign and the fixed-time determinant
sign is $\sigma_\Gamma$, as follows by evaluating their product in
genus $H$ at the transverse normalized graph. The scalar orientation
and quotient determinant comparison above give the same product in
genus $k$. Thus this graph contributes
$\sigma_\Gamma[v^{2k-2}]x^{H-1}$ to the projective pairing upstairs.
Summing by \eqref{p10-eq:desc-sum}, multiplying by two for the full
Euler change, and dividing by the stack degree two of the sheet
labelling gives \eqref{p10-eq:desc-contribution}. All isolating
neighborhoods and perturbation sizes were chosen at fixed genus.
In particular the formula includes tails and bridges of every
arithmetic genus allowed by $k$.
\end{proof}

\section{The two estimates in the smoothing at a cubic branch}\label{app:joined-estimates}
We use the parameters $h,(q_i)$, the nodal derivative $J$, and its
remaining real adjoint $\Lambda$ from \S\ref{p14-sec:joins}.
The different orders of the matching and adjoint equations determine
the degree at a join. Here we prove \eqref{p14-eq:fullestimate}
and \eqref{p14-eq:rowestimate} for the stabilized equation.

Choose thick stabilizing forms representing $P_I$ away from the
moving disks and fix a real normalization of the height, equal to one
on $s$. Solving the stabilized equation and taking its forcing
coordinate gives a continuous column $F(h,(q_i))$ with values in
$P_I$. The model \eqref{p14-eq:cubicpolynomial} permits computation in
a fixed vector bundle over $C$. Values at the three roots identify
its sections with sections of the pullback away from the branches.
The conjugates $\bar z,\bar z^2$ are replaced by their polynomial
interpolants of degree at most two. At the central fiber this imposes
the ordinary matching of the two normalizations at the node. For
$2<p_0<4$ these matrices are continuous as $L^{p_0}$ multipliers
in the parameters, by the divided-difference bound
$O(1+|w+p^2/3|^{-1/2})$. Bounded sections on the source extend over
the branches. Interpolation identifies the central stabilized kernel
by even and odd decomposition on the double sheet and matching
to the third sheet. At the smooth branch factors the same argument
applies, and the indices agree by the complex-linear finite projection formula.
Forms on thick unramified disks span both cokernels by adjoint
continuation and the matching sequence. Choose a minimal common
supplement. Both stabilized problems have a line of solutions nearby
by continuity of the pullback stabilization \cite[Appendix~B]{bai}.
Evaluation identifies the lines and their forcing coordinates.
The details of this comparison for the algebra \eqref{p14-eq:cubicpolynomial} are in \S\ref{p10-sec:nodes}.

Set $Q_i=q_i\log(1/|q_i|)$, with $Q_i=0$ when $q_i=0$.
First, in the entire supplement one has
\begin{equation}\label{p14-eq:fullestimate}
 F=Jh+2\iota\big((\overline Q_i B_i v_i)_{i\in I}\big)
                     +o(|h|+\max_i|Q_i|).
\end{equation}
Here is the coefficient of the logarithm. On a single disk put
$\ell=\max(|p|,|q|^{1/3})$. Outside $|w|\le A\ell^2$, take
$r^2=w$ and expand the roots in $p/r,q/r^3$. The quadratic
coefficient of the interpolant of $\bar z$ has linear $q$-term
\[
  -\frac32\left(\frac{\bar q}{|w|^2}
                   -\frac{\bar r}{r}\frac{q}{w^2}\right).
\]
Only its first summand has nonzero angular mean. For a fixed
H\"older test the integrated logarithmic term is
$-2\pi\bar Q\delta_0$ in the $z^2$ component. Indeed the inner
integral is $O(\ell^3)$, the outer quadratic error is
$O(|p|^2+|pq|/\ell+|q|^2/\ell^3)$, and if $|q|<|p|^3$ the
logarithmic terms are already $o(|p|)$. All other entries have
integrable $q$-derivatives. The full matrix difference has $L^1$
norm $O(|p|+|Q|)$. Continuous convergence of the normalized
stabilized solution in $W^{1,p_0}\subset C^0$ therefore permits
using its central value in this pairing to first order. The section
$\chi z^2 S(w)a/w$, with $S(0)=1$ a local real frame of solutions,
has ordinary forcing of class $\iota(a)$ and point forcing
$\pi a z^2\delta_0$. Their sum pairs to zero with a central
adjoint. Apply this to the generator with derivative $v_i$.
The disks are disjoint, which proves \eqref{p14-eq:fullestimate}.

The missing row requires a smaller remainder. Put
$R_0=\max_i|q_i|$. On each cone $|h|\le C R_0\log(1/R_0)$ on
which all $|q_i|$ are comparable and positive, one has
\begin{equation}\label{p14-eq:rowestimate}
                       \Lambda F=\sum_{i\in I} H_i(q_i)+o(R_0).
\end{equation}
We give the estimate on one disk. The half-trace of the forcing
$a_0+a_1z+a_2z^2$ is $a_0+w a_2$. Write $\varepsilon=|q|^{1/3}$,
$R=|w|^{1/2}$. On $|w|\le A\varepsilon^2$ the roots are
$O(\varepsilon)$ and their diameter is bounded below by
$c\varepsilon$: their mean is $O(p)=o(\varepsilon)$ and their
product is $-q$. The constant and quadratic coefficients of the
interpolant of $\bar z$ are $O(\varepsilon)$ and $O(1/\varepsilon)$.
The weighted inner integral is therefore $O(\varepsilon^5)$.
The central odd height is $z u_1(w)$ with $u_1(0)=v_i$ and
$u_1(w)-v_i=O(|w|)$, by oddness and the pulled-back equation.
On the exterior the weighted quadratic error from the three roots is
$O(R(|p|/R+|q|/R^3)^2)$. Integration against $O(R^3\,dR)$ gives
$O(|p|^2+|pq|+|q|^2/\varepsilon)=o(|q|)$. The linear terms are bounded
by $C(|p|+|q|/|w|)$; restoring their inner integrals costs $o(|q|)$.
Smooth center and time changes of the coefficients and collars have
Taylor remainder \(O(|h|^2)\). Their products with the varying
interpolation matrix have integral \(O(|h|(|h|+R_0\log(1/R_0)))\)
by the full matrix bound. Both are \(o(R_0)\) on the comparable cone. Other simple branch centers may be fixed by
translation on their disjoint source disks. Without translation the
angular multiplier on such a disk has first and second derivative
$O(|w|^{-1})$ and $O(|w|^{-2})$, giving integrated Taylor error
$O(|h|^2(1+|\log|h||))$.

For the varying generator, the worst unweighted matrix change in
$L^{p_0}$ on a joining disk is bounded by
\[
 C\big(\varepsilon^{4/p_0-1}
       +|p|\varepsilon^{4/p_0-2}\big).
\]
Use the root diameter on the inner disk and the outer bound
$CR^{-1}(|p|/R+|q|/R^3)$ elsewhere. Thus the bounded normalized
stabilized inverse gives $O(R_0^\delta)$ convergence of the generator
in $W^{1,p_0}$ on the comparable cone, for some $\delta>0$.
The product with the $L^1$ multiplier variation
$O(|h|+R_0\log(1/R_0))$ is $o(R_0)$ in the adjoint pairing; the
principal-part variation on the thick collars satisfies the same
bound. This justifies differentiation of the frozen height in
the half-trace pairing. At fixed regular target point the two
sextic roots move by $-q/(2w)$ to first order. Their averaged
variation has pole $-qv_i/(2w)$ at $y_i$ on the quotient cubic.
Away from that point the differentiated forcing is the cubic
operator applied to this sheet variation, including its collar cutoff.
Green's formula against the adjoint gives minus $\pi$ times
the real pairing with the residue, since
$\partial_{\bar w}(1/w)=\pi\delta_0$. The bounded remainder has
zero limiting boundary pairing. The $h$-derivative pairs to zero
by $\Lambda J=0$. Adding the integrals proves \eqref{p14-eq:rowestimate}.

\section{Comparison with the double-cover count on a closed threefold}\label{sec:normalization}
We turn to a second compact consequence. The chamber weight for smooth double covers is normalized at a complex-linear operator. In the elementary cluster expansion the reference is the balanced local curve. Computing the two references gives a correction in twice the primitive class.
Let $v$ be the genus parameter. Write $s=2\sin(v/2)$ and $x=s^2$. For a partition $\lambda$ and a box
$u$ of its diagram, let $\ell(u)$ denote its hook length. The connected
elementary series of a curve of genus $g$ is
\[
 G_g(q,v)=\log\!\left(\sum_{\lambda}q^{|\lambda|}
            \prod_{u\in\lambda}(2\sin(\ell(u)v/2))^{2g-2}\right).
\]
The empty diagram contributes one before taking the logarithm.
Ionel--Parker's $e^{\rm IP}_{B,h}$ is the coefficient of
$G_h(Q^B,v)$ in the connected Gromov--Witten series
\cite[Proposition~3.3 and (8.1)]{ip}. In particular
$[q]G_g=x^{g-1}$. The normal weights in the elementary series are
opposite. The class variable $Q$ records integral homology classes.

For a generic compatible almost complex structure, denote the
determinant sign of the normal operator $D_C$ of an embedded curve $C$
by $\epsilon_C$. The Bai--Swaminathan count in genus $h$ is
\[
 \operatorname{Gr}^{\rm BS}_{2A,h}
 =\sum_{\substack{[C]=2A\\g(C)=h}}\epsilon_C
       +\sum_{[C]=A}\epsilon_C w_{2,h}(D_C).
\]
Here $w_{2,h}$ is the oriented intersection number with the wall of
smooth genus-$h$ double covers carrying an anti-invariant kernel
\cite[(1.5), Definitions~5.8 and~5.17]{bai}. Its value at a
double-cover-rigid complex-linear reference is zero. We use the
coorientation and path convention of that definition throughout.
We extend $w_{2,h}$ by zero when $h<g(C)$.

\begin{theorem}\label{p9-thm:global}
Let $X$ be a closed symplectic six-manifold with $c_1(TX)=0$ and torsion-free $H_2(X;\ZZ)$. For every primitive class $A$ of
positive area and every $h\ge0$,
\begin{equation}\label{p9-eq:global}
 e^{\rm IP}_{2A,h}(X)-\operatorname{Gr}^{\rm BS}_{2A,h}(X)
       =\sum_{g\ge0} e^{\rm IP}_{A,g}(X)\,a_{g,h}.
\end{equation}
The coefficient $a_{g,h}$ vanishes for $g=0,1$. For $c=g-1\ge1$ and
$j=h-(2g-1)$ it is
\begin{equation}\label{p9-eq:coeff}
 a_{g,h}=
 \begin{cases}
  2\,4^{c-j}\left(\binom{c}{2j}-(-1)^j\binom{c}{j}\right),&1\le j\le c,\\[2pt]
  0,&\text{otherwise}.
 \end{cases}
\end{equation}
A binomial coefficient outside its usual range is zero.
\end{theorem}

The first two nonzero rows of the correction, written in the variable
$x$ with genus $h$ contributing $x^{h-1}$, are $2x^3$ and
$24x^5-2x^6$. On the compact threefold of Section~\ref{p9-sec:example},
a primitive class $A$ has $e^{\rm IP}_{A,3}=-14$ and no higher-genus
primitive coefficient. Only the term $(-14)(-2)$ contributes in
genus seven. Corollary~\ref{p9-cor:example} gives a family on the same
threefold for which the absolute differences grow quadratically.

\subsection{The fixed-target comparison}
Let $C$ be a smooth closed connected Riemann surface of genus $g$ and
let $N\to C$ be a complex bundle of rank two and degree $2g-2$. Let $D$
be a real Cauchy--Riemann operator on $N$. Its index on $C$ is zero.
When $D$ is invertible, $\epsilon(D)$ is its sign for the determinant
orientation obtained by transport from complex-linear operators.
Suppose also that $f^*D$ is injective for every smooth connected
double cover $f$ of $C$. We call such an operator rigid through degree
two. On a nodal stable map of degree two one can normalize the positive
components, each of degree one or two. Kernel sections vanish there,
and the constants on each connected contracted subcurve vanish by
matching. Hence all stable pullbacks are injective.

Write $M_k=\overline{\mathcal M}_k(C,2)$ for the stack of connected
degree-two stable maps with unmarked source of arithmetic genus $k$.
The oriented real cokernels of $f^*D$ give a bundle on $M_k$ of real rank
$2r$, where $r=2k-2-2(2g-2)$ is the complex virtual dimension of $M_k$. Define
\[
 L_{2,k}(D)=\int_{[M_k]^{\rm vir}} e(\operatorname{coker} f^*D),
 \qquad B_{2,g}^{\rm bal}(v)=[q^2]G_g(q,v).
\]
The superscript $\rm bal$ denotes the balanced antidiagonal series.
The integral is zero for empty moduli. We use the orientation of the
full pullback determinant for $L_{2,k}$, including its sign in rank
zero. In the orientation relative to the base operator, this number
is $\epsilon(D)L_{2,k}(D)$ \cite[Remark~6.4]{bai}.

\begin{theorem}\label{p9-thm:local}
For every operator rigid through degree two one has
\begin{equation}\label{p9-eq:local}
 \epsilon(D)\sum_{k\ge0}L_{2,k}(D)v^{2k-2}
    = B_{2,g}^{\rm bal}(v)+
       \sum_{h\ge0}\bigl(w_{2,h}(D)+a_{g,h}\bigr)x^{h-1}.
\end{equation}
The equality is coefficientwise in $v$. In genus zero the series may
have a pole of order two.
\end{theorem}

\subsection{Comparison with the chamber weight}\label{p9-sec:comparison}

\subsubsection{Transport from a complex-linear operator}
Fix $C$ and $N$ as in Theorem~\ref{p9-thm:local}. Use an invertible
complex-linear operator $\beta$ rigid through degree two as a reference.
Such operators exist by the complex-operator transversality in
\cite[Lemmas~5.3, 5.6 and~5.15(ii), Definition~5.17]{bai}. Base acyclicity is
open and is achieved by taking a direct sum of degree-$g-1$ line
bundles off the theta divisor, or $\mathcal O(-1)^{\oplus2}$ when
$g=0$. The reference need not split holomorphically.

Choose a generic path from $\beta$ to $D$ in the affine space of real
Cauchy--Riemann operators with fixed symbol. At an ordinary crossing
of the base wall the kernel is a real line of nowhere-zero sections.
At an ordinary crossing of a double-cover wall its anti-invariant
kernel line gives an embedded graph and its intrinsic derivative
\eqref{p9-eq:wall-differential} is invertible. Here one uses the
one-parameter transversality and embeddedness of the kernel graphs
\cite[Lemmas~5.6, 5.11, 5.15--5.16]{bai}. The variations in these tests
are variations of the zeroth-order real tensor on an unramified disk.
The fixed-symbol space also has the proper degree-at-most-two wall of
\cite[Definition~2.6, Proposition~2.7 and Theorem~3.40]{dw}. In
particular, up to any prescribed genus there are only finitely many
crossings.

We give the transversality detail that separates a base crossing from
an odd crossing. Mark real kernel lines on the base and in the
anti-invariant part of a smooth cover with $b$ branch points. For a
fixed cover the Fredholm indices of the two operators are $0$ and
$-2b$. Projectivizing both kernels and allowing the branch values to
vary gives index $-2$. For surjectivity of the universal section,
an adjoint pair annihilating all variations of the operators would,
on an unramified disk, annihilate all pairs of values
$(A\xi_0,A\xi_1)$, where $A$ is a localized real tensor and the two
generators are evaluated on one sheet. The odd kernel and its
adjoint both change sign on the second sheet, so the contributions
of the sheets add. The generators are real independent on a disk.
Otherwise on a disk of nonvanishing a real multiple relating two
solutions has zero $\bar\partial$, hence is constant. Unique
continuation on the connected cover then contradicts their opposite
deck characters. Arbitrary variations of $A$ on the disk force both
adjoints to vanish there and hence globally. Thus a generic path
misses this incidence. Taking the countable intersection over source
genera preserves genericity.

Put $E_{2,k}(D)=\epsilon(D)L_{2,k}(D)$ on regular slices. At a base
crossing with no odd kernel this reduced Euler number extends across
the crossing. To see this, for a smooth positive component of degree
two the invariant sections in the pullback kernel descend away from the
branch values. They are bounded there in a target trivialization.
Removable singularities for the real Cauchy--Riemann equation extend
them across the branch values. Since there is no odd kernel, all kernel
sections are therefore pulled back from the base. If there are two
degree-one components, their base sections agree at a point, either at
a common node or along the chain of contracted components joining
them. Their difference vanishes everywhere because a nonzero base
kernel section is nowhere zero. Constants on each contracted subcurve
are determined by their nodal values. The same description on a single
positive component is compatible with any further node. Consequently
the kernel on every stable map is precisely the base pullback.

Represent the base and pullback families by compatible stabilized
two-term complexes of finite-dimensional bundles as in
\cite[Appendix~B, (B.7)--(B.10)]{bai}. The reduced complex is the
quotient in this finite-dimensional presentation. To see that its
degree-zero cohomology vanishes, it remains, by the cohomology exact
sequence and the kernel computation above, to check that pullback
is injective on cokernels. Represent a base cokernel class by a smooth
form $\alpha$. If $f^*D u=f^*\alpha$ on a single effective double
normalization, averaging $u$ over the two sheets gives a base solution
of $D\bar u=\alpha$ off the branch values. It extends by bounded
removability for the inhomogeneous equation. On two degree-one components
either restriction supplies such a solution. Thus the base class
vanishes. The finite-dimensional reduced differential is injective
over the whole compact stable-map space at the crossing.
Its cokernels form an oriented bundle across the base crossing. The
determinant identity is
\[
 \det(f^*D)\simeq\det D\otimes\det((f^*D)_{\rm red}).
\]
In particular the reduced integral on either side is $E_{2,k}$.

At a double-cover crossing the base sign is constant. In the convention
of \cite[(5.59)--(5.60)]{bai}, the forward change of $w_{2,H}$ is
$\epsilon(D_0)\eta_f$ for a cover of genus $H$. Applying
Theorem~\ref{p9-thm:wall} and summing the jumps within each source genus
gives
\begin{equation}\label{p9-eq:transport}
 \sum_{k\ge0}\bigl(E_{2,k}(D)-E_{2,k}(\beta)\bigr)v^{2k-2}
       =\sum_{h\ge0} w_{2,h}(D)x^{h-1}.
\end{equation}
At the initial endpoint $w_{2,h}(\beta)=0$ by the constant path in its
definition. The equality follows for every formal coefficient using
compactness only through that coefficient's genus.

\subsubsection{The Euler integral at the reference}
Let $\pi$ and $f$ now denote the universal curve and evaluation map
over $M_k$. At $\beta$ the class $-R\pi_* f^*N$ is represented by its
actual complex cokernel. Its rank is $r=\dim_{\CC}^{\rm vir}M_k$ and
the base determinant has positive sign. Thus its ordinary top Chern
number is $E_{2,k}(\beta)$. To compute it give $N$ scalar torus weight
$t$. If $V$ is the virtual index class of rank $r$, its equivariant Euler
class is $t^r c_{1/t}(V)$, with $c_z$ the total Chern polynomial.
Integration on the virtual cycle selects $c_r(V)$. This description
also permits deformation of the holomorphic structure to a sum of
two lines of degree $g-1$, since the family index class extends along
the deformation, including at split operators with kernel.

Equal weights compute the reference number, whereas opposite weights
compute the elementary contribution already assigned to the base curve.
The scalar and antidiagonal specializations have the same degree-one
term $x^{g-1}$. In degree two taking the connected logarithm subtracts
the same half-square from both. The scalar formula of
\cite[\S8]{bp08} and the two antidiagonal hook products in
\cite[Corollary~7.2]{bp08} therefore give
\begin{align*}
 \sum_{k\ge0} E_{2,k}(\beta)v^{2k-2}-B_{2,g}^{\rm bal}(v)
  =x^{2g-2}\bigl((4-2s)^{g-1}+(4+2s)^{g-1}
                           -2(4-x)^{g-1}\bigr).
\end{align*}
The bracket vanishes for $g=1$. For $g=0$ the sum of the first
two reciprocals is $2/(4-x)$. For $c=g-1>0$ the odd powers of $s$
cancel. The coefficient of $x^{2c+j}$ is \eqref{p9-eq:coeff} for
$h=2g-1+j$, and the constant term in the bracket vanishes. Together
with \eqref{p9-eq:transport} this proves Theorem~\ref{p9-thm:local}.

\subsubsection{Passage to the compact invariants}
Choose a generic compatible almost complex structure for which all
simple curves of index zero are embedded, regular, super-rigid and
pairwise disjoint \cite[\S1]{bai}. The image of a connected stable
map in class $2A$ lies on a single such support. Indeed, two distinct
effective images cannot meet, and contracted components can only join
them at a common image point. If the map has degree $m$ on its support
of class $B$, then $mB=2A$. Primitivity and the torsion-free lattice
give $m=1$, $B=2A$, or $m=2$, $B=A$. The same argument in class $A$
gives only degree one.

Through any fixed source genus there are finitely many contributing
simple supports, by compactness at fixed genus and super-rigidity
\cite[\S\S3,5]{bai}. The degree-one series on a support of genus $g$
is $\epsilon_C x^{g-1}$, by the degree-one calculation in
Theorem~\ref{p9-thm:wall}. Hence $e^{\rm IP}_{A,g}$ is the sum of signs
of genus-$g$ supports of class $A$. In class $2A$ first subtract their
degree-two elementary terms $\epsilon_C B_{2,g}^{\rm bal}$. By
\eqref{p9-eq:local}, the remaining coefficient of $x^{h-1}$ on each such
support is $\epsilon_C(w_{2,h}(D_C)+a_{g,h})$. A support in $2A$
contributes its degree-one determinant sign in its own genus. The
latter signs and the terms $\epsilon_C w_{2,h}$ give exactly
\cite[(1.5)]{bai}. The remaining terms sum to \eqref{p9-eq:global},
proving Theorem~\ref{p9-thm:global}.
\subsection{A compact threefold with unbounded discrepancy}\label{p9-sec:example}

Let $X_9$ be a general nondegenerate anticanonical hypersurface in the smooth
projective toric fourfold with rays and maximal cones
\cite[(23)--(24)]{WueCounter}. We equip it with a K\"ahler form.
It has trivial canonical bundle and torsion-free $H_2(X_9;\ZZ)$
\cite[\S4.1]{WueCounter}. The divisor corresponding to ray $8$ cuts out
a surface $T\simeq\mathbb F_1$. Write $s$ for its negative section, $f$ for a
ruling fiber and $i:T\hookrightarrow X_9$ for the inclusion. For $d\ge3$ let
$D_d=ds+(d+1)f$ and $A_d=i_*D_d$. Set
\begin{equation}\label{p9-eq:compact-classes}
 G=\frac{d(d-1)}2,\qquad N=\frac{(d+1)(d+4)}2 .
\end{equation}

\begin{corollary}\label{p9-cor:example}
The class $A_d$ is primitive, and
\[
 e^{\rm IP}_{2A_d,3G-2}(X_9)-\operatorname{Gr}^{\rm BS}_{2A_d,3G-2}(X_9)
     =(-1)^{d+1}(d+1)(d+4).
\]
In particular, the two counts differ by $28$ in class $2i_*(3s+4f)$
and genus seven. Their differences are unbounded on $X_9$.
\end{corollary}

\subsubsection{The surface classes}
The semiample divisor $H_4$ of \cite[\S4.2]{WueCounter} contracts curves
only in $T$, in a second surface $E\simeq\operatorname{Bl}_{p_1,p_2}\PP^2$,
or in two rational curves $C_1,C_2$. On $E$ the effective cone is generated
by its exceptional curves $e_1,e_2$ and the strict transform of the line
through the two points, which equals $s=E\cap T$. The six integral nef
divisors of \cite[Lemma~4.1 and (30)]{WueCounter} have degrees
\[
\begin{array}{c|rrrrrr}
 & H_1&H_2&H_3&H_4&H_5&H_6\\ \hline
 s   &0&1&0&0&0&0\\
 f   &0&0&1&0&0&0\\
 e_1 &1&0&0&0&1&1\\
 e_2 &0&0&0&0&0&1\\
 C_1 &0&0&0&0&1&0\\
 C_2 &1&0&0&0&0&0
\end{array}
 .
\]
In particular $H_3-H_2$ pairs to one with $A_d$, and $H_2,H_3$
distinguish all classes on $T$. Each effective summand of $A_d$ has zero
$H_1,H_4,H_5,H_6$ degrees. These conditions exclude $C_1,C_2$ and force
any component on $E$ to have class proportional to its negative curve $s$.
Such a component equals $s$. Thus every summand is represented entirely
on $T$. This is the argument of \cite[Proposition~4.2]{WueCounter}
for arbitrary $d$.

\subsubsection{The primitive coefficient}
Adjunction gives $N_{T/X_9}=K_T$. To compute the connected invariants on
these classes, we deform $T\subset X_9$ to $Z=\operatorname{Tot}(K_T)$. Blowing up $T\times 0$ in $X_9\times\mathbb A^1$ and removing
the strict transform of the central fiber gives a smooth family $V$ with
fibers $X_9$ away from zero and $Z$ over zero. It contains the Cartier
divisor $\mathcal T=T\times\mathbb A^1$ and has a natural morphism to $X_9$.
Use this morphism to label curve classes. In $Z$ every complete irreducible
curve lies in $T$. Its projection must be nonconstant since the fibers are
affine lines, and the tautological fiber coordinate on its normalization
is a section of a line bundle of negative degree.

For every positive class on this face, stable maps from connected domains to
the fibers of $V$ factor through $\mathcal T$ in families. Indeed on a
geometric fiber the image lies in $T$ by the preceding support calculation.
The pullback of $\mathcal O_V(\mathcal T)$ then has negative degree on
each nonconstant component and is trivial on each contracted component.
Matching at the nodes makes its space of sections zero. For a family of
nodal domains the derived direct image of this line bundle is locally a
two-term complex of vector bundles in degrees zero and one. Fiberwise
injectivity of its differential implies that its degree-zero cohomology
sheaf vanishes after every base change. The defining section of $\mathcal T$
therefore pulls back to zero, giving the asserted factorization also over
nonreduced bases. The relative stable-map stack in the specified class
and genus is consequently the product of a proper stable-map stack of $T$
with $\mathbb A^1$. Deformation invariance of its relative virtual class
identifies the connected invariants of $X_9$ and $Z$ in every genus.

Write $D=D_d$. On $T=\operatorname{Bl}_p\PP^2$ it is the class of a plane
curve of degree $d+1$ through $p$ with one imposed simple base point.
Thus $p_a(D)=G$ and $h^0(T,\mathcal O_T(D))=N$. We claim that the least
power of $z$ in the stable-pair partition function of $Z$ in class $D$ is
$z^{1-G}$, with coefficient $(-1)^{N-1}N$, and that this term survives
in the logarithm. The case $d=3$ of the following moduli calculation appears
in \cite[Propositions~5.2--5.3]{WueCounter}.

For any splitting $D=\sum_j D_j$ into at least two nonzero effective classes
on $T$, write $D_j=a_js+b_jf$ with $a_j,b_j\ge0$. From
$s^2=-1$, $sf=1$ and $f^2=0$ we get
\[
 \sum_j\chi(\mathcal O_{D_j})-\chi(\mathcal O_D)
  =\sum_{j<k}D_jD_k
  =\sum_j b_j(d-a_j)-\sum_{j<k}a_ja_k>0.
\]
Indeed, let $m=\max a_j$ and $r=d-m$. If $r>0$, the positive sum is
at least $(d+1)r$, and the sum being subtracted is at most
$mr+r(r-1)/2$. The difference is at least $r(r+3)/2$.
If $r=0$, the remaining terms are positive multiples of $f$, giving a
difference of at least $d$.

The image of a stable pair is the structure sheaf of a pure curve $B$,
and its cokernel has finite length. Filter $\mathcal O_B$ by powers of
the ideal of $T$. Its $j$th quotient is
$K_T^{-j}\otimes(\mathcal O_T/J_j)$ for some ideal $J_j$, since the
conormal bundle is $K_T^{-1}$. Write
$J_j^{\vee\vee}=\mathcal O_T(-D_j)$ and
$\ell_j=\operatorname{length}(J_j^{\vee\vee}/J_j)$. Then
\[
 \chi(\mathcal O_B)=
 \sum_j\bigl(\chi(\mathcal O_{D_j})-jK_TD_j+\ell_j\bigr),
 \qquad \sum_j[D_j]=[B].
\]
The correction terms are nonnegative because $-K_T$ is ample. If there
are several divisorial layers, the strict inequality above gives
$\chi(\mathcal O_B)>1-G$ in class $D$. If there is only one, it is the
zeroth layer by Nakayama's lemma at the generic points. The ideal of
$T$ then generates a subsheaf of $\mathcal O_B$ of finite support, which
vanishes by purity. Hence equality occurs precisely for quotient pairs
$\mathcal O_Z\to\mathcal O_B$ with $B\in|D|$. For a product of pairs
in two or more nonzero classes summing to $D$, their combined filtrations
have at least two nonzero divisorial layers. The same strict inequality
shows that all product terms in the logarithm occur above $z^{1-G}$.

The quotient pairs give a closed $\PP^{N-1}=|D|$ in the minimum
pair space. They exhaust its closed points, and surjectivity of their
sections persists in every infinitesimal deformation by Nakayama's lemma.
For every $B\in|D|$, including nonreduced $B$, the normal sequence is
\[
 0\longrightarrow\mathcal O_B(D)\longrightarrow N_{B/Z}
   \longrightarrow K_T|_B\longrightarrow0.
\]
Now $H^0(K_T)=0$ and $H^1(K_T-D)=H^1(\mathcal O_T(D))^*=0$ by Kodaira
vanishing, since $D-K_T=(d+2)s+(d+4)f$ is ample.
Thus $H^0(K_T|_B)=0$, while $h^0(\mathcal O_B(D))=N-1$ by rationality
of $T$. The Hilbert tangent dimension is $N-1$ at every closed point.
It equals the local dimension, so the pair space is regular and the closed
immersion is an isomorphism. Its symmetric obstruction theory has virtual
degree $(-1)^{N-1}\chi_{\rm top}(\PP^{N-1})=(-1)^{N-1}N$
\cite[\S1]{pt}.

Finally, apply the toric Gromov--Witten/stable-pairs correspondence
\cite[Theorem~1]{pp} to $Z$, with empty insertions and $-z=e^{iu}$.
The Calabi--Yau condition gives prefactor one. The stable-map and
stable-pair spaces in these classes are proper. For pairs, compactify $Z$
to the projective line bundle over $T$. In a flat family over a valuation
ring the support is the closure of the generic support, so a pair
supported set-theoretically on the zero section specializes there as well.
Thus the open space of pairs supported in $Z$ is proper, and the
equivariant integrals give ordinary virtual degrees. The disconnected
Gromov--Witten series in this correspondence omits degree-zero connected
components, so taking its logarithm gives the connected series.
The equality applies to every effective summand on the face.

On the projective fiber, the algebraic and symplectic virtual degrees
agree \cite{siebertgw}. By primitivity, the coefficient of class $A_d$ in
the connected elementary expansion is the finite sum of degree-one terms
\cite[Proposition~3.3 and Theorem~8.1]{ip}, \cite{diw}. It follows that
\[
 [Q^D]\log Z_{\rm PT}^{Z}(z,Q)
     =\sum_{g\ge0} e^{\rm IP}_{A_d,g}(X_9)
           \left(\frac{(1+z)^2}{z}\right)^{g-1}.
\]
Both sides are rational in $z$. The correspondence identifies their expansions
at $-z=e^{iu}$, since $(2\sin(u/2))^2=(1+z)^2/z$, and therefore also
their Laurent expansions at zero. Each summand of genus $g$ has lowest
power $z^{1-g}$ with leading factor one. The minimum computed above gives
\begin{equation}\label{p9-eq:compact-primitive}
 e^{\rm IP}_{A_d,G}(X_9)=(-1)^{N-1}N,\qquad
 e^{\rm IP}_{A_d,g}(X_9)=0\quad(g>G).
\end{equation}

\begin{proof}[Proof of Corollary~\ref{p9-cor:example}]
In \eqref{p9-eq:global} a row with base genus $g<G$ vanishes in
genus $3G-2$ by \eqref{p9-eq:coeff}. All rows with $g>G$ vanish by
\eqref{p9-eq:compact-primitive}. The remaining coefficient in
\eqref{p9-eq:coeff} is $a_{G,3G-2}=-2(-1)^{G-1}$. Multiplication by
$(-1)^{N-1}N$ gives the formula, since $G+N-2=d(d+2)$.
For $d=3$ the minimum pair space is $\PP^{13}$ and its virtual degree
is $-14$, giving $(-2)(-14)=28$. For $d=15$ the difference is $304$
in genus $313$. As $d$ ranges through $12\ell+3$ these positive differences
tend to infinity.
\end{proof}

\section{The first genera above a ramified double wall}\label{sec:quartic}
Let \(f:Y\to C\) be an ordinary double crossing with a single odd kernel line of embedded graph, \(b=2\ell>0\) simple branch values, and sign \(\eta_f\). Put \(h=H_2+\ell\). Use the generic nonvanishing of the projected antilinear jet at each branch, the classification of smooth central kernels as pullbacks of the fixed graph line, and invertibility of the unramified upper normal twists used below. For clarity about genericity, on a regular disk one varies the real order-zero tensor while annihilating the graph vector. Perturbing a good Petri jet, locally solving for a nearby section, and using a complex gauge sending it to the prescribed section preserves that vector. The adjoint test on irreducible monodromy systems then leaves only its sign-character summand; the moving branch positions are included only at genuinely ramified punctures of each system. The fixed graph, polar and quartic transversality proofs for the complete residual intersection are given in Calculations G1--G3 of \cite{WueSupplement}.
For any regular point of \(C\) put
\[
 B_m=\sum_{[u:T\to C]}\frac{m}{|\operatorname{Aut}_C(u)|}\epsilon(u^*D_0)
\]
over connected unramified covers of degree \(m\). Denote by \(A\) the normal operator of the graph in its homogeneous target, \(T=\tfrac12\sum_{\lambda\ne0}\epsilon(A^\lambda)\) over the sign twists of \(Y\), and \(N_-=\#\{\lambda\ne0:\epsilon(A^\lambda)=-1\}\).

\begin{proposition}
For \(d\ge3\) there is no change below \(k_d=H_d+\ell+1\). At this genus the connected and elementary jumps, respectively, are
\[
 -b\eta_f B_{d-2}+\mathbf1_{(d,\ell)=(4,1)}\eta_f T,\qquad
 -b\eta_f B_{d-2}-\mathbf1_{(d,\ell)=(4,1)}\eta_f N_- .
\]
The comparison is coefficientwise. For elementary exponents the endpoints admit the balanced lift.
\end{proposition}
\begin{proof}
On the positive normalizations denote the total degree of the nonzero ones over \(Y\) by \(a\), their number together with the zero-height positive components by \(n\), and their total ramification over their respective targets \(Y,C\) by \(R\). If \(G\) sums the contracted genera and \(\beta\) is the graph Betti number, the genus identity is \(k-H_d=n-1+a\ell+R/2+G+\beta\). When \(n\ge2\), equality at the minimum forces \(n=2,a=1,R=G=\beta=0\). Thus a single branch of \(Y\) joins a pointed étale cover at zero height. The smoothing degree is \(-\eta_f\epsilon(u^*D_0)\) per marked orbit, by the one-node logarithmic calculation of \S\ref{p10-sec:nodes} on that sheet. The deck action is free on the marked fiber. Summing yields the first term. When \(n=1\) equality forces \((d,\ell)=(4,1)\) and an unramified upper double. Its odd determinant is that of the graph normal twist and its covering involution divides the count by two. This adds \(\eta_f T\). Proper-divisor contributions from a genus excess \(e\) are supported above excess \(je\) on a \(j\)-fold multiple. Only the exceptional pair requires subtraction at equality, namely \(\eta_f(2^{2h}-1)/2\), the elementary double of the lower curve. Since \(T-(2^{2h}-1)/2=-N_-\), the claim follows.
\end{proof}

Modulo two the pointed sum counts index-\(m\) subgroups of \(\pi_1(C)\). Its parity is odd precisely for \(m\) a square or twice a square, by the surface-group parity theorem of M\"uller and Puchta \cite[Theorem~1]{mp}. Outside the exception the jump in these degrees is an odd multiple of \(b\). For a cubic crossing the central classification of \S\ref{sec:heptic} gives no changes in degrees one, two, four or five. At \(d=6+m,\ 3\nmid d\), the first possible genus is \(H_d+3\ell+1\), with a sextic closure joined to a pointed unramified \(m\)-cover. A normalized section that reaches the zero section is a negative eigenvector of a peripheral transposition; its stabilizer as a vector in the permutation plane is trivial. Its component must therefore cover the sextic closure. Applying the genus identity with this component gives the stated bound and forces just one join at equality. The adjoint and polar equations yield odd local degree for each pointed cover, with the upper determinant included. More precisely, writing \(a_i\) for the signed degree with one identity sheet at the \(i\)-th branch, each pointed upper cover \(u\) contributes \(a_i\epsilon(u^*D_0)\epsilon(D_0)\). Hence the complete first-genus jump is \(\epsilon(D_0)B_m\sum_i a_i\) for both connected and elementary counts. Calculations H2--H3 of \cite{WueSupplement} derive the discriminant adjoint line and its sign. The complete series at \(d=7\) is Theorem~\ref{p14-thm:wall}.

\subsection{The next quartic genus}
Set \(k_0=h+2g-1,\ k_1=k_0+1\), and put \(S_C=B_2\). At each lower ramification point \(r\), let \(E_r\) be the integral of the Euler class of the pole-augmented upper obstruction line along the ramification divisor of the marked double-cover surface. The upper section is allowed the simple pole fixed by the projected derivative at \(r\). Write \(F_2(D)\) for the corresponding ordinary upper finite-map Euler integral over the two-branch surface. For \(\ell\ge2\) the complete elementary coefficient at \(k_1\) is
\[
 \Delta n_{4,k_1}=\eta_f\left[(2S_C+1)\binom b2-2bF_2(D_0)-\sum_r E_r
                   -\mathbf1_{\ell=2}N_-\right].
\]
Two distinct flags give the binomial term. Each marked surface contributes \(-\eta_f(2F_2+E_r)\). A tower reaches this genus only for four lower branch values. For two lower branch values the two copies of the lower graph vary in their relative real height. Displace on one common path. Denote by \(D_\theta\) the base operator displaced in this direction and by \(A_\delta\) the displaced upper graph operator. Define
\[
 R=\eta_f\left[\epsilon(A_\delta)F_2(A_\delta)-4F_2(D_\theta)
                   -\sum_r E(D_\theta;r)\right].
\]
Then \(\Delta n_{4,k_1}=R+\eta_f[2S_C+1+(h-1)2^{2h-3}]\).
In both cases the full connected coefficient is recovered by adding the lower degree-two factor and the next Taylor coefficient of the preceding degree-four factor. The detailed two-height calculation in the companion proves independence of the sign of displacement. Its interior chart has transverse degree \(-3\). The common-parameter perturbation excludes smoothings at unequal and zero heights; the equal-height endpoint uses the graph operator and includes all contracted expansions of its regular towers. The same comparison gives an upper ordinary odd Euler jump \(-3b\eta_A\) with at least three smooth crossings near each lower branch after a generic perturbation. For \(2k-1\) prescribed plus-axis poles the signed degree in genus \(2g-1+k\) is \((-1)^{k+1}(2k-1)2^{2g}/4\). The pole residue and its stack normalization are worked out in the accompanying calculations on marked covers and polar crossings.

\begin{thebibliography}{Wue26f}
\bibitem[BS24]{bai} S.~Bai and M.~Swaminathan,
\emph{Bifurcations of embedded curves and toward an extension of Taubes's
Gromov invariant to Calabi--Yau $3$-folds}, Duke Math. J. \textbf{173}
(2024), 2947--3057,
\href{https://doi.org/10.1215/00127094-2023-0070}{doi:10.1215/00127094-2023-0070}.
Numbering as in \href{https://arxiv.org/abs/2106.01206v4}{arXiv:2106.01206v4}.

\bibitem[BP08]{bp08} J.~Bryan and R.~Pandharipande,
\emph{The local Gromov--Witten theory of curves}, J. Amer. Math. Soc.
\textbf{21} (2008), 101--136,
\href{https://doi.org/10.1090/S0894-0347-06-00545-5}{doi:10.1090/S0894-0347-06-00545-5}.

\bibitem[DIW26]{diw} A.~Doan, E.-N.~Ionel and T.~Walpuski,
\emph{The Gopakumar--Vafa finiteness conjecture},
Ann. of Math. (2) \textbf{204} (2026), 819--871,
\href{https://doi.org/10.4007/annals.2026.204.3.1}{doi:10.4007/annals.2026.204.3.1}.
Numbering as in \href{https://arxiv.org/abs/2103.08221v2}{arXiv:2103.08221v2}.

\bibitem[DW23]{dwbn} A.~Doan and T.~Walpuski,
\emph{Equivariant Brill--Noether theory for elliptic operators and superrigidity of
$J$-holomorphic maps}, Forum Math. Sigma \textbf{11} (2023), e3,
\href{https://doi.org/10.1017/fms.2022.104}{doi:10.1017/fms.2022.104}.

\bibitem[DW25]{dw} A.~Doan and T.~Walpuski,
\emph{Chambered invariants of real Cauchy--Riemann operators},
\href{https://arxiv.org/abs/2410.21057v3}{arXiv:2410.21057v3} (2025).

\bibitem[ES25]{ekholm} T.~Ekholm and V.~Shende,
\emph{Ghost bubble censorship}, Comm. Anal. Geom. \textbf{33}
(2025), 1815--1863,
\href{https://doi.org/10.4310/CAG.260215235638}{doi:10.4310/CAG.260215235638}.

\bibitem[IP18]{ip} E.-N.~Ionel and T.~H.~Parker,
\emph{The Gopakumar--Vafa formula for symplectic manifolds},
Ann. of Math. (2) \textbf{187} (2018), 1--64,
\href{https://arxiv.org/abs/1306.1516}{arXiv:1306.1516}.

\bibitem[Küh23]{kuehne} L.~K\"uhne,
\emph{The universality of the resonance arrangement and its Betti numbers},
Combinatorica \textbf{43} (2023), 277--298,
\href{https://doi.org/10.1007/s00493-023-00006-x}{doi:10.1007/s00493-023-00006-x}.

\bibitem[MP03]{mp} T.~W.~M\"uller and J.-C.~Puchta, \emph{Parity patterns in one-relator groups}, J. Group Theory \textbf{6} (2003), no.~2, \href{https://doi.org/10.1515/JGTH.2003.018}{doi:10.1515/JGTH.2003.018}.

\bibitem[Pan99]{pandhodge} R.~Pandharipande,
\emph{Hodge integrals and degenerate contributions},
Comm. Math. Phys. \textbf{208} (1999), 489--506,
\href{https://doi.org/10.1007/s002200050766}{doi:10.1007/s002200050766}.

\bibitem[PP14]{pp} R.~Pandharipande and A.~Pixton,
\emph{Gromov--Witten/pairs descendent correspondence for toric 3-folds},
Geom. Topol. \textbf{18} (2014), 2747--2821,
\href{https://doi.org/10.2140/gt.2014.18.2747}{doi:10.2140/gt.2014.18.2747}.

\bibitem[PT10]{pt} R.~Pandharipande and R.~P.~Thomas,
\emph{Stable pairs and BPS invariants}, J. Amer. Math. Soc.
\textbf{23} (2010), 267--297,
\href{https://arxiv.org/abs/0711.3899v4}{arXiv:0711.3899v4}.

\bibitem[Par16]{pardonatlas} J.~Pardon,
\emph{An algebraic approach to virtual fundamental cycles on moduli spaces
of pseudo-holomorphic curves}, Geom. Topol. \textbf{20} (2016), 779--1034,
\href{https://doi.org/10.2140/gt.2016.20.779}{doi:10.2140/gt.2016.20.779}.

\bibitem[Par25]{pardon} J.~Pardon,
\emph{Universally counting curves in Calabi--Yau threefolds},
\href{https://arxiv.org/abs/2308.02948v3}{arXiv:2308.02948v3} (2025).

\bibitem[Sie98]{siebert} B.~Siebert,
\emph{Gromov--Witten invariants of general symplectic manifolds},
\href{https://arxiv.org/abs/dg-ga/9608005v2}{arXiv:dg-ga/9608005v2} (1998).

\bibitem[Sie99]{siebertgw} B.~Siebert,
\emph{Algebraic and symplectic Gromov--Witten invariants coincide},
Ann. Inst. Fourier (Grenoble) \textbf{49} (1999), 1743--1795,
\href{https://arxiv.org/abs/math/9804108v1}{arXiv:math/9804108v1}.

\bibitem[Stacks]{stacks} The Stacks Project Authors,
\emph{The Stacks Project}, The Cotangent Complex,
\href{https://stacks.math.columbia.edu/tag/08SL}{Tag~08SL}.

\bibitem[Swa21]{swam} M.~Swaminathan,
\emph{Rel--$C^\infty$ structures on Gromov--Witten moduli spaces},
J. Symplectic Geom. \textbf{19} (2021), 413--473,
\href{https://doi.org/10.4310/JSG.2021.v19.n2.a4}{doi:10.4310/JSG.2021.v19.n2.a4}.

\bibitem[Wen23]{wendl} C.~Wendl,
\emph{Transversality and super-rigidity for multiply covered holomorphic curves},
Ann. of Math. (2) \textbf{198} (2023), 93--230,
\href{https://doi.org/10.4007/annals.2023.198.1.2}{doi:10.4007/annals.2023.198.1.2}.

\bibitem[Wue26a]{WueCounter} B.~J.~Wuebben,
\emph{A counterexample to integral product exponents in universal curve counting}, 28 October 2026, revised 21 November 2026,
\href{https://doi.org/10.5281/zenodo.23750048}{doi:10.5281/zenodo.23750048}.

\bibitem[Wue26b]{WueProduct} B.~J.~Wuebben,
\emph{A geometric product formula for universal curve counting on Calabi--Yau threefolds}, 23 November 2026,
\href{https://doi.org/10.5281/zenodo.24455518}{doi:10.5281/zenodo.24455518}.

\bibitem[Wue26e]{WueGenerators} B.~J.~Wuebben,
\emph{Integral generators for curve counting on threefolds}, 24 November 2026,
\href{https://doi.org/10.5281/zenodo.24510962}{doi:10.5281/zenodo.24510962}.

\bibitem[Wue26f]{WueSpin} B.~J.~Wuebben,
\emph{Spin structures and local Gromov--Witten invariants of curves}, 24 November 2026,
\href{https://doi.org/10.5281/zenodo.24510976}{doi:10.5281/zenodo.24510976}.

\bibitem[Wue26g]{WueSupplement} B.~J.~Wuebben,
\emph{Supplementary calculations for branched covers and spin structures in local curve counting}, Calculations A1--H3, 24 November 2026, archived with \emph{Branched triple covers and wall crossing for local curve counts},
\href{https://doi.org/10.5281/zenodo.24510974}{doi:10.5281/zenodo.24510974}.

\bibitem[Wue26j]{WueUnramified} B.~J.~Wuebben,
\emph{Unramified covers and wall crossing for local Gromov--Witten invariants}, 24 November 2026,
\href{https://doi.org/10.5281/zenodo.24510972}{doi:10.5281/zenodo.24510972}.

\bibitem[Zin11]{zinger} A.~Zinger,
\emph{A comparison theorem for Gromov--Witten invariants in the
symplectic category}, Adv. Math. \textbf{228} (2011), 535--574,
\href{https://doi.org/10.1016/j.aim.2011.05.021}{doi:10.1016/j.aim.2011.05.021}.
\end{thebibliography}

\end{document}
